% Options for packages loaded elsewhere
% Options for packages loaded elsewhere
\PassOptionsToPackage{unicode}{hyperref}
\PassOptionsToPackage{hyphens}{url}
\PassOptionsToPackage{dvipsnames,svgnames,x11names}{xcolor}
%
\documentclass[
  letterpaper,
  DIV=11,
  numbers=noendperiod]{scrreprt}
\usepackage{xcolor}
\usepackage{amsmath,amssymb}
\setcounter{secnumdepth}{2}
\usepackage{iftex}
\ifPDFTeX
  \usepackage[T1]{fontenc}
  \usepackage[utf8]{inputenc}
  \usepackage{textcomp} % provide euro and other symbols
\else % if luatex or xetex
  \usepackage{unicode-math} % this also loads fontspec
  \defaultfontfeatures{Scale=MatchLowercase}
  \defaultfontfeatures[\rmfamily]{Ligatures=TeX,Scale=1}
\fi
\usepackage{lmodern}
\ifPDFTeX\else
  % xetex/luatex font selection
\fi
% Use upquote if available, for straight quotes in verbatim environments
\IfFileExists{upquote.sty}{\usepackage{upquote}}{}
\IfFileExists{microtype.sty}{% use microtype if available
  \usepackage[]{microtype}
  \UseMicrotypeSet[protrusion]{basicmath} % disable protrusion for tt fonts
}{}
\makeatletter
\@ifundefined{KOMAClassName}{% if non-KOMA class
  \IfFileExists{parskip.sty}{%
    \usepackage{parskip}
  }{% else
    \setlength{\parindent}{0pt}
    \setlength{\parskip}{6pt plus 2pt minus 1pt}}
}{% if KOMA class
  \KOMAoptions{parskip=half}}
\makeatother
% Make \paragraph and \subparagraph free-standing
\makeatletter
\ifx\paragraph\undefined\else
  \let\oldparagraph\paragraph
  \renewcommand{\paragraph}{
    \@ifstar
      \xxxParagraphStar
      \xxxParagraphNoStar
  }
  \newcommand{\xxxParagraphStar}[1]{\oldparagraph*{#1}\mbox{}}
  \newcommand{\xxxParagraphNoStar}[1]{\oldparagraph{#1}\mbox{}}
\fi
\ifx\subparagraph\undefined\else
  \let\oldsubparagraph\subparagraph
  \renewcommand{\subparagraph}{
    \@ifstar
      \xxxSubParagraphStar
      \xxxSubParagraphNoStar
  }
  \newcommand{\xxxSubParagraphStar}[1]{\oldsubparagraph*{#1}\mbox{}}
  \newcommand{\xxxSubParagraphNoStar}[1]{\oldsubparagraph{#1}\mbox{}}
\fi
\makeatother

\usepackage{color}
\usepackage{fancyvrb}
\newcommand{\VerbBar}{|}
\newcommand{\VERB}{\Verb[commandchars=\\\{\}]}
\DefineVerbatimEnvironment{Highlighting}{Verbatim}{commandchars=\\\{\}}
% Add ',fontsize=\small' for more characters per line
\usepackage{framed}
\definecolor{shadecolor}{RGB}{241,243,245}
\newenvironment{Shaded}{\begin{snugshade}}{\end{snugshade}}
\newcommand{\AlertTok}[1]{\textcolor[rgb]{0.68,0.00,0.00}{#1}}
\newcommand{\AnnotationTok}[1]{\textcolor[rgb]{0.37,0.37,0.37}{#1}}
\newcommand{\AttributeTok}[1]{\textcolor[rgb]{0.40,0.45,0.13}{#1}}
\newcommand{\BaseNTok}[1]{\textcolor[rgb]{0.68,0.00,0.00}{#1}}
\newcommand{\BuiltInTok}[1]{\textcolor[rgb]{0.00,0.23,0.31}{#1}}
\newcommand{\CharTok}[1]{\textcolor[rgb]{0.13,0.47,0.30}{#1}}
\newcommand{\CommentTok}[1]{\textcolor[rgb]{0.37,0.37,0.37}{#1}}
\newcommand{\CommentVarTok}[1]{\textcolor[rgb]{0.37,0.37,0.37}{\textit{#1}}}
\newcommand{\ConstantTok}[1]{\textcolor[rgb]{0.56,0.35,0.01}{#1}}
\newcommand{\ControlFlowTok}[1]{\textcolor[rgb]{0.00,0.23,0.31}{\textbf{#1}}}
\newcommand{\DataTypeTok}[1]{\textcolor[rgb]{0.68,0.00,0.00}{#1}}
\newcommand{\DecValTok}[1]{\textcolor[rgb]{0.68,0.00,0.00}{#1}}
\newcommand{\DocumentationTok}[1]{\textcolor[rgb]{0.37,0.37,0.37}{\textit{#1}}}
\newcommand{\ErrorTok}[1]{\textcolor[rgb]{0.68,0.00,0.00}{#1}}
\newcommand{\ExtensionTok}[1]{\textcolor[rgb]{0.00,0.23,0.31}{#1}}
\newcommand{\FloatTok}[1]{\textcolor[rgb]{0.68,0.00,0.00}{#1}}
\newcommand{\FunctionTok}[1]{\textcolor[rgb]{0.28,0.35,0.67}{#1}}
\newcommand{\ImportTok}[1]{\textcolor[rgb]{0.00,0.46,0.62}{#1}}
\newcommand{\InformationTok}[1]{\textcolor[rgb]{0.37,0.37,0.37}{#1}}
\newcommand{\KeywordTok}[1]{\textcolor[rgb]{0.00,0.23,0.31}{\textbf{#1}}}
\newcommand{\NormalTok}[1]{\textcolor[rgb]{0.00,0.23,0.31}{#1}}
\newcommand{\OperatorTok}[1]{\textcolor[rgb]{0.37,0.37,0.37}{#1}}
\newcommand{\OtherTok}[1]{\textcolor[rgb]{0.00,0.23,0.31}{#1}}
\newcommand{\PreprocessorTok}[1]{\textcolor[rgb]{0.68,0.00,0.00}{#1}}
\newcommand{\RegionMarkerTok}[1]{\textcolor[rgb]{0.00,0.23,0.31}{#1}}
\newcommand{\SpecialCharTok}[1]{\textcolor[rgb]{0.37,0.37,0.37}{#1}}
\newcommand{\SpecialStringTok}[1]{\textcolor[rgb]{0.13,0.47,0.30}{#1}}
\newcommand{\StringTok}[1]{\textcolor[rgb]{0.13,0.47,0.30}{#1}}
\newcommand{\VariableTok}[1]{\textcolor[rgb]{0.07,0.07,0.07}{#1}}
\newcommand{\VerbatimStringTok}[1]{\textcolor[rgb]{0.13,0.47,0.30}{#1}}
\newcommand{\WarningTok}[1]{\textcolor[rgb]{0.37,0.37,0.37}{\textit{#1}}}

\usepackage{longtable,booktabs,array}
\usepackage{calc} % for calculating minipage widths
% Correct order of tables after \paragraph or \subparagraph
\usepackage{etoolbox}
\makeatletter
\patchcmd\longtable{\par}{\if@noskipsec\mbox{}\fi\par}{}{}
\makeatother
% Allow footnotes in longtable head/foot
\IfFileExists{footnotehyper.sty}{\usepackage{footnotehyper}}{\usepackage{footnote}}
\makesavenoteenv{longtable}
\usepackage{graphicx}
\makeatletter
\newsavebox\pandoc@box
\newcommand*\pandocbounded[1]{% scales image to fit in text height/width
  \sbox\pandoc@box{#1}%
  \Gscale@div\@tempa{\textheight}{\dimexpr\ht\pandoc@box+\dp\pandoc@box\relax}%
  \Gscale@div\@tempb{\linewidth}{\wd\pandoc@box}%
  \ifdim\@tempb\p@<\@tempa\p@\let\@tempa\@tempb\fi% select the smaller of both
  \ifdim\@tempa\p@<\p@\scalebox{\@tempa}{\usebox\pandoc@box}%
  \else\usebox{\pandoc@box}%
  \fi%
}
% Set default figure placement to htbp
\def\fps@figure{htbp}
\makeatother





\setlength{\emergencystretch}{3em} % prevent overfull lines

\providecommand{\tightlist}{%
  \setlength{\itemsep}{0pt}\setlength{\parskip}{0pt}}



 


\KOMAoption{captions}{tableheading}
\makeatletter
\@ifpackageloaded{tcolorbox}{}{\usepackage[skins,breakable]{tcolorbox}}
\@ifpackageloaded{fontawesome5}{}{\usepackage{fontawesome5}}
\definecolor{quarto-callout-color}{HTML}{909090}
\definecolor{quarto-callout-note-color}{HTML}{0758E5}
\definecolor{quarto-callout-important-color}{HTML}{CC1914}
\definecolor{quarto-callout-warning-color}{HTML}{EB9113}
\definecolor{quarto-callout-tip-color}{HTML}{00A047}
\definecolor{quarto-callout-caution-color}{HTML}{FC5300}
\definecolor{quarto-callout-color-frame}{HTML}{acacac}
\definecolor{quarto-callout-note-color-frame}{HTML}{4582ec}
\definecolor{quarto-callout-important-color-frame}{HTML}{d9534f}
\definecolor{quarto-callout-warning-color-frame}{HTML}{f0ad4e}
\definecolor{quarto-callout-tip-color-frame}{HTML}{02b875}
\definecolor{quarto-callout-caution-color-frame}{HTML}{fd7e14}
\makeatother
\makeatletter
\@ifpackageloaded{bookmark}{}{\usepackage{bookmark}}
\makeatother
\makeatletter
\@ifpackageloaded{caption}{}{\usepackage{caption}}
\AtBeginDocument{%
\ifdefined\contentsname
  \renewcommand*\contentsname{Table of contents}
\else
  \newcommand\contentsname{Table of contents}
\fi
\ifdefined\listfigurename
  \renewcommand*\listfigurename{List of Figures}
\else
  \newcommand\listfigurename{List of Figures}
\fi
\ifdefined\listtablename
  \renewcommand*\listtablename{List of Tables}
\else
  \newcommand\listtablename{List of Tables}
\fi
\ifdefined\figurename
  \renewcommand*\figurename{Figure}
\else
  \newcommand\figurename{Figure}
\fi
\ifdefined\tablename
  \renewcommand*\tablename{Table}
\else
  \newcommand\tablename{Table}
\fi
}
\@ifpackageloaded{float}{}{\usepackage{float}}
\floatstyle{ruled}
\@ifundefined{c@chapter}{\newfloat{codelisting}{h}{lop}}{\newfloat{codelisting}{h}{lop}[chapter]}
\floatname{codelisting}{Listing}
\newcommand*\listoflistings{\listof{codelisting}{List of Listings}}
\usepackage{amsthm}
\theoremstyle{definition}
\newtheorem{definition}{Definition}[chapter]
\theoremstyle{plain}
\newtheorem{theorem}{Theorem}[chapter]
\theoremstyle{plain}
\newtheorem{proposition}{Proposition}[chapter]
\theoremstyle{definition}
\newtheorem{example}{Example}[chapter]
\theoremstyle{plain}
\newtheorem{corollary}{Corollary}[chapter]
\theoremstyle{plain}
\newtheorem{lemma}{Lemma}[chapter]
\theoremstyle{remark}
\AtBeginDocument{\renewcommand*{\proofname}{Proof}}
\newtheorem*{remark}{Remark}
\newtheorem*{solution}{Solution}
\newtheorem{refremark}{Remark}[chapter]
\newtheorem{refsolution}{Solution}[chapter]
\makeatother
\makeatletter
\makeatother
\makeatletter
\@ifpackageloaded{caption}{}{\usepackage{caption}}
\@ifpackageloaded{subcaption}{}{\usepackage{subcaption}}
\makeatother
\makeatletter
\@ifpackageloaded{algorithm}{}{\usepackage{algorithm}}
\makeatother
\makeatletter
\@ifpackageloaded{algpseudocode}{}{\usepackage{algpseudocode}}
\makeatother
\makeatletter
\@ifpackageloaded{caption}{}{\usepackage{caption}}
\makeatother
    \makeatletter
    \newcommand\fs@nocaption{
      \def\@fs@cfont{\bfseries}
      \let\@fs@capt\floatc@plain
      \def\@fs@pre{}%
      \def\@fs@post{\kern2pt\hrule}%
      \def\@fs@mid{\hrule\kern2pt}%
      \let\@fs@iftopcapt\iftrue}
    \makeatother
  
\captionsetup[algorithm]{justification=raggedright}
\usepackage{bookmark}
\IfFileExists{xurl.sty}{\usepackage{xurl}}{} % add URL line breaks if available
\urlstyle{same}
\hypersetup{
  pdftitle={Phy2Q Roadmap},
  pdfauthor={Albert Zhang},
  colorlinks=true,
  linkcolor={blue},
  filecolor={Maroon},
  citecolor={Blue},
  urlcolor={Blue},
  pdfcreator={LaTeX via pandoc}}


\title{Phy2Q Roadmap}
\author{Albert Zhang}
\date{2025-08-09}
\begin{document}
\maketitle

\floatname{algorithm}{Algorithm}

\numberwithin{algorithm}{chapter}

\renewcommand*\contentsname{Table of contents}
{
\hypersetup{linkcolor=}
\setcounter{tocdepth}{2}
\tableofcontents
}

\bookmarksetup{startatroot}

\chapter*{Preface}\label{preface}
\addcontentsline{toc}{chapter}{Preface}

\markboth{Preface}{Preface}

This is a Quarto book.

To learn more about Quarto books visit
\url{https://quarto.org/docs/books}.

\part{Numerical Optimization}

\chapter{Introduction}\label{introduction}

\chapter{Theory of Unconstrained
Optimization}\label{theory-of-unconstrained-optimization}

\chapter{Line Search Methods}\label{line-search-methods}

The line search methods choose a direction \(\mathbf{p}_k\) and search
along this direction for next point hoping the function value decree,
and continue this process until converge. In general, we hope to find
the step length s.t. \[
\min_{\alpha>0}f(\mathbf{x}_k+\alpha\mathbf{p}_k)
\]

\begin{refremark}
\leavevmode

\begin{itemize}
\tightlist
\item
  Through the minimization problem can be solved exactly, but this is
  typically not desired as it is expensive to find the exact solution.
  (methods attained \(\arg\min_{\alpha>0}\) are called \emph{exact line
  search})
\item
  Practically, the step length is approximately solved through limit
  number of trail steps tile meeting certain conditions, referred as
  \emph{inexact line search}.
\end{itemize}

\label{rem-lineSearch1}

\end{refremark}

For line search methods, the algorithm generally has two major steps:

\begin{enumerate}
\def\labelenumi{\arabic{enumi}.}
\item
  Pick descent direction
\item
  Choose appropriate step length \(\alpha\) to update
\end{enumerate}

We first discuss generally a descent direction, then in details about
step length selection.

\begin{tcolorbox}[enhanced jigsaw, bottomrule=.15mm, title=\textcolor{quarto-callout-tip-color}{\faLightbulb}\hspace{0.5em}{What is a descent direction}, left=2mm, colback=white, toprule=.15mm, bottomtitle=1mm, titlerule=0mm, leftrule=.75mm, colbacktitle=quarto-callout-tip-color!10!white, opacitybacktitle=0.6, coltitle=black, colframe=quarto-callout-tip-color-frame, rightrule=.15mm, breakable, toptitle=1mm, arc=.35mm, opacityback=0]

In general, we want the update \(x_{k+1}=x_k+\alpha_k\mathbf{p}_k\) to
always decrease the function value, for small enough step length
\(\alpha_k\).

\end{tcolorbox}

\begin{definition}[Descent
Direction]\protect\hypertarget{def-descentDir}{}\label{def-descentDir}

\(\mathbf{p}\in\mathbb{R}^n\) is called a descent direction for
\(f:\mathbb{R}^n\rightarrow\mathbb{R}\) at point
\(\mathbf{x}\in\mathbb{R}^n\) if \[
f(\mathbf{x} + \alpha p )<f(x),\quad\text{for sufficient small } \alpha>0
\]

\end{definition}

\begin{refremark}
\leavevmode

\begin{itemize}
\tightlist
\item
  The Definition~\ref{def-descentDir} is equivalent as stating
  \(\mathbf{p}^T\nabla f(x)<0\), which is often used to select descent
  direction in practice. This can be seen from the Taylor expansion of
  \(f(\mathbf{x+\alpha \mathbf{p}})\)
\end{itemize}

\label{rem-descentDir}

\end{refremark}

Once the descent direction has been selected, we still need to pick up a
appropriate step length \(\alpha\).

In the inexact line search methods, we need certain condition that stop
the iteration of finidng \(\alpha\). The simplest choice would be just
require a decrease in function value,
i.e.~\(f(\mathbf{x}_k + \alpha_k\mathbf{p_k})< f(\mathbf{x_k})\).
However, this condition is not sufficient to ensure the convergence to a
local minimizer.

\begin{example}[]\protect\hypertarget{exm-tooSmall}{}\label{exm-tooSmall}

example of too large step length here

\end{example}

The Example~\ref{exm-tooSmall} shows that the too small function value
drop should be avoid otherwise the algotihem would not be efficient.

\section{Sufficient Decrease
Condition}\label{sufficient-decrease-condition}

To avoid this issue, a \emph{sufficient decrease condition} can be
imposed. As the name suggested, the condition imposed the the drop of
function value for a specific chosed step length should be suffificently
large.

One popular measure of sufficient decrease is the \emph{Armijo
condition}, which state as follow

\begin{definition}[Armijo
Condition]\protect\hypertarget{def-armijo}{}\label{def-armijo}

For \(\mathbf{x}_k\in\mathbb{R}^n\) and descent direction
\(\mathbf{p}_k\in\mathbb{R}^n\), step length \(\alpha_k\) satisfied
Armijo condition if \[
f(\mathbf{x}_k+\alpha_k\mathbf{p}_k)\leq f(\mathbf{x}_k) + c_1\alpha_k\nabla^Tf(\mathbf{x}_k)\mathbf{p}_k,\quad \text{for some } c_1\in (0, 1)
\]

\end{definition}

\begin{tcolorbox}[enhanced jigsaw, bottomrule=.15mm, title=\textcolor{quarto-callout-tip-color}{\faLightbulb}\hspace{0.5em}{Idea behind}, left=2mm, colback=white, toprule=.15mm, bottomtitle=1mm, titlerule=0mm, leftrule=.75mm, colbacktitle=quarto-callout-tip-color!10!white, opacitybacktitle=0.6, coltitle=black, colframe=quarto-callout-tip-color-frame, rightrule=.15mm, breakable, toptitle=1mm, arc=.35mm, opacityback=0]

We can denote the LHS of the inequality in Definition~\ref{def-armijo}
as a function of \(\alpha\) as \(\phi(\alpha)\), and the RHS as a linear
function \(l(\alpha)\) with a negative slop
\(c_1\nabla^Tf(\mathbf{x}_k)\mathbf{p}_k\). Note that this slop is
between the horizontal line and the tangent line at \(\phi(0)\). The
\emph{Armijo condition} is thus requiring the selected step length drop
the function value larger than a predifined linear function, which
derease slower than the steepest local function drop.

\end{tcolorbox}

Figure~\ref{fig-armijo} shows the acceptable region of the \emph{Armijo
line search}.

\begin{figure}

\centering{

\pandocbounded{\includegraphics[keepaspectratio]{numerical-optimization/lineSearch_files/figure-pdf/fig-armijo-output-1.pdf}}

}

\caption{\label{fig-armijo}Sufficient-decrease (Armijo) condition:
\(\phi(\alpha)=f(\mathbf{x}_k+\alpha\mathbf{p}_k)\) (blue), the tangent
at \(\alpha=0\) with slope
\(\phi^\prime(0)=\nabla^{\mathsf T}f_k\,\mathbf{p}_k\) (dotted), and the
sufficient-decrease line
\(l(\alpha)=f(\mathbf{x}_k)+c_1\alpha\,\nabla^{\mathsf T}f_k\,\mathbf{p}_k\)
(dashed). Armijo accepts the step lengths where
\(\phi(\alpha)\le l(\alpha)\) (shaded).}

\end{figure}%

In practice, the \emph{Armijo condition} can be satisified using
\emph{Backtracing line search} algorithm

\subsection{Backtracking Algorithm}\label{backtracking-algorithm}

The idea of the backtracking line search is to start with a relatively
large step length and shrink \(\alpha\) every iteration till the
\emph{Armijo condition} is satisfied.

\begin{Shaded}
\begin{Highlighting}[]
\ImportTok{import}\NormalTok{ numpy }\ImportTok{as}\NormalTok{ np}
\ImportTok{from}\NormalTok{ numpy.typing }\ImportTok{import}\NormalTok{ NDArray}
\ImportTok{from}\NormalTok{ typing }\ImportTok{import}\NormalTok{ Callable}

\KeywordTok{def}\NormalTok{ backtracking\_search(}
\NormalTok{    f: Callable,}
\NormalTok{    gradf\_xk: NDArray,}
\NormalTok{    xk: NDArray,}
\NormalTok{    pk: NDArray,}
\NormalTok{    alpha0: }\BuiltInTok{float}\NormalTok{,}
\NormalTok{    c1: }\BuiltInTok{float} \OperatorTok{=} \FloatTok{1e{-}4}\NormalTok{,}
\NormalTok{    tau: }\BuiltInTok{float} \OperatorTok{=} \FloatTok{0.5}\NormalTok{,}
\NormalTok{) }\OperatorTok{{-}\textgreater{}} \BuiltInTok{float}\NormalTok{:}
    \CommentTok{"""Step length search using Armijo backtracking methods}

\CommentTok{    Args:}
\CommentTok{        f (Callable): objective function}
\CommentTok{        gradf (Callable): gradient of objective function}
\CommentTok{        xk (NDArray): current position}
\CommentTok{        pk (NDArray): descent direction}
\CommentTok{        alpha0 (float): initial step length. Default to 1. }
\CommentTok{        c1 (float): control variable for Armijo condition. Default to 1e{-}5}
\CommentTok{        tau (float): contraction coefficient for step length. Default to 0,5. }

\CommentTok{    Returns:}
\CommentTok{        float: final step length}
\CommentTok{    """}
\NormalTok{    alpha }\OperatorTok{=}\NormalTok{ alpha0}
\NormalTok{    f\_xk }\OperatorTok{=}\NormalTok{ f(xk)}
    \ControlFlowTok{while}\NormalTok{ f(xk }\OperatorTok{+}\NormalTok{ alpha }\OperatorTok{*}\NormalTok{ pk) }\OperatorTok{\textgreater{}}\NormalTok{ f\_xk }\OperatorTok{+}\NormalTok{ c1 }\OperatorTok{*}\NormalTok{ alpha }\OperatorTok{*}\NormalTok{ np.dot(gradf\_xk, pk):}
\NormalTok{        alpha }\OperatorTok{=}\NormalTok{ tau }\OperatorTok{*}\NormalTok{ alpha}
    \ControlFlowTok{return}\NormalTok{ alpha}
\end{Highlighting}
\end{Shaded}

\begin{tcolorbox}[enhanced jigsaw, bottomrule=.15mm, title=\textcolor{quarto-callout-note-color}{\faInfo}\hspace{0.5em}{Note}, left=2mm, colback=white, toprule=.15mm, bottomtitle=1mm, titlerule=0mm, leftrule=.75mm, colbacktitle=quarto-callout-note-color!10!white, opacitybacktitle=0.6, coltitle=black, colframe=quarto-callout-note-color-frame, rightrule=.15mm, breakable, toptitle=1mm, arc=.35mm, opacityback=0]

\begin{itemize}
\tightlist
\item
  The initial guess of step length typically is 1, but might vary depens
  on \(\mathbf{p}_k\)
\item
  Require only function evaluation during the search
\item
  Will terminated in finite steps
\item
  The backtracking approach avoid too small step length as it star with
  large step and will usually stop before the step length becomes too
  small.
\end{itemize}

\end{tcolorbox}

Though the \emph{Armijo condition} ensure the sufficient decrease of the
function value. However, the selected step length could be too small and
leading to inefficient progress. The following \emph{curvature
condition} make sure the the step length would not be too short.

\section{Curvature Condition}\label{curvature-condition}

\begin{definition}[Curvature
Condition]\protect\hypertarget{def-curvature}{}\label{def-curvature}

For \(\mathbf{x}_k\in\mathbb{R}^n\) and descent direction
\(\mathbf{p}_k\in\mathbb{R}^n\), step length \(\alpha_k\) satisfied
curvature condition if \[
\nabla^Tf(\mathbf{x}_k+\alpha_k\mathbf{p}_k)\mathbf{p_k}\geq c_2\nabla f(\mathbb{x}_k),\quad c_2\in(c_1, 1)
\] \(c_1\) is the constant from Definition~\ref{def-armijo}

\end{definition}

\begin{tcolorbox}[enhanced jigsaw, bottomrule=.15mm, title=\textcolor{quarto-callout-tip-color}{\faLightbulb}\hspace{0.5em}{Idea behind}, left=2mm, colback=white, toprule=.15mm, bottomtitle=1mm, titlerule=0mm, leftrule=.75mm, colbacktitle=quarto-callout-tip-color!10!white, opacitybacktitle=0.6, coltitle=black, colframe=quarto-callout-tip-color-frame, rightrule=.15mm, breakable, toptitle=1mm, arc=.35mm, opacityback=0]

Again if we consider \(\phi(\alpha) = f(\mathbf{x} + \alpha\mathbf{p})\)
as a function if \(\alpha\). Then the LHS of the inequality in
Definition~\ref{def-curvature} is \(\phi^\prime(\alpha_k)\), which is
the slope of the tangent line at the accepted step length. Since the RHS
is negative, the curvature condition requires the curvature of \(\phi\)
at the accepted step length should be ``flat'' comparing to that at the
starting point, which is measured by the coefficient \(c_2\). Thus, the
accepted step length can not be too close to the staring point

\end{tcolorbox}

Figure~\ref{fig-curvature} shows the acceptable region satisfying the
curvature condition.

\begin{figure}

\centering{

\pandocbounded{\includegraphics[keepaspectratio]{numerical-optimization/lineSearch_files/figure-pdf/fig-curvature-output-1.pdf}}

}

\caption{\label{fig-curvature}Curvature condition: the accepted step
lengths (shaded) are those where the slope
\(\phi^\prime(\alpha)\ge c_2\,\phi^\prime(0)\) is flat enough, i.e.~not
too steeply negative. Short segments show the slope of \(\phi\) at the
two ends of the accepted interval.}

\end{figure}%

Combine the two condition together, we get the \emph{Wolf conditions}

\section{Wolf Condition}\label{wolf-condition}

\begin{definition}[Wolf
Condition]\protect\hypertarget{def-wolf}{}\label{def-wolf}

For \(\mathbf{x}_k\in\mathbb{R}^n\) and descent direction
\(\mathbf{p}_k\in\mathbb{R}\), step length \(\alpha_k\) satisfied Wolf
condition if \[
\begin{aligned}
f(\mathbf{x}_k+\alpha_k\mathbf{p}_k)&\leq f(\mathbf{x}_k) + c_1\alpha_k\nabla^Tf(\mathbf{x}_k)\mathbf{p}_k\\
\nabla^Tf(\mathbf{x}_k+\alpha_k\mathbf{p}_k)\mathbf{p_k}&\geq c_2\nabla f(\mathbb{x}_k)
\end{aligned}
\] for some \(0<c_1<c_2<1\)

\end{definition}

\begin{tcolorbox}[enhanced jigsaw, bottomrule=.15mm, title=\textcolor{quarto-callout-note-color}{\faInfo}\hspace{0.5em}{Note}, left=2mm, colback=white, toprule=.15mm, bottomtitle=1mm, titlerule=0mm, leftrule=.75mm, colbacktitle=quarto-callout-note-color!10!white, opacitybacktitle=0.6, coltitle=black, colframe=quarto-callout-note-color-frame, rightrule=.15mm, breakable, toptitle=1mm, arc=.35mm, opacityback=0]

\begin{itemize}
\tightlist
\item
  If \(c_1\geq c_2\), the exsistance of Wolf condition might not be
  satisfied
\item
  The Wolf condition requires evaluations of both function and gradient
  values at each iteration
\end{itemize}

\end{tcolorbox}

Figure~\ref{fig-wolf} shows the acceptable region satisfying the
\emph{Wolf condition}, which is the intersection of the sufficient
decrease and curvature conditions.

\begin{figure}

\centering{

\pandocbounded{\includegraphics[keepaspectratio]{numerical-optimization/lineSearch_files/figure-pdf/fig-wolf-output-1.pdf}}

}

\caption{\label{fig-wolf}Wolf condition: the accepted region (green) is
the intersection of the Armijo/sufficient-decrease region
(\(\phi(\alpha)\le l(\alpha)\)) and the curvature region
(\(\phi^\prime(\alpha)\ge c_2\phi^\prime(0)\)).}

\end{figure}%

The \emph{Wolf condition} can be strengthened to the \emph{strong Wolf
condition}.

\begin{definition}[Strong Wolf
Condition]\protect\hypertarget{def-strong-wolf}{}\label{def-strong-wolf}

Change the curvature condition in Definition~\ref{def-wolf} to \[
|\nabla^{\mathsf T}f(\mathbf{x}_k+\alpha_k\mathbf{p}_k)\mathbf{p}_k| \leq c_2\,|\nabla^{\mathsf T}f(\mathbf{x}_k)\mathbf{p}_k|
\]

\end{definition}

This restrict too large positive slope, which is also away from minimum.

\subsection{Line Search Algorithm Satisfying Wolf
Condition}\label{line-search-algorithm-satisfying-wolf-condition}

\begin{Shaded}
\begin{Highlighting}[]
\KeywordTok{def}\NormalTok{ wolf\_search(}
\NormalTok{        f: Callable[[NDArray], }\BuiltInTok{float}\NormalTok{],}
\NormalTok{        gradf: Callable[[NDArray], NDArray],}
\NormalTok{        xk: NDArray,}
\NormalTok{        pk: NDArray,}
\NormalTok{        c1: }\BuiltInTok{float}\OperatorTok{=}\FloatTok{1e{-}4}\NormalTok{,}
\NormalTok{        c2: }\BuiltInTok{float}\OperatorTok{=}\FloatTok{0.9}\NormalTok{,}
\NormalTok{        max\_iter: }\BuiltInTok{int}\OperatorTok{=}\DecValTok{1000}
\NormalTok{) }\OperatorTok{{-}\textgreater{}} \BuiltInTok{float}\NormalTok{:}
    \CommentTok{"""wolf line search methods}

\CommentTok{    Args:}
\CommentTok{        f (Callable[[NDArray], float]): objective function}
\CommentTok{        gradf (Callable[[NDArray], float]): gradient of objective function}
\CommentTok{        xk (NDArray): current position for starting the line search}
\CommentTok{        pk (NDArray): current search direction}
\CommentTok{        c1 (float, optional): parameter for Amijor contion. Defaults to 1e{-}4.}
\CommentTok{        c2 (float, optional): parameter for curvature condition. Defaults to 0.9.}
\CommentTok{        max\_iter(float, optional): maximum iteration of line search. Defaults to 1000.}
\CommentTok{    Returns:}
\CommentTok{        float: return the step length}
\CommentTok{    """}    
    
\NormalTok{    alpha }\OperatorTok{=} \DecValTok{1}
\NormalTok{    alpha\_l }\OperatorTok{=} \DecValTok{0}
\NormalTok{    alpha\_u }\OperatorTok{=}\NormalTok{ np.inf}

\NormalTok{    f\_xk }\OperatorTok{=}\NormalTok{ f(xk)}
\NormalTok{    gradf\_xk }\OperatorTok{=}\NormalTok{ gradf(xk)}
\NormalTok{    phi\_prime\_0 }\OperatorTok{=}\NormalTok{ np.dot(gradf\_xk, pk)}

    \CommentTok{\# check c1 \textless{} c2:}
    \ControlFlowTok{if}\NormalTok{ c1 }\OperatorTok{\textgreater{}}\NormalTok{ c2:}
        \ControlFlowTok{raise} \PreprocessorTok{ValueError}\NormalTok{(}\StringTok{"c1 must smaller than c2"}\NormalTok{)}
    
    \CommentTok{\# if the input search direction is not descent direction, just return}
    \ControlFlowTok{if}\NormalTok{ phi\_prime\_0 }\OperatorTok{\textgreater{}=} \DecValTok{0}\NormalTok{:}
        \ControlFlowTok{return} \DecValTok{1}

\NormalTok{    counter }\OperatorTok{=} \DecValTok{0}
    \ControlFlowTok{while} \VariableTok{True}\NormalTok{:}
        \CommentTok{\# break the line search if the iteration take to much}
        \ControlFlowTok{if}\NormalTok{ counter }\OperatorTok{\textgreater{}}\NormalTok{ max\_iter:}
            \ControlFlowTok{return}\NormalTok{ alpha}
        \CommentTok{\# if armijo condition not satisfied, shrink the upper bond of allowed step length}
        \ControlFlowTok{if}\NormalTok{ f(xk }\OperatorTok{+}\NormalTok{ alpha }\OperatorTok{*}\NormalTok{ pk) }\OperatorTok{\textgreater{}}\NormalTok{ f\_xk }\OperatorTok{+}\NormalTok{ c1 }\OperatorTok{*}\NormalTok{ alpha }\OperatorTok{*}\NormalTok{ phi\_prime\_0:}
\NormalTok{            alpha\_u }\OperatorTok{=}\NormalTok{ alpha}

        \ControlFlowTok{else}\NormalTok{:}

            \CommentTok{\# if armojo satsified but curvature condition is not, increase the lower bond of allowed step length}
            \ControlFlowTok{if}\NormalTok{ np.dot(gradf(xk }\OperatorTok{+}\NormalTok{ alpha }\OperatorTok{*}\NormalTok{ pk), pk) }\OperatorTok{\textless{}}\NormalTok{ c2 }\OperatorTok{*}\NormalTok{ phi\_prime\_0:}
\NormalTok{                alpha\_l }\OperatorTok{=}\NormalTok{ alpha}
            \ControlFlowTok{else}\NormalTok{:}
                \ControlFlowTok{return}\NormalTok{ alpha}
        
        \CommentTok{\# if the Armijo condition has not been violated once, increase the lower bound}
        \ControlFlowTok{if}\NormalTok{ alpha\_u }\OperatorTok{\textless{}}\NormalTok{ np.inf:}
\NormalTok{            alpha }\OperatorTok{=}\NormalTok{ (alpha\_l }\OperatorTok{+}\NormalTok{ alpha\_u) }\OperatorTok{/} \DecValTok{2}
        \ControlFlowTok{else}\NormalTok{:}
\NormalTok{            alpha }\OperatorTok{=} \DecValTok{2} \OperatorTok{*}\NormalTok{ alpha}
\NormalTok{        counter }\OperatorTok{+=} \DecValTok{1}
\end{Highlighting}
\end{Shaded}

\section{Existence and Convergence of Line
Search}\label{existence-and-convergence-of-line-search}

We now discuss when the step-length conditions can actually be
satisfied.

\begin{theorem}[Existence of Armijo
Condition]\protect\hypertarget{thm-armijo-existence}{}\label{thm-armijo-existence}

If \(f:\mathbb{R}^n\rightarrow\mathbb{R}\) is \(C^1\) and bounded below,
then for a descent direction \(\mathbf{p}_k\) at \(\mathbf{x}_k\), there
exists \(\alpha^\prime>0\) s.t. the Armijo condition is satisfied for
\(\forall\alpha\in[0, \alpha_{\max}]\).

In addition, if \(\nabla f\) is \(L\)-Lipschitz continuous, then \[
\alpha^\prime = \frac{2(c_1-1)\,\nabla^{\mathsf T}f(\mathbf{x}_k)\mathbf{p}_k}{L\|\mathbf{p}_k\|^2}
\]

\end{theorem}

\begin{refremark}
\leavevmode

\begin{itemize}
\tightlist
\item
  Notice that \(\phi(0)+c_1\alpha\phi^\prime(0)\) is a line go to
  unbounded below, thus \(l(\alpha)\) must intersect with
  \(\phi(\alpha)\), thus there exists \(\alpha_{\max}\) which is the
  smallest intersection, s.t. \(\alpha^\prime<\alpha_{\max}\) is the one
  we need.
\item
  To show the second part, can use inequality \[
  f(\mathbf{y}) \leq f(\mathbf{x}) + \nabla^{\mathsf T}f(\mathbf{x})(\mathbf{y}-\mathbf{x}) + \frac{L}{2}\|\mathbf{y}-\mathbf{x}\|^2
  \]
\end{itemize}

\label{rem-armijo-existence}

\end{refremark}

\begin{theorem}[Existence of Wolf
Condition]\protect\hypertarget{thm-wolf-existence}{}\label{thm-wolf-existence}

If \(f\) is \(C^1\) and is bounded below, for a descent direction at
\(\mathbf{x}_k\), then there exists an interval for \(\alpha_k\) s.t.
the Wolf condition is satisfied.

\end{theorem}

\begin{refremark}
\leavevmode

\begin{itemize}
\tightlist
\item
  The existence for Armijo condition show the exist interval bounded
  above by line \(l(\alpha)\), use Taylor's theorem can show tangent
  line.
\item
  Continuity of \(f\) can show the existence of interval satisfying Wolf
  condition.
\end{itemize}

\label{rem-wolf-existence}

\end{refremark}

\chapter{Gradient Descent Methods}\label{gradient-descent-methods}

\chapter{Newton's Methods}\label{newtons-methods}

\chapter{Quasi Newton Methods}\label{quasi-newton-methods}

\textbf{Newton Method}: Require calculation of Hessian matrix
\(H_{ij}=\frac{\partial^2f}{\partial x_i\partial x_j}\)

\textbf{Steepest Descent}: Find the descent direction using local
quadratic approximation with identity matrix as Hessian. \[
\underset{\mathrm{argmin}}{\mathbf{p}\in\mathbb{R}^n}f(\mathbf{x}_k)+\nabla f^T(\mathbf{x}_k)\mathbf{p}+1/2\mathbf{p}^TI\mathbf{p}
\]

\begin{tcolorbox}[enhanced jigsaw, bottomrule=.15mm, title=\textcolor{quarto-callout-tip-color}{\faLightbulb}\hspace{0.5em}{Idea of Quasi-Newton Method}, left=2mm, colback=white, toprule=.15mm, bottomtitle=1mm, titlerule=0mm, leftrule=.75mm, colbacktitle=quarto-callout-tip-color!10!white, opacitybacktitle=0.6, coltitle=black, colframe=quarto-callout-tip-color-frame, rightrule=.15mm, breakable, toptitle=1mm, arc=.35mm, opacityback=0]

Instead of using precise second-order derivative information to
construct quadratic approximation, use a matrix \(B_k\) that approximate
the Hessian and update at each iteration

\end{tcolorbox}

\section{Generic Approach of Quasi-Newton
Methods}\label{generic-approach-of-quasi-newton-methods}

\subsection{Quadratic Objective
Function}\label{quadratic-objective-function}

The descent direction \(\mathbf{p}_k\) at each iteration is found by
finding the minimizer of the following qudratic approximation of the
objective function. \[
\min_{\mathbf{p}}f_k+\nabla f_k^T\mathbf{p} + \frac{1}{2}\mathbf{p}^TB_k\mathbf{p}
\]

\begin{tcolorbox}[enhanced jigsaw, bottomrule=.15mm, title=\textcolor{quarto-callout-note-color}{\faInfo}\hspace{0.5em}{Approximate Hessian \(B_k\)}, left=2mm, colback=white, toprule=.15mm, bottomtitle=1mm, titlerule=0mm, leftrule=.75mm, colbacktitle=quarto-callout-note-color!10!white, opacitybacktitle=0.6, coltitle=black, colframe=quarto-callout-note-color-frame, rightrule=.15mm, breakable, toptitle=1mm, arc=.35mm, opacityback=0]

\begin{itemize}
\tightlist
\item
  symmetric (a Hessian matrix would be symmetric)
\item
  Positive definite (to make sure existence of minimizer)
\end{itemize}

\end{tcolorbox}

\begin{refremark}
The minimizer can be found as \(\mathbf{p}_k=-B^{-1}_k\nabla f_k\),
which would be a descent direction if \(B_k\) is PD.

\label{rem-quasiDescent}

\end{refremark}

\begin{tcolorbox}[enhanced jigsaw, bottomrule=.15mm, title=\textcolor{quarto-callout-note-color}{\faInfo}\hspace{0.5em}{Note}, left=2mm, colback=white, toprule=.15mm, bottomtitle=1mm, titlerule=0mm, leftrule=.75mm, colbacktitle=quarto-callout-note-color!10!white, opacitybacktitle=0.6, coltitle=black, colframe=quarto-callout-note-color-frame, rightrule=.15mm, breakable, toptitle=1mm, arc=.35mm, opacityback=0]

Remark~\ref{rem-quasiDescent} can be verified using
Remark~\ref{rem-descentDir}.

\end{tcolorbox}

\subsection{Secant Condition}\label{secant-condition}

The idea of \emph{secant} condition in quasi Newton method is to
approximate the Hessian matrix with gradient difference. In 1D case,
this is simply

\[
\frac{\mathrm{d}^2f}{\mathrm{d}^2x}\approx\frac{\left.\frac{\mathrm{d}f}{\mathrm{d}x}\right|_{x_{k+1}} - \left.\frac{\mathrm{d}f}{\mathrm{d}x}\right|_{x_k}}{x_{k+1}-x_k}
\]

Generalized to multi-variable case, using Taylor expansion of
\(\nabla f(\mathbf{x}_{k+1})\) \[
\nabla f(\mathbf{x}_{k+1}) \approx \nabla f(\mathbf{x}_k) + \nabla ^2 f(\mathbf{x}_k)(\mathbf{x}_{k+1} - \mathbf{x}_k)
\] therefor, we can approximate the Hessian with matrix
\(B_{k+1}\in\mathbb{R}^n\) satisfy the \emph{secant condition} \[
\underbrace{\nabla f(\mathbf{x}_{k+1}) - \nabla f(\mathbf{x}_k)}_{\mathbf{y}_k} = B_{k+1}\underbrace{(\mathbf{x}_{k+1}-\mathbf{x}_k)}_{\mathbf{s}_k} 
\]

\subsection{Positive Definite
Condition}\label{positive-definite-condition}

As we mentioned before, \(B_{k+1}\) has to be positive definite so that
the quadratic objective function has minimizer and \(\mathbf{p}_k\)
found is a descent direction. If the quasi Newton methods shown later
(DFP and BFGS), it can be shown that if previous step \(B_{k}\) is PD,
then \(B_{k+1}\) is PD if \[
\mathbf{s}^T\mathbf{y}_k>0
\]

\subsection{Additional Condition}\label{additional-condition}

It is obvouse that \(B_k\) should be symmeteric, i.e.~\(B^T_k = B_k\).
However, the \(B_k\) satisfied the secant condition and symmetric is not
unique.

To uniquely detemines \(B_k\), additional condition should be imposed
and one resonable would be requring the \(B_{k+1}\), the Hessian for
next interation be a symmetric matrix satisfying secant condition and
also closest to previouse Hessian in some sence, i.e. \[
B_{k+1} = \mathrm{argmin}\lVert B - B_k \rVert, 
\] where \(\lVert\cdot\rVert\) be some matrix norm.

\subsection{Summary}\label{summary}

The quasi Newton methods search the descent direction for next step by
approximate the function locally with quadratic function

\chapter{Conjugate Gradient Methods}\label{conjugate-gradient-methods}

\chapter{Theory of Constrained
Optimization}\label{theory-of-constrained-optimization}

\section{COnstrained Optimization
Problem}\label{constrained-optimization-problem}

The constrained optimization problem can be formulated as follow

\section{Some Terminology}\label{some-terminology}

\section{Optimmality COndition}\label{optimmality-condition}

\subsection{First Order COnsideration}\label{first-order-consideration}

\subsection{Second Order
Consideration}\label{second-order-consideration}

\chapter{Quadratic Programming}\label{quadratic-programming}

The Quadratic Programming solves constrain optimization problem with
quadratic objective and linear constrains

\section{Equality Contsrained Quadratic Programming
(EQP)}\label{equality-contsrained-quadratic-programming-eqp}

The first order KKT condition of the

\chapter{Sequential Quadratic
Programming}\label{sequential-quadratic-programming}

Sequential Quadratic Programming (SQP) is extended method of quadratic
programming to solve general constrained optimization problem using
Newton's method and QP at each iterations

\section{Equality COnstrained Case}\label{equality-constrained-case}

\section{Inequality COnstrained Case}\label{inequality-constrained-case}

\chapter{Penalty Methods}\label{penalty-methods}

\chapter{Interior Point Methods}\label{interior-point-methods}

\part{Monte-Carlo Methods}

\chapter{Monte-Carlo Method}\label{monte-carlo-method}

Monte-Carlo method play important role

\part{Machine Learning and Data Science Basics}

\chapter{Recap on Probability}\label{recap-on-probability}

\section{Concentration Inequalities}\label{concentration-inequalities}

The guiding theme: a sample average concentrates around its expectation,
i.e.

\[
\Big|\,\tfrac{1}{n}\sum_{i} X_i - \mathbb{E}\big[\tfrac{1}{n}\sum_{i=1}^{n} X_i\big]\,\Big|
\quad \text{is small, guaranteed with high probability.}
\]

\section{Markov's Inequality}\label{markovs-inequality}

\begin{theorem}[Markov's
inequality]\protect\hypertarget{thm-markov}{}\label{thm-markov}

Let \(X\) be a non-negative random variable. Then for all \(a > 0\), \[
P(X \ge a) \le \frac{\mathbb{E}[X]}{a}.
\]

\end{theorem}

\begin{remark}
The probability of \(X\) being ``too large'', \(P(X \ge a)\), is
bounded.
\end{remark}

\begin{proof}
If \(X\) is a continuous R.V. with pdf \(f_X(x)\), then \[
\mathbb{E}(X) = \int_0^{+\infty} x\, f_X(x)\,dx
= \int_0^{a} x\, f_X(x)\,dx + \int_a^{\infty} x\, f_X(x)\,dx
\ge \int_a^{\infty} x\, f_X(x)\,dx
\ge a \int_a^{\infty} f_X(x)\,dx
= a\, P(X \ge a).
\]
\end{proof}

\begin{remark}
\leavevmode

\begin{itemize}
\tightlist
\item
  \(+\) : almost always works (\(X\) non-negative).
\item
  \(-\) : the bound is weak (not precise).
\end{itemize}

\textbf{Power trick \(\to\) sharper bound:} \[
P(X^r \ge a^r) \le \frac{\mathbb{E}(X^r)}{a^r}
\qquad (\text{assuming the } r\text{-th moment } \mathbb{E}(X^r) \text{ exists}).
\]
\end{remark}

\section{Chebyshev's Inequality}\label{chebyshevs-inequality}

\begin{theorem}[Chebyshev's
inequality]\protect\hypertarget{thm-chebyshev}{}\label{thm-chebyshev}

Let \(X\) be a random variable. Then for all \(c > 0\), \[
P\big(|X - \mathbb{E}(X)| \ge c\big) \le \frac{\operatorname{Var}(X)}{c^2}.
\]

\end{theorem}

\begin{proof}
Apply the power trick of Markov's inequality: \[
P\big(|X - \mathbb{E}(X)| \ge c\big)
= P\big(|X - \mathbb{E}(X)|^2 \ge c^2\big)
\le \frac{\mathbb{E}\big[(X - \mathbb{E}(X))^2\big]}{c^2}
= \frac{\operatorname{Var}(X)}{c^2}.
\]
\end{proof}

\section{Master Tail Bound}\label{master-tail-bound}

\begin{theorem}[Master tail
bound]\protect\hypertarget{thm-master-tail}{}\label{thm-master-tail}

Let \(X = X_1 + X_2 + \cdots + X_n\), where \(X_1, \dots, X_n\) are
i.i.d. and

\begin{itemize}
\tightlist
\item
  \(\mathbb{E}(X_i) = 0\),
\item
  \(\mathbb{E}(X_i^2) \le \sigma^2\),
\item
  \(\big|\mathbb{E}(X_i^s)\big| \le \sigma^2\, s!\) for
  \(s = 3, 4, \dots, \big\lfloor \tfrac{a^2}{4n\sigma^2}\big\rfloor\).
\end{itemize}

Then, for \(0 \le a \le \sqrt{2}\, n\sigma^2\), \[
P(|X| \ge a) \le 3\, e^{-\frac{a^2}{12 n \sigma^2}}.
\]

\end{theorem}

\begin{proof}
\textbf{(Sketch --- the full proof is deferred to an appendix.)}

\textbf{(1)} Choose \(r\) an even integer such that \(0 < r \le s\).
Then \[
P(|X| \ge a) = P(X^r \ge a^r) \le \frac{\mathbb{E}(X^r)}{a^r}.
\]

\textbf{(2)} By the multinomial theorem and independence, \[
\mathbb{E}(X^r) = \mathbb{E}\Big[(X_1 + \cdots + X_n)^r\Big]
= \sum_{\substack{r_1 + \cdots + r_n = r \\ r_i \ge 0}} \frac{r!}{r_1!\, r_2! \cdots r_n!}\, \mathbb{E}\big(X_1^{r_1} \cdots X_n^{r_n}\big)
= \sum_{\substack{r_1 + \cdots + r_n = r \\ r_i \ge 0}} \frac{r!}{r_1! \cdots r_n!}\, \prod_i \mathbb{E}\big(X_i^{r_i}\big).
\]

\textbf{(3)} Since \(\mathbb{E}(X_i) = 0\),
\(\mathbb{E}(X_i^2) \le \sigma^2 \le 2\sigma^2\), and
\(\mathbb{E}(X_i^{r_i}) \le r_i!\, \sigma^2\), any term with some
\(r_i = 1\) vanishes, so the sum is restricted to indices with
\(r_i = 0\) or \(r_i \ge 2\). Let \(t = \#\{i : r_i \neq 0\}\); then \[
1 \le t \le \tfrac{r}{2} \qquad (\text{each non-zero } r_i \text{ is at least } 2).
\]

\textbf{(4)} Enumerate the multi-indices by \(t\): for a fixed \(t\)
there are \(\binom{n}{t}\) ways to choose which \(r_i\) are non-zero,
and (after that choice) \(\binom{r-t-1}{t-1}\) ways to distribute the
remaining mass.

\textbf{(5)} Hence, with \(\sigma^{2t}\) from the moment bounds, \[
\mathbb{E}(X^r) \le r! \sum_{t=1}^{r/2} \binom{n}{t}\binom{r-t-1}{t-1} \sigma^{2t},
\qquad \text{using } \binom{n}{t} \le \frac{n^t}{t!}, \quad \binom{r-t-1}{t-1} \le 2^{\,r-t-1}.
\]

\textbf{(6)} This leads to \[
\mathbb{E}(X^r) \le \Big(\frac{2\, r\, n\, \sigma^2}{a^2}\Big)^{r/2} =: g(r),
\] and \(g(r)\) decreases while \(r - 1 \le \tfrac{a^2}{4n\sigma^2}\).
Choosing \(r\) to be the largest even integer with
\(r \le \tfrac{a^2}{6 n \sigma^2}\) gives \[
P(|X| \ge a) \le \big(\tfrac{1}{3}\big)^{r/2} \le e^{-r/2} \le e \cdot e^{-\frac{a^2}{12 n \sigma^2}} \le 3\, e^{-\frac{a^2}{12 n \sigma^2}}.
\]
\end{proof}

\begin{remark}
The master tail bound is a general probability tail bound for a sum of
i.i.d. random variables that gives tighter concentration than the CLT.
Tail bounds for various distributions can be derived from it.
\end{remark}

\section{Chernoff Bound}\label{chernoff-bound}

\begin{example}[Chernoff bound (i.i.d.
Bernoulli)]\protect\hypertarget{exm-chernoff}{}\label{exm-chernoff}

Let
\(Y_1, \dots, Y_n \overset{\text{iid}}{\sim} \mathrm{Bernoulli}(p)\),
i.e.~\(Y_i = \begin{cases} 1, & p \\ 0, & 1-p \end{cases}\), and let
\(Y = Y_1 + \cdots + Y_n\). Then for all \(c \in [0,1]\), \[
P\Big[\,\big|\tfrac{Y}{n} - p\big| \ge \sqrt{2}\, c\, p(1-p)\,\Big] \le 3\, e^{-\frac{n\, c^2\, p(1-p)}{6}}.
\]

\end{example}

\begin{proof}
Apply the master tail bound to \(X_i = Y_i - p\), so
\(\mathbb{E}(X_i) = 0\) and \[
\mathbb{E}(X_i^2) = \mathbb{E}\big((Y_i - p)^2\big) = \mathbb{E}(Y_i^2) + p^2 - 2p\,\mathbb{E}(Y_i) = \mathbb{E}(Y_i^2) - \mathbb{E}(Y_i)^2 = \operatorname{Var}(Y_i),
\] \[
\big|\mathbb{E}(X_i^s)\big| = \big|\, p(1-p)^s + (1-p)(-p)^s \,\big| \le p(1-p) = \operatorname{Var}(Y_i).
\] The Chernoff bound is then obtained directly from the master tail
bound with \(a = n\sqrt{2}\, c\, p(1-p)\).
\end{proof}

\section{Concentration via the Moment Generating
Function}\label{concentration-via-the-moment-generating-function}

Instead of applying Markov's inequality to \(|X|^r\), apply it to
\(e^{\lambda X}\): \[
P(X \ge t) = P\big(e^{\lambda X} \ge e^{\lambda t}\big) \le \frac{\mathbb{E}[e^{\lambda X}]}{e^{\lambda t}}, \qquad \forall\, \lambda > 0
\;\;\Longrightarrow\;\;
P(X \ge t) \le \inf_{\lambda > 0} e^{-\lambda t}\, \mathbb{E}[e^{\lambda X}],
\] noting that \(\mathbb{E}[e^{\lambda X}]\) is the moment generating
function (MGF).

\begin{theorem}[Tail bound for a
Gaussian]\protect\hypertarget{thm-gaussian-tail}{}\label{thm-gaussian-tail}

For \(X \sim N(0, \sigma^2)\),
\(\mathbb{E}(e^{\lambda X}) = e^{\frac{1}{2}\sigma^2 \lambda^2}\), hence
\[
P(X \ge t) \le \inf_{\lambda > 0} e^{-\lambda t}\, e^{\frac{1}{2}\sigma^2 \lambda^2} = e^{-\frac{t^2}{2\sigma^2}}.
\]

\end{theorem}

\section{Hoeffding's Inequality}\label{hoeffdings-inequality}

\begin{theorem}[Hoeffding's
theorem]\protect\hypertarget{thm-hoeffding}{}\label{thm-hoeffding}

Let \(X_i\) be independent random variables with \(X_i \in [a_i, b_i]\)
almost surely, and let \(Y = \sum_i X_i\). Then for all \(t \ge 0\), \[
P\big(|Y - \mathbb{E}(Y)| \ge t\big) \le 2\, e^{-\frac{2 t^2}{\sum_i (b_i - a_i)^2}}.
\]

\end{theorem}

\begin{lemma}[Hoeffding's
lemma]\protect\hypertarget{lem-hoeffding}{}\label{lem-hoeffding}

If \(X\) is a bounded random variable with \(X \in [a, b]\) almost
surely, then \[
\mathbb{E}\big(e^{\lambda(X - \mathbb{E}(X))}\big) \le e^{\frac{1}{2}\left(\frac{b-a}{2}\right)^2 \lambda^2} = e^{\frac{\lambda^2 (b-a)^2}{8}}.
\]

\end{lemma}

\begin{proof}
\textbf{Hoeffding's theorem.} Use the Markov tail bound for the
exponential: \[
P\Big(\sum_i X_i - \mathbb{E}(X_i) \ge t\Big) \le e^{-\lambda t}\, \mathbb{E}\Big(e^{\lambda\left(\sum_i X_i - \mathbb{E}(X_i)\right)}\Big)
= e^{-\lambda t} \prod_i \mathbb{E}\big(e^{\lambda(X_i - \mathbb{E}(X_i))}\big).
\] By Hoeffding's lemma (Lemma~\ref{lem-hoeffding}),
\(\mathbb{E}\big(e^{\lambda(X_i - \mathbb{E}(X_i))}\big) \le e^{\frac{1}{2}\frac{(b_i - a_i)^2}{4}\lambda^2}\),
and minimizing the resulting exponent over \(\lambda\) gives \[
P\Big(\sum_i X_i - \mathbb{E}(X_i) \ge t\Big) \le e^{-\frac{2 t^2}{\sum_i (b_i - a_i)^2}}.
\] The same argument applied to \(-X_i\) shows
\(P\big(\sum_i X_i - \mathbb{E}(X_i) \le -t\big)\) obeys the same bound,
so \[
P\Big(\big|\textstyle\sum_i X_i - \mathbb{E}(X_i)\big| \ge t\Big) \le 2\, e^{-\frac{2 t^2}{\sum_i (b_i - a_i)^2}}.
\]
\end{proof}

\section{\texorpdfstring{\(\chi^2\) Tail
Bound}{\textbackslash chi\^{}2 Tail Bound}}\label{chi2-tail-bound}

\begin{theorem}[\(\chi^2\) tail
bound]\protect\hypertarget{thm-chi2-tail}{}\label{thm-chi2-tail}

Let \(Z_i \overset{\text{iid}}{\sim} N(0,1)\) and
\(X = \sum_i Z_i^2 \sim \chi^2(n)\), so \(\mathbb{E}(X) = n\). Then for
all \(t \ge 0\), \[
P\big(|X - \mathbb{E}(X)| \ge t\big) \le 2\exp\!\Big[-c \min\big(t, \tfrac{t^2}{n}\big)\Big].
\]

\end{theorem}

\begin{proof}
Upper tail: with \(\lambda \in (0, \tfrac12)\), \[
P\big(X - \mathbb{E}(X) \ge t\big) \le e^{-\lambda t}\, e^{-n\lambda}\, \mathbb{E}\big(e^{\lambda \sum_i Z_i^2}\big)
= e^{-\lambda t}\, e^{-n\lambda} \prod_i \mathbb{E}\big(e^{\lambda Z_i^2}\big),
\] where \(\mathbb{E}(e^{\lambda Z_i^2})\) is the MGF of \(\chi^2(1)\),
equal to \((1 - 2\lambda)^{-1/2}\). Thus \[
P\big(X - \mathbb{E}(X) \ge t\big) \le e^{-\lambda t}\big(1 - 2\lambda\big)^{-n/2} e^{-n\lambda}
= e^{-\lambda t}\Big(\underbrace{e^{-\lambda}(1 - 2\lambda)^{-1/2}}\Big)^{n}
\le e^{-\lambda t + 2 n \lambda^2}, \qquad \forall\, \lambda \in (0, \tfrac14),
\] so that \[
\min_{\lambda \in (0, 1/4)} e^{-\lambda t + 2 n \lambda^2} = \exp\!\Big(-c \min\big(t, \tfrac{t^2}{n}\big)\Big)
\;\Longrightarrow\;
P\big(X - \mathbb{E}(X) \ge t\big) \le \exp\!\Big(-c \min\big(t, \tfrac{t^2}{n}\big)\Big).
\] Lower tail: similarly, \[
P\big(X - n \le -t\big) = P\big(-(X - n) \ge t\big) \le e^{-\lambda t}\, e^{n\lambda}\, \mathbb{E}\big(e^{-\lambda X}\big)
= e^{-\lambda t}\, e^{n\lambda}(1 + 2\lambda)^{-n/2}
\le e^{-\lambda t}\, e^{n \lambda^2}, \quad \forall\, \lambda > 0,
\] which gives \(P(X - n \le -t) \le e^{-\frac{t^2}{4n}}\). Combining
the two tails, \[
P\big(|X - \mathbb{E}(X)| \ge t\big) \le 2\exp\!\Big(-c \min\big(t, \tfrac{t^2}{n}\big)\Big).
\]
\end{proof}

\chapter{Properties of High Dimensional
Data}\label{properties-of-high-dimensional-data}

\section{Curse of dimensionality}\label{curse-of-dimensionality}

To approximate function using interpolations, need a grid in
\(\mathbb{R}^n\), which grows exponentially with respect to \(n\).

\begin{tcolorbox}[enhanced jigsaw, bottomrule=.15mm, title=\textcolor{quarto-callout-tip-color}{\faLightbulb}\hspace{0.5em}{Sampling in high dimension}, left=2mm, colback=white, toprule=.15mm, bottomtitle=1mm, titlerule=0mm, leftrule=.75mm, colbacktitle=quarto-callout-tip-color!10!white, opacitybacktitle=0.6, coltitle=black, colframe=quarto-callout-tip-color-frame, rightrule=.15mm, breakable, toptitle=1mm, arc=.35mm, opacityback=0]

To sample grid with resolution \(\epsilon\), requires points of order
\textasciitilde{}\((1/\epsilon)^n\)

\end{tcolorbox}

\section{Random Sampling in unit
ball}\label{random-sampling-in-unit-ball}

We star with a useful problem, how to draw a random sample uniformlly
from a unit ball in \(n\) dimension.

\begin{tcolorbox}[enhanced jigsaw, bottomrule=.15mm, title=\textcolor{quarto-callout-tip-color}{\faLightbulb}\hspace{0.5em}{Close Ball in \(\mathbb{R}^n\)}, left=2mm, colback=white, toprule=.15mm, bottomtitle=1mm, titlerule=0mm, leftrule=.75mm, colbacktitle=quarto-callout-tip-color!10!white, opacitybacktitle=0.6, coltitle=black, colframe=quarto-callout-tip-color-frame, rightrule=.15mm, breakable, toptitle=1mm, arc=.35mm, opacityback=0]

\[
\mathbb{B^n}\coloneqq\{\mathbf{x}\in\mathbb{R}^n: x_1^2+x_2^2+\cdots+x_n^2\leq 1\}
\]

\end{tcolorbox}

\begin{example}[]\protect\hypertarget{exm-unitSphere}{}\label{exm-unitSphere}

Consider a 2D case, we can draw \(X_1, X_2\~ \mathrm{Unif}(-1, 1)\), and
return \(\mathbf{X}=(X_1, X_2)^T\) if \(X_1^2+X_2^2\leq 1\).

\pandocbounded{\includegraphics[keepaspectratio]{ml-ds/highDimData_files/figure-pdf/cell-2-output-1.pdf}}

\end{example}

If we generalized the approch in Example~\ref{exm-unitSphere} to \(n\)
dimension case, that means draw each coordinate
\(X_i\sim \mathrm{Unif}(-1, 1)\) and keep the sample if
\(\lVert\mathbf{X}\rVert_2^2\leq 1\). However, this will not work as the
volume of the unit ball approaches to zero.

\begin{tcolorbox}[enhanced jigsaw, bottomrule=.15mm, title=\textcolor{quarto-callout-tip-color}{\faLightbulb}\hspace{0.5em}{Volume of \(\mathbb{B}^n\)}, left=2mm, colback=white, toprule=.15mm, bottomtitle=1mm, titlerule=0mm, leftrule=.75mm, colbacktitle=quarto-callout-tip-color!10!white, opacitybacktitle=0.6, coltitle=black, colframe=quarto-callout-tip-color-frame, rightrule=.15mm, breakable, toptitle=1mm, arc=.35mm, opacityback=0]

\[
V(n)=\frac{2\pi^{n/2}}{n\Gamma(n)}\xrightarrow{n\rightarrow\infty}0
\]

\end{tcolorbox}

This means the acceptance probability of the sample is
\(V(n)/\mathrm{Vol}[-1, 1]^d=V(n)/2^n\) approaches zero when the
dimensality is high.

\begin{tcolorbox}[enhanced jigsaw, bottomrule=.15mm, title=\textcolor{quarto-callout-note-color}{\faInfo}\hspace{0.5em}{Note}, left=2mm, colback=white, toprule=.15mm, bottomtitle=1mm, titlerule=0mm, leftrule=.75mm, colbacktitle=quarto-callout-note-color!10!white, opacitybacktitle=0.6, coltitle=black, colframe=quarto-callout-note-color-frame, rightrule=.15mm, breakable, toptitle=1mm, arc=.35mm, opacityback=0]

This low probability means that we only accept very few point from a
large draw of Monte-Carlo samples from uniform \([-1, 1]^n\)

\end{tcolorbox}

\chapter{Randomized Projection}\label{randomized-projection}

\textbf{Idea:} Compress high dimension data to lower dimension to reduce
computation / storage.

\textbf{Approach:} map the original data \(X \in \mathbb{R}^d\) into
\(\mathbb{R}^r\), with \(r < d\), while maintaining their pairwise
distance unchanged (or at least approximately).

\section{\texorpdfstring{Projection map
\(\phi : \mathbb{R}^d \to \mathbb{R}^r\)}{Projection map \textbackslash phi : \textbackslash mathbb\{R\}\^{}d \textbackslash to \textbackslash mathbb\{R\}\^{}r}}\label{projection-map-phi-mathbbrd-to-mathbbrr}

Let the original data lie in \(\mathbb{R}^d\), and \(S_r\) an
\(r\)-dimensional subspace of \(\mathbb{R}^d\) with orthonormal basis
\((v_1, v_2, \dots, v_r)\) spanning \(S_r\).

\begin{definition}[Linear projection
map]\protect\hypertarget{def-projection}{}\label{def-projection}

A linear projection map \(\phi : \mathbb{R}^d \to \mathbb{R}^r\) is
defined as

\[
Y = \phi(x) = \big(\langle x, v_1\rangle,\ \langle x, v_2\rangle,\ \dots,\ \langle x, v_r\rangle\big)^{\mathsf{T}}
= \begin{pmatrix} -\,v_1^{\mathsf{T}}\,- \\ -\,v_2^{\mathsf{T}}\,- \\ \vdots \\ -\,v_r^{\mathsf{T}}\,- \end{pmatrix} x
= \Phi\, x,
\quad \Phi \in \mathbb{R}^{r \times d},\ x \in \mathbb{R}^d,\ y \in \mathbb{R}^r .
\]

\end{definition}

A \textbf{random projection map} \(\phi\) is when \(S_r\) is a random
\(r\)-dimensional subspace.

\section{Approx. distance
preservation}\label{approx.-distance-preservation}

\(y_i = \Phi x_i\), we wish for small \(\varepsilon > 0\)

\[
(1-\varepsilon)\,\|x_i - x_j\|_2^2 \ \le\ \|y_i - y_j\|_2^2 \ \le\ (1+\varepsilon)\,\|x_i - x_j\|_2^2,
\qquad \forall\, i, j \in \{1, \dots, n\}.
\]

\section{Gaussian Projectors}\label{gaussian-projectors}

Consider random projection
\(\phi = \begin{pmatrix} -\,u_1^{\mathsf{T}}\,- \\ -\,u_2^{\mathsf{T}}\,- \\ \vdots \\ -\,u_r^{\mathsf{T}}\,- \end{pmatrix}\)
where \(u_i \in \mathbb{R}^d\) and
\(u_i \overset{\text{iid}}{\sim} N(0, I_d)\) spherical Gaussian vectors.

\begin{theorem}[Gaussian
projection]\protect\hypertarget{thm-gaussian-proj}{}\label{thm-gaussian-proj}

For \(\forall v \in \mathbb{R}^d\), \(\exists\, c > 0\) s.t.
\(\forall \varepsilon \in (0,1)\)

\[
P\Big(\big|\,\|\phi v\|_2 - \sqrt{r}\,\|v\|_2\,\big| \ge \varepsilon \sqrt{r}\,\|v\|_2 \Big) \le 3 e^{-c r \varepsilon^2}.
\]

\end{theorem}

\begin{remark}
The projected vector \(\phi v\) will be close to \(\sqrt{r}\,\|v\|\)
length, with high probability.
\end{remark}

\begin{proof}
\[
\phi_{ij} \sim N(0,1) \;\Rightarrow\; \phi_{ij} v_j \sim N(0, v_j^2)
\;\Rightarrow\; \sum_j \phi_{ij} v_j \sim N(0, \|v\|^2)
\]

thus

\[
\frac{1}{\|v\|}\sum_j \phi_{ij} v_j \sim N(0,1)
\;\Rightarrow\; \frac{1}{\|v\|}(\phi v)_i \sim N(0,1)
\;\Rightarrow\; \frac{1}{\|v\|}\phi v \sim N(0, I_r)
\]

using the \textbf{Gaussian Annulus Theorem}.

\[
\begin{align}
P\Big(\big|\,\|\frac{1}{\|v\|_2}\phi v\|_2 - \sqrt{r}\,\,\big| \ge \varepsilon\sqrt{r}\, \Big) &\le 3 e^{-c r \varepsilon^2} \\
\Rightarrow P\Big(\big|\,\|\phi v\|_2 - \sqrt{r}\,\|v\|_2\,\big| \ge \varepsilon\sqrt{r}\,\|v\|_2 \Big) &\le 3 e^{-c r \varepsilon^2}.
\end{align}
\]
\end{proof}

\section{Johnson--Lindenstrauss lemma}\label{johnsonlindenstrauss-lemma}

\begin{lemma}[Johnson--Lindenstrauss]\protect\hypertarget{lem-jl}{}\label{lem-jl}

Let \(x_1, \dots, x_n \in \mathbb{R}^d\), and Gaussian random projection
\(\phi = \begin{pmatrix} -\,u_1^{\mathsf{T}}\,- \\ -\,u_2^{\mathsf{T}}\,- \\ -\,u_r^{\mathsf{T}}\,- \end{pmatrix}\).
then \(\exists\, c > 0\) s.t. \(\forall \varepsilon \in (0,1)\) and
\(\delta > 0\), if

\[
r \ge \frac{c}{\varepsilon^2}\,\log\!\left(\frac{3 n^2}{\delta}\right),
\]

then

\[
P\Big(\big|\,\|\phi (x_j -  x_k)\| - \sqrt{r}\,\|x_j - x_k\|\,\big| \le \varepsilon \sqrt{r}\,\|x_j - x_k\| \Big) \ge 1 - \delta,
\qquad \forall\, j, k .
\]

\end{lemma}

\begin{remark}
The JL lemma states that the distance between projected points

\[
(1-\varepsilon)\sqrt{r}\,\|x_i - x_j\| \ \le\ \|\phi(x_i - x_j)\| \ \le\ (1+\varepsilon)\sqrt{r}\,\|x_i - x_j\|
\]

with high probability.
\end{remark}

\begin{proof}
Let \(v_{ij} = x_i - x_j\), from the Gaussian projection theorem
(Theorem~\ref{thm-gaussian-proj}), we know
\(\exists\, c_{ij},\ \varepsilon_{ij} \in (0,1)\) s.t.

\[
P\Big(\big|\,\|\phi v_{ij}\| - \sqrt{r}\,\|v_{ij}\|\,\big| \ge \varepsilon_{ij}\sqrt{r}\,\|v_{ij}\| \Big) \le 5 e^{-c_{ij} r \varepsilon_{ij}^2}.
\]

to show JL lemma, we need a union bound of event

\[
E_{ij}:\quad \big|\,\|\phi v_{ij}\| - \sqrt{r}\,\|v_{ij}\|\,\big| \ge \varepsilon_{ij}\sqrt{r}\,\|v_{ij}\|
\]

so that above is true for \(\forall i, j\).

\[
P\Big(\bigcup_{ij} E_{ij}\Big) \le \sum_{ij} P(E_{ij}) \le \sum_{ij} 3 e^{-c_{ij} r \varepsilon_{ij}^2} \le \binom{n}{2} e^{-c r \varepsilon^2} 3
\]

where \(c = \min\{c_{ij}\}\),
\(\varepsilon = \min\{\varepsilon_{ij}\}\), \(\sum_{ij}\) has
\(\binom{n}{2}\) terms. Therefore

\[
P\Big(\bigcup_{ij} E_{ij}\Big) \le 3\binom{n}{2} e^{-c r \varepsilon^2} \le 3\tfrac{1}{2} n^2 e^{-c r \varepsilon^2} \le \delta
\]

\[
\Rightarrow\ \delta \ge 3 n^2 e^{-c r \varepsilon^2}
\ \Rightarrow\ r \ge \frac{1}{c \varepsilon^2}\,\log\!\left(\frac{3n^2}{\delta}\right).
\]
\end{proof}

\section{Computing Random Projection}\label{computing-random-projection}

\textbf{Direct computing} requires \(O(nrd)\) time, \(n\) is sample
size.

\textbf{Idea:} if \(X = (x_1, \dots, x_d)\) has iid entry, we can
randomly select \(r\) entry and approximate
\(\|x\|^2_2 \approx \tfrac{d}{r}\sum_{k}^rv_{ik}\), where \(i_k\) is
randomly selected in \(\{1,2,\cdots,d\}\). Thus
\(\phi = \sqrt{\tfrac{d}{r}}\, P\), where \(P\) is a random selection
matrix should approximatly satisfy \(\|\phi x\|_2 \approx \|x\|_2\).

\textbf{Issue:} \(X\) is usually not iid entry, but tends to be sparse.

\textbf{Fix:} consider if \(x\) is sparse (\(\delta\)-function peaks),
its Fourier transform would be smooth (sin function). We can find a
similar transform matrix s.t. the transformed \(X\) is flat.

\section{Hadamard transform}\label{hadamard-transform}

\begin{definition}[Hadamard
transform]\protect\hypertarget{def-hadamard}{}\label{def-hadamard}

Hadamard transform matrix \(H_m \in 2^m \times 2^m\) is defined
recursively as

\[
\begin{align}
H_1 &= \frac{1}{\sqrt{2}}\begin{pmatrix} 1 & 1 \\ 1 & -1 \end{pmatrix}, \qquad
H_2 = \frac{1}{\sqrt{2}}\begin{pmatrix} H_1 & H_1 \\ H_1 & -H_1 \end{pmatrix},
\quad \dots \quad\\
H_m &= \frac{1}{\sqrt{2}}\begin{pmatrix} H_{m-1} & H_{m-1} \\ H_{m-1} & -H_{m-1} \end{pmatrix},
\end{align}
\]

each entry \(\dfrac{1}{\sqrt{2^m}}\).

\end{definition}

\begin{remark}
\leavevmode

\begin{itemize}
\tightlist
\item
  \(H_m\) is orthogonal,
  \(H_m H_m^{\mathsf{T}} = H_m^{\mathsf{T}} H_m = I\).
\item
  \(H_m x\) can be computed using a similar algorithm as FFT, with
  \(O(2^m \log(2^m))\) time.
\item
  \(H_m\) maps a sparse vector to a flat vector.
\item
  flatness can be further improved by randomization.
\end{itemize}

\end{remark}

\[
D_{ij} = \begin{cases} \delta_{ij}, & p = \tfrac{1}{2} \\[4pt] \delta_{ij} & p = \tfrac{1}{2} \end{cases}
\]

\(H_m D\) is the \textbf{randomized Hadamard transform} (independent
random sign flip).

\begin{lemma}[Flatness of
\(H_m D x\)]\protect\hypertarget{lem-hadamard-flat}{}\label{lem-hadamard-flat}

\[
\forall x \in \mathbb{R}^d, \qquad P\big(\|H D x\|_\infty \ge t\big) \le 2d\, e^{-\frac{d\, t^2}{2\|x\|^2}}.
\]

\end{lemma}

\begin{proof}
\(u_i := (H D x)_i=\sum_{j=1}^da_jx_j\), with

\[
a_{ij} = \begin{cases} \tfrac{1}{\sqrt{d}}, & p = \tfrac{1}{2} \\ -\tfrac{1}{\sqrt{d}}, & p = \tfrac{1}{2} \end{cases}
\]

for each entry \(i\):

\[
P(|u_i| \ge t) \le 2 \exp\!\left(-\frac{t^2}{2 \cdot \frac{1}{d}\|x\|^2 \cdot \tfrac{1}{?}}\right)
\quad \text{(Hoeffding)}
= 2 \exp\!\left(-\frac{d\, t^2}{2\|x\|^2}\right)
\]

\[
P(\|u\|_\infty \ge t) \le d\, P(|u_i| \ge t) \le 2d \exp\!\left(-\frac{d\, t^2}{2\|x\|^2}\right).
\]
\end{proof}

\begin{remark}
\(\frac{1}{\sqrt{d}}\|H D x\|_2 = \frac{1}{\sqrt{d}}\|x\|_2\); if
\(t = k \frac{\|x\|}{\sqrt{d}}\) then the probability of \(\|H D x\|_2\)
not flat (one component is very large comparing to other
\(\leftrightarrow \|x\|/\sqrt{d}\) with large \(k\)) would be small.
\end{remark}

\section{Fast JLT}\label{fast-jlt}

Since \(H D x\) is flat with high probability, then
\(\phi = \sqrt{\tfrac{d}{r}}\, P H D\) would preserve the norm with high
prob

\[
\|\phi x\| = \sqrt{\tfrac{d}{r}}\,\|P H D x\| \approx \|H D x\| \approx \|x\|_2 .
\]

\begin{remark}
The computation cost for Fast JLT is

\[
O\big(n(d \log d + d)\big) \ \longrightarrow\ O(n d \log d).
\]

\begin{itemize}
\tightlist
\item
  \(Hx\) takes \(O(d\log d)\)
\item
  \(P\) and \(D\) tales \(O(d)\)
\end{itemize}

\end{remark}

\chapter{Singular Value
Decomposition}\label{singular-value-decomposition}

\section{Singular Vectors}\label{singular-vectors}

\begin{definition}[First right singular
vector]\protect\hypertarget{def-1stSingularVec}{}\label{def-1stSingularVec}

The first right singular vector \(v_1\) of
\(X \in \mathbb{R}^{n \times d}\) is defined as:
\(v_1 \in \mathbb{R}^d\), and
\(v_1 = \text{argmax}_{\|v\|_2=1} \|Xv\|_2\) and the first singular
value is \(\sigma_1 = \|Xv_1\|_2\)

\end{definition}

\begin{proposition}[]\protect\hypertarget{prp-equivSingular}{}\label{prp-equivSingular}

Definition~\ref{def-1stSingularVec} is equivalent as \[
v_1 = \underset{\|v\|_2=1}{\mathrm{argmin}} \sum_{j=1}^{n} \|x_j - P_{v_1}(x_j)\|_2^2,
\] where \(P_{v_1}(x_j)\) is the projected vector of \(x_j\) in the
subspace spanned by \(v_1\). Here, we have \[
X \in \mathbb{R}^{n \times d} = 
\begin{pmatrix} -x_1^T- \\ -x_2^T- \\ \vdots \\ -x_n^T- \end{pmatrix}
\]

\end{proposition}

\begin{proof}
\[
\begin{align}
\sum_{j} \|x_j - P_{v}(x_j)\|_2^2 &= \sum_{ij}\left(x_{i,j}-\langle\mathbf{v}, \mathbf{x}_j\rangle v_i\right)^2 \\
&=\sum_{ij}x^2_{i, j} + \langle\mathbf{v}, \mathbf{x}_j\rangle v_i^2 - 2x_{i,j}v_i\langle\mathbf{v}, \mathbf{x}_j \rangle\\
&=\sum_{ij}x^2_{i,j} + \sum_j\left(\langle\mathbf{v},\mathbf{x}_j\rangle ^2 -2\langle\mathbf{v},\mathbf{x}_j\rangle^2\right)\\
&=C - \sum_j\langle\mathbf{v},\mathbf{x}_j\rangle^2=C-\|X\mathbf{v}\|^2_2
\end{align}
\] where \(C\) is a constant. Thus we have the \[
\underset{\|v\|_2=1}{\mathrm{argmin}} \sum_{j=1}^{n} \|x_j - P_{v_1}(x_j)\|_2^2 = \underset{\|v\|_2=1}{\mathrm{argmax}} \|Xv\|_2^2
\]
\end{proof}

\begin{tcolorbox}[enhanced jigsaw, bottomrule=.15mm, title=\textcolor{quarto-callout-note-color}{\faInfo}\hspace{0.5em}{Note}, left=2mm, colback=white, toprule=.15mm, bottomtitle=1mm, titlerule=0mm, leftrule=.75mm, colbacktitle=quarto-callout-note-color!10!white, opacitybacktitle=0.6, coltitle=black, colframe=quarto-callout-note-color-frame, rightrule=.15mm, breakable, toptitle=1mm, arc=.35mm, opacityback=0]

Proposition~\ref{prp-equivSingular} shows that the first singular vector
can be understood as a line that minimizes least square error of
projected data, i.e., a best fitting line.

\end{tcolorbox}

Next, we continue to define the second singular vector.

\begin{definition}[Second Singular
Vector]\protect\hypertarget{def-2ndSingularVec}{}\label{def-2ndSingularVec}

The second right singular vector \(v_2 \in \mathbb{R}^d\) is defined as
\[
v_2 = \underset{{v \perp v_1, \|v\|_2=1}}{\mathrm{argmin}} \|Xv\|_2,
\] and second singular value defined correspondingly as \[
\sigma_2 = \|Xv_2\|_2
\]

\end{definition}

\begin{proposition}[Best-Fitting 2D
Subspace]\protect\hypertarget{prp-best2dSubspace}{}\label{prp-best2dSubspace}

Let \(v_1, v_2\) be the first and second right singular vectors of
\(X \in \mathbb{R}^{n \times d}\). Let
\(V = \text{span}(v_1, v_2) \subseteq \mathbb{R}^d\) be the
two-dimensional subspace they span.

Then \(V\) is the solution to the following minimization: \[
V = \underset{W:\,\dim(W)=2}{\mathrm{argmin}} \sum_{j=1}^{n} \|x_j - P_W(x_j)\|_2^2
\] where \(P_W:\mathbb{R}^d\rightarrow\mathbb{R}^2\) is the orthogonal
projection onto the subspace \(W\).

\end{proposition}

\begin{proof}
Let \(w_1, w_2\) be two orthonormal bases for any 2-dimensional subspace
\(W\). We can always select \(w_2\) s.t. \(\langle w_2,v_1 \rangle = 0\)
(such \(w_2\) can be found by the intersection between \(W\) and
subspace that perpendicular to \(v_1\)).

\textbf{A figure perhaps showing the constructiong of w1 and w2} \[
\begin{align}
&\sum_j \|x_j - \langle P_W(x_j)\rangle\|_2^2 = \sum_j \langle x_j-P_W(x_j), x_j-P_W(x_j) \rangle\\ 
={}&\sum_j\langle x_j - P_w(x_j), x_j\rangle + \sum_j\langle x_j - P_W(x_j), -P_W(x_j)\rangle\\
={}&\sum_j\|x_j\|^2_2 - 2\sum_j\langle P_W(x_j), x_j\rangle + \|P_W(x_j)\|^2_2\\
\end{align}.
\] Noticed that \[
\begin{align}
\langle P_W(x_j), x_j \rangle &= \langle w_1\langle w_1, x_j\rangle + w_2\langle w_2, x_j\rangle, x_j\rangle\\
& = \langle w_1, x_j\rangle^2 + \langle w_2, x_j\rangle^2=\|P_W(x_j)\|^2_2,
\end{align}
\] thus \[
\sum_j\|x_j - P_W(x_j)\|^2_2 = C - (\sum_j \langle w_1, x_j \rangle^2 + \langle w_2, x_j \rangle^2),
\] where \[
C = \sum_{j=1}^n \| x_j \|_2^2
\] is a constant.

Since \(v_1\) maximizes \(\|Xv\|_2^2\) among all \(\|v\|_2=1\) , and
\(v_2\) maximizes \(\|Xv\|_2^2\) among \(v \perp v_1\) and
\(\|v\|_2=1, thus\)\$
\sum\_j\langle w\_1,x\_j\rangle\leq sum\_j\langle v\_1, x\_j\rangle,
\sum\_j\langle w\_2, x\_j\rangle \leq \langle v\_2, x\_j\rangle \$\$ and
thus \(V\) minimize \(\|x_j-P_W(x_j)\|^2_2\) among all 2-dim subspace.
\end{proof}

\begin{tcolorbox}[enhanced jigsaw, bottomrule=.15mm, title=\textcolor{quarto-callout-note-color}{\faInfo}\hspace{0.5em}{Note}, left=2mm, colback=white, toprule=.15mm, bottomtitle=1mm, titlerule=0mm, leftrule=.75mm, colbacktitle=quarto-callout-note-color!10!white, opacitybacktitle=0.6, coltitle=black, colframe=quarto-callout-note-color-frame, rightrule=.15mm, breakable, toptitle=1mm, arc=.35mm, opacityback=0]

Proposition~\ref{prp-best2dSubspace} meaning that the \(v_1\) and
\(v_2\) we constructed form the best fitting subspace (2-dimensional)
that minimizes projected least square error.

\end{tcolorbox}

\textbf{Posible code example for demonstrating the best fitting sub
space}

\begin{definition}[r-right Singular
Vector]\protect\hypertarget{def-rSingularVec}{}\label{def-rSingularVec}

The leading \(r\) right singular vectors
\(v_1, \dots, v_r \in \mathbb{R}^d\) are defined as:

\[
v_k := \underset{
  \substack{
  \|v\|_2=1,\, \\ v \perp \text{span}\{v_1, \dots, v_{k-1}\}
  }
  }{\mathrm{argmax}}{} \|Xv\|_2
\]

with the corresponding singular values as \(\sigma_k = \|Xv_k\|_2\).

\end{definition}

\begin{proposition}[Greedy construction for best-fit
subspace]\protect\hypertarget{prp-bestRSubspace}{}\label{prp-bestRSubspace}

Let \(v_1, \dots, v_r\) be the leading right singular vectors
constructed by Definition~\ref{def-rSingularVec}, and
\(V = \mathrm{span}(v_1, \dots, v_r) \subset \mathbb{R}^d\) be the
subspace spanned by them. Then \[
V \in \underset{
  \substack{W\subseteq \mathbb{R}^d \\ \mathrm{dim}(W) = r \\ \text{linear subspace}}
}{\mathrm{argmin}}\sum_j \|x_j - P_W(x_j)\|_2^2
\] i.e., \(V\) is the best-fit linear subspace in terms of least square
error.

\end{proposition}

\begin{refremark}
\leavevmode

\begin{itemize}
\tightlist
\item
  Proposition~\ref{prp-bestRSubspace} can be proved similarly as in the
  2 dimensional case (Proposition~\ref{prp-best2dSubspace}).
\item
  Proposition~\ref{prp-bestRSubspace} shows that the greedy construction
  of r-right singularvectors formed a best-fit subspace in terms of
  least square error.
\end{itemize}

\label{rem-bestRsubspace}

\end{refremark}

\section{Singular Value
Decomposition}\label{singular-value-decomposition-1}

\begin{theorem}[Singular Value Decomposition
(SVD)]\protect\hypertarget{thm-SVD}{}\label{thm-SVD}

For any matrix \(X \in \mathbb{R}^{n \times d}\), there exist orthogonal
matrices \(U \in \mathbb{R}^{n \times n}\) and
\(V \in \mathbb{R}^{d \times d}\) and a non-negative diagonal matrix
\(\Sigma \in \mathbb{R}^{n \times d}\) such that

\[
X_{n \times d} = U_{n \times n}\, \Sigma_{n \times d}\, V_{d \times d}^T
\]

\end{theorem}

\begin{refremark}
\leavevmode

\begin{itemize}
\tightlist
\item
  The diagonals of \(\Sigma\), arranged in non-increasing order, are the
  singular values defined in Proposition~\ref{prp-bestRSubspace}.
\item
  Columns of \(U\) are the left singular vectors of \(X\).
\item
  The left singular vectors \(u_i := \frac{1}{\sigma_i} X v_i\), where
  \(\sigma_i\) is the \(i\)-th diagonal element of \(\Sigma\).
\item
  Columns of \(V\) are the right singular vectors of \(X\) defined in
  Proposition~\ref{prp-bestRSubspace}.
\item
  SVD can be viewed as `generalized diagonolization' of an arbitrary
  matrix.
\end{itemize}

\label{rem-SVD}

\end{refremark}

The existence of SVD can be shown as follow

\begin{refproof}
Let \(X \in \mathbb{R}^{n \times d}\) and
\(C = X^T X \in \mathbb{R}^{d \times d}\), then \(C\) is symmetric and
positive semidefinite. Let the orthogonal decomposition of \(C\) be: \[
C = V \Lambda V^T,
\] where \(V^T V = I\) and
\(\Lambda = \mathrm{diag}(\lambda_1, \lambda_2, \dots, \lambda_d), \lambda_i \geq 0 ,\forall i\in\{1, 2, \cdots, d\}\).

Consider \(\Sigma \in \mathbb{R}^{n \times d}\) such that
\(\Sigma^T \Sigma = \Lambda\). If \(\mathrm{rank}(X) = r \leq d\), we
can arrange the diagonal element of \(\Lambda\) s.t.
\(\lambda_1\geq\lambda_2\cdots\lambda_r\geq0,\lambda_{r+1} = \cdots=\lambda_d=0\).
Then \[
\Sigma = \begin{pmatrix}
\mathrm{diag}(\sqrt{\lambda_1},\cdots,\sqrt{\lambda_r}) & 0\\
0 & 0
\end{pmatrix}
\]

We define the left singular vector \(u_i\) as: \[
u_i := \frac{1}{\sigma_i} X v_i,\quad i = 1, \dots, r.
\] Then \(u_i\) are orthonormal: \[
u_i^T u_j = \frac{1}{\sigma_i \sigma_j} v_i^T X^T X v_j = \frac{1}{\sigma_i \sigma_j} v_i^T C v_j = \frac{1}{\sigma_i \sigma_j} \lambda_j \delta_{ij} = \delta_{ij}.
\] We can then extend \(u_{r+1}, \dots, u_n\) such that
\(U = [u_1\,|\,u_2\,|\,\dots\,|\,u_n] \in \mathbb{R}^{n\times n}\) is
orthogonal.\\
Thus, we have \(\sigma_i u_i = X v_i\) for \(i=1,\ldots, r\) and
\(X v_i = 0\) for \(i > r\), which gives \(X = U \Sigma V^T\).

\label{prf-SVD}

\end{refproof}

\begin{refremark}
This proof can be extended to \(X \in \mathbb{C}^{n \times d}\) with
unitary matrices \(U\) and \(V\).

\label{rem-svd-complex}

\end{refremark}

\section{Greedy Construction and SVD}\label{greedy-construction-and-svd}

\begin{proposition}[]\protect\hypertarget{prp-svd-greedy}{}\label{prp-svd-greedy}

The SVD of \(X = U \Sigma V^T\) is the same as the greedy construction;
i.e., \(V\) and \(\Sigma\) are right singular vectors and singular
values constructed by the greedy definitions.

\end{proposition}

\begin{refproof}
Let \(v_1' := \mathrm{argmax}_{\|v\|_2=1} \|Xv\|_2\) be the first right
singular vector.\\
Write \(X = U \Sigma V^T\) where
\(V = [v_1, \dots, v_d] \in \mathbb{R}^{d \times d}\) and
\(v_i^T v_j = \delta_{ij}\).\\
Then: \[
\|Xv\|_2 = \|U \Sigma V^T v\|_2 = \|\Sigma V^T v\|_2
\] Let \(z = V^T v\), then \(\|z\|_2 = 1\) since \(V\) is orthogonal and
\(\|v\|_2 = 1\).\\
So \[
\|Xv\|_2 = \|\Sigma z\|_2 = \sqrt{\sum_{i} (\sigma_i z_i)^2}
\] Suppose \(\sigma_1 \geq \sigma_2 \geq \dots \geq \sigma_r\). Notice
\[
\sum \sigma_i^2 z_i^2 \leq \sigma_1^2 \left(\sum z_i^2\right) = \sigma_1^2
\] Equality is attained when \(z_1 = 1\)
\(\Rightarrow v = v_1 \Rightarrow v_1' = v_1\).

Next, \(v_2' := \mathrm{argmax}_{\|v\|_2=1,\, v \perp v_1} \|Xv\|_2\).
We proceed similarly, but now \(z_1 = 0\) (since \(v \perp v_1\)).\\
Then \[
\|Xv\|_2 = \sqrt{\sum_{i=2} (\sigma_i z_i)^2} \leq \sigma_2
\] Equality when \(z_2 = 1 \Rightarrow v_2' = v_2\).\\
Similarly, we find all \(v_i\) by the greedy construction.

\label{prf-svd-greedy}

\end{refproof}

\section{Skinny SVD and Full SVD}\label{skinny-svd-and-full-svd}

\textbf{Full SVD:}\\
\[
X_{n \times d} = U_{n \times n} \Sigma_{n \times d} V_{d \times d}^T
\] where \(\mathrm{rank}(X) = r \leq \min(n, d)\).\\
Notice \((\sigma_{r+1}, \dots, \sigma_{\min(d,n)}) = 0\).

\textbf{Truncated SVD (Skinny SVD):}\\
Let - \(U_r := [u_1, \dots, u_r] \in \mathbb{R}^{n \times r}\) -
\(V_r := [v_1, \dots, v_r] \in \mathbb{R}^{d \times r}\) -
\(\Sigma_r = \mathrm{diag}(\sigma_1, \dots, \sigma_r)\)

then \[
X_{n \times d} = U_r \Sigma_r V_r^T = \sum_{i=1}^{r} \sigma_i u_i v_i^T
\]

\section{Computing SVD}\label{computing-svd}

\subsection{Power Method}\label{power-method}

\textbf{Idea:} Let \(C = X^T X \in \mathbb{R}^{d \times d}\) and
\(X = U \Sigma V^T = \sum_{i=1}^{r} \sigma_i u_i v_i^T\). Then

\[
C = \sum_{i=1}^{r} \sigma_i^2 v_i v_i^T
\quad\xrightarrow{\text{iterate}}\quad
C^2 = \sum_{i=1}^{r} \sigma_i^4 v_i v_i^T
\quad\longrightarrow\quad
C^k = \sum_{i=1}^{r} \sigma_i^{2k} v_i v_i^T .
\]

If \(\sigma_1 > \sigma_2\), then for \(k\) large enough
\(\sigma_1^{2k} \gg \sigma_2^{2k}\), so

\[
C^k \approx \sigma_1^{2k}\, v_1 v_1^T
\;\Rightarrow\;
C^k \begin{pmatrix} 1 \\ 0 \\ \vdots \\ 0 \end{pmatrix}
= C^k e_1 \approx \sigma_1^{2k}\, v_1(1)\, v_1
\;\Rightarrow\;
\frac{C^k e_1}{\|C^k e_1\|_2} = \pm\, v_1 .
\]

\textbf{Linear rate of convergence.} To check the convergence rate, we
look at

\[
\min\Big\{\, \|v^{(k)} - v_1\|_2^2,\ \|v^{(k)} + v_1\|_2^2 \,\Big\} = 2 - 2\,\big|\langle v^{(k)}, v_1\rangle\big|,
\qquad v^{(k)} := \frac{C^k e_1}{\|C^k e_1\|_2}.
\]

Since the \(v_i\) are orthonormal,

\[
1 = \|v^{(k)}\|_2^2 = \sum_{i} \big|\langle v^{(k)}, v_i\rangle\big|^2 = \big|\langle v^{(k)}, v_1\rangle\big|^2 + \sum_{i=2}^{r} \big|\langle v^{(k)}, v_i\rangle\big|^2 .
\]

For convergence we wish \(|\langle v^{(k)}, v_1\rangle| \approx 1\),
i.e.~the projection onto the other subspace is small:

\[
\frac{\sum_{i=2}^{r} \big|\langle v^{(k)}, v_i\rangle\big|^2}{\big|\langle v^{(k)}, v_1\rangle\big|^2}
= \frac{\sum_{i=2}^{r} \sigma_i^{4k}\, v_i^2(1)}{\sigma_1^{4k}\, v_1^2(1)}
\le \frac{\sum_{i=2}^{r} \sigma_i^{4k}}{\sigma_1^{4k}\, v_1^2(1)}
\le \frac{(r-1)\, \sigma_2^{4k}}{\sigma_1^{4k}\, \delta^2},
\qquad \text{if } |v_1(1)| \ge \delta > 0 ,
\]

hence

\[
\frac{1 - \big|\langle v^{(k)}, v_1\rangle\big|^2}{\big|\langle v^{(k)}, v_1\rangle\big|^2}
\le (r-1)\, \frac{\sigma_2^{4k}}{\sigma_1^{4k}\, \delta^2}
\;\longrightarrow\; \text{linear convergence, as } \sigma_2 < \sigma_1 .
\]

\begin{remark}
Computing \(C^k\) is costly \(\to\) compute \(C^k q\) instead; each step
only needs a matrix--vector product (efficient when one dimension of
\(X\) is small, or \(X\) is sparse).

\end{remark}

The same analysis holds for an arbitrary starting vector \(q\) (not just
\(e_1\)):

\begin{proposition}[Convergence to the near-top singular
subspace]\protect\hypertarget{prp-power-general}{}\label{prp-power-general}

Let \(q \in \mathbb{R}^d\) with
\(|\langle q, v_1\rangle| \ge \delta > 0\), and let \(U\) be the linear
subspace spanned by the \(v_i\) such that
\(\sigma_i \ge (1-\varepsilon)\sigma_1\). Let
\(v^{(k)} = \dfrac{C^k q}{\|C^k q\|_2}\) be the \(k\)-th power
iteration. Then it can be shown that

\[
\big\|\operatorname{Proj}_{U^{\perp}}\big(v^{(k)}\big)\big\|_2 \le \varepsilon,
\qquad \forall\, \varepsilon > 0, \quad \text{when } k \ge \frac{\log(1/\varepsilon\delta)}{2\varepsilon}.
\]

\end{proposition}

\begin{proof}
Expand \(q = \sum_{i} c_i v_i\) in the orthonormal basis
\(\{v_1, v_2, \dots, v_d\}\) of \(\mathbb{R}^d\). Then

\[
C^k q = \sum_{i=1}^{r} \sigma_i^{2k} v_i v_i^T \sum_{j} c_j v_j = \sum_{i=1}^{r} \sigma_i^{2k}\, c_i\, v_i,
\qquad
\|C^k q\|_2^2 = \sum_{i=1}^{r} \sigma_i^{4k}\, c_i^2 .
\]

For the first (single-vector) case, with
\(\sigma_1 \ge \sigma_2 \ge \cdots\),

\[
\frac{\sum_{i=2}^{r} \big|\langle v^{(k)}, v_i\rangle\big|^2}{\big|\langle v^{(k)}, v_1\rangle\big|^2}
= \frac{\sum_{i=2}^{r} \sigma_i^{4k}\, c_i^2}{c_1^2\, \sigma_1^{4k}}
\le \frac{\sum_{i=2}^{r} \sigma_i^{4k}}{c_1^2\, \sigma_1^{4k}}
\le \Big(\tfrac{\sigma_2}{\sigma_1}\Big)^{4k}\, (r-1)\, \frac{1}{\delta^2},
\]

where we assume \(\|q\|_2 = 1\) and
\(|c_1| = |\langle v_1, q\rangle| \ge \delta > 0\).

For the subspace statement, suppose \(\sigma_1, \dots, \sigma_m\)
satisfy \(\sigma_i \ge (1-\varepsilon)\sigma_1\), so that the component
of \(v^{(k)}\) orthogonal to \(U\) is composed of indices
\(i = m+1, \dots, r\). Then

\[
\big\|\operatorname{Proj}_{U^{\perp}}\big(v^{(k)}\big)\big\|_2^2
= \frac{\sum_{i=m+1}^{r} \sigma_i^{4k}\, c_i^2}{\sum_{i=1}^{r} c_i^2\, \sigma_i^{4k}}
\le \frac{(1-\varepsilon)^{4k}\, \sigma_1^{4k}}{c_1^2\, \sigma_1^{4k}}
\le \frac{(1-\varepsilon)^{4k}\, \sigma_1^{4k}}{\delta^2\, \sigma_1^{4k}},
\]

hence

\[
\big\|\operatorname{Proj}_{U^{\perp}}\big(v^{(k)}\big)\big\|_2
\le \frac{(1-\varepsilon)^{2k}}{\delta}
\le \frac{e^{-2k\varepsilon}}{\delta}
\le \varepsilon
\qquad \text{when } k \ge \frac{\ln(1/\varepsilon\delta)}{2\varepsilon}.
\]
\end{proof}

\begin{remark}
Note the above convergence analysis assumes \(\sigma_1 > \sigma_2\)
strictly. In fact, the power method converges to some vector in the span
of the singular vectors with ``nearly highest'' singular value.

\(q\) can be chosen at random, e.g.~\(q \sim N(0, I_d)\). When chosen
randomly, \(|\langle q, v_1\rangle| \ge \delta > 0\) can be shown to
hold with high probability.
\end{remark}

\subsection{Subspace (Orthogonal Iteration) Power
Method}\label{subspace-orthogonal-iteration-power-method}

Finding the right singular vectors is the same as finding the
eigenvectors of \(X^T X\). One can use the \textbf{orthogonal iteration}
algorithm for that.

Let \(M \in \mathbb{R}^{d \times d}\) be positive semidefinite with
eigendecomposition
\(M = V \Lambda V^T = \sum_{i=1}^{d} \lambda_i v_i v_i^T\).

\begin{tcolorbox}[enhanced jigsaw, leftrule=.75mm, bottomrule=.15mm, left=2mm, colback=white, colframe=quarto-callout-note-color-frame, toprule=.15mm, breakable, arc=.35mm, rightrule=.15mm, opacityback=0]
\begin{minipage}[t]{5.5mm}
\textcolor{quarto-callout-note-color}{\faInfo}
\end{minipage}%
\begin{minipage}[t]{\textwidth - 5.5mm}

\textbf{Orthogonal iteration}

Initialize \(Q_0 \in \mathbb{R}^{d \times \ell}\), a random
\(\ell\)-dimensional subspace of \(\mathbb{R}^d\).

Repeat until convergence:

\begin{itemize}
\tightlist
\item
  \(t + 1 \leftarrow t\)
\item
  \(Z_{t+1} = M \cdot Q_t\)
\item
  \(Q_{t+1} = \mathrm{QR}[Z_{t+1}]\), where \(\mathrm{QR}[Z] := Q\) when
  \(Z = QR\) is the QR decomposition.
\end{itemize}

\(Q\) then represents the first \(\ell\) right singular vectors.

\end{minipage}%
\end{tcolorbox}

\section{\texorpdfstring{Best rank-\(k\)
Approximation}{Best rank-k Approximation}}\label{best-rank-k-approximation}

\subsection{Recap on matrix norms}\label{recap-on-matrix-norms}

\textbf{Frobenius norm} \(\|\cdot\|_F\): \[
\|X\|_F = \sqrt{\sum_{i,j} x_{ij}^2} = \sqrt{\sum_{i=1}^{r} \sigma_i^2}
\quad (\sigma_i \text{ singular values}).
\]

\textbf{Schatten \(p\)-norm} \(\|\cdot\|_{S_p}\): operates on the
singular values. For \(X \in \mathbb{R}^{n \times d}\) with
\(r \le \min(n,d)\) equal to \(\mathrm{rank}(X)\), \[
\|X\|_{S_p} := \Big(\sum_{i=1}^{r} \sigma_i^p(X)\Big)^{1/p}, \qquad 1 \le p < \infty .
\]

\begin{itemize}
\tightlist
\item
  The spectral norm \(\|\cdot\|_2 = \|\cdot\|_{S_\infty} = \sigma_1\),
  and \(\|\cdot\|_F = \|\cdot\|_{S_2}\).
\item
  \(\displaystyle \|X - X_k\|_{S_p} = \min_{B:\ \mathrm{rank}(B) \le k} \|X - B\|_{S_p} = \Big(\sum_{i=k+1}^{r} \sigma_i^p\Big)^{1/p}\).
\end{itemize}

\subsection{\texorpdfstring{Best rank-\(k\)
approximation}{Best rank-k approximation}}\label{best-rank-k-approximation-1}

\begin{definition}[\(k\)-truncated
SVD]\protect\hypertarget{def-truncated-svd}{}\label{def-truncated-svd}

Let \(X \in \mathbb{R}^{n \times d}\) with
\(\mathrm{rank}(X) = r \le \min\{n, d\}\). The \textbf{\(k\)-truncated
SVD} of \(X\) is \[
X_k = \sum_{i=1}^{k} \sigma_i u_i v_i^T, \qquad k \in \{1, 2, \dots, r\},
\] where \(X = \sum_{i=1}^{r} \sigma_i u_i v_i^T\) is the full SVD.

\end{definition}

\begin{theorem}[Best rank-\(k\) approximation
(Eckart--Young--Mirsky)]\protect\hypertarget{thm-eckart-young}{}\label{thm-eckart-young}

\(X_k\) is the ``best'' rank-\(k\) approximation of \(X\) in terms of
any orthogonally invariant norm \(\|\cdot\|\): \[
X_k = \underset{B \in \mathbb{R}^{n \times d},\ \mathrm{rank}(B) \le k}{\arg\min} \|X - B\| ,
\] i.e.~\(X_k\) is the closest rank-\(k\) matrix to \(X\).

\end{theorem}

\begin{remark}
A norm \(\|\cdot\|\) is \textbf{orthogonally invariant} if
\(\|O_1 X O_2\| = \|X\|\) for all \(X \in \mathbb{R}^{n \times d}\) and
all orthogonal \(O_1 \in \mathbb{R}^{n \times n}\),
\(O_2 \in \mathbb{R}^{d \times d}\). Examples: the Frobenius norm
\(\|\cdot\|_F\) and the spectral norm \(\|\cdot\|_2\).
\end{remark}

\begin{proof}
\textbf{Spectral norm.} (In the manuscript \(A\) is used for \(X\) in
this proof.) Let \(X \in \mathbb{R}^{n \times d}\) with
\(\mathrm{rank}(X) = r \le \min\{n,d\}\),
\(X = \sum_{i=1}^{r} \sigma_i u_i v_i^T\) and
\(X_k = \sum_{i=1}^{k} \sigma_i u_i v_i^T\). Then \[
\|X - X_k\|_2 = \Big\|\sum_{i=k+1}^{r} \sigma_i u_i v_i^T\Big\|_2 = \sigma_{k+1},
\] so we hope to show \(\|X - B\|_2 \ge \sigma_{k+1}\) for any \(B\)
with \(\mathrm{rank}(B) \le k\).

Consider \(w \in \mathbb{R}^d\) with unit norm \(\|w\|_2 = 1\). We can
choose \(w \in \mathrm{Null}(B)\), i.e.~\(Bw = 0\); on the other hand,
we can also choose \(w \in \mathrm{span}\{v_1, \dots, v_{k+1}\}\). In
this case \[
\|X w\|_2 = \sqrt{\sigma_1^2 \cdots \sigma_{k+1}^2} \ge \sigma_{k+1} .
\] Thus if
\(w \in \mathrm{Null}(B) \cap \mathrm{span}\{v_1, \dots, v_{k+1}\}\),
then \[
\|X - B\|_2 = \max_{v \in \mathbb{R}^d,\, \|v\|=1} \|(X-B)v\|_2 \ge \|(X-B)w\|_2 = \|X w\|_2 \ge \sigma_{k+1} .
\]
\end{proof}

\begin{proof}
\textbf{Frobenius norm.} Consider \(\|X - B\|_F\), where
\(\mathrm{rank}(B) = p \le k \le \mathrm{rank}(X) = r\). Denote \[
B = \begin{pmatrix} -\,b_1^T\,- \\ -\,b_2^T\,- \\ \vdots \\ -\,b_n^T\,- \end{pmatrix}, \qquad
X = \begin{pmatrix} -\,x_1^T\,- \\ -\,x_2^T\,- \\ \vdots \\ -\,x_n^T\,- \end{pmatrix}
\;\Rightarrow\;
\|X - B\|_F^2 = \|x_1 - b_1\|_2^2 + \cdots + \|x_n - b_n\|_2^2 .
\]

Let \(\{v_1, \dots, v_p\}\) be an orthonormal basis of the row subspace
of \(B\), and extend to an orthonormal basis
\(\{v_1, \dots, v_p, \dots, v_r\}\) for (the row space of) \(X\). Then
expand \[
b_i = \beta_1^i v_1 + \beta_2^i v_2 + \cdots + \beta_p^i v_p, \qquad
x_i = \chi_1^i v_1 + \chi_2^i v_2 + \cdots + \chi_r^i v_r ,
\] so \[
\|b_i - x_i\|_2^2 = (\beta_1^i - \chi_1^i)^2 + \cdots + (\beta_p^i - \chi_p^i)^2 + \sum_{l=p+1}^{r} (\chi_l^i)^2 ,
\] which is minimized when \(b_i\) is the projection of \(x_i\) onto the
subspace,
i.e.~\(\beta_1^i = \chi_1^i,\ \beta_2^i = \chi_2^i,\ \dots,\ \beta_p^i = \chi_p^i\).
(This can be visualized in the 3D case, where the distance between two
vectors is minimized when \(b_i\) is the projection of \(x_i\) onto the
plane.)

Since we are looking over all subspaces \(V\) of dimension at most
\(k\), there could be a specific subspace that is better. As claimed
previously, the greedy construction of singular vectors minimizes the
distance between the vectors and the orthogonal subspace, so
\(\{v_1, \dots, v_k\}\) should be the first \(k\) singular vectors. The
projection of \(X\) onto this subspace is \[
\hat{x}_j = (x_j^T v_1)v_1 + (x_j^T v_2)v_2 + \cdots + (x_j^T v_k)v_k , \qquad j = 1, \dots, n,
\] which gives \[
\hat{X} = \begin{pmatrix} -\,\hat{x}_1^T\,- \\ \vdots \\ -\,\hat{x}_n^T\,- \end{pmatrix}
= X \sum_{i=1}^{k} v_i v_i^T
= \sum_{i=1}^{k} \sigma_i u_i v_i^T = X_k .
\] Therefore \[
\|X - B\|_F \ge \|X - P_W^{p}(X)\|_F \ge \|X - P_W^{k}(X)\|_F \ge \|X - X_k\|_F .
\]
\end{proof}

\chapter{Clustering}\label{clustering}

Given \(X = \{x_1, x_2, \dots, x_n\} \subseteq \mathbb{R}^d\), find a
disjoint partition \(X = \bigcup_i X_i\), with
\(X_i \cap X_j = \emptyset\) for \(i \ne j\), such that points within
\(X_i\) are ``similar'' to each other and ``dissimilar'' to points in
\(X_j\) (\(i \ne j\)).

\begin{remark}
\leavevmode

\begin{itemize}
\tightlist
\item
  The ``similarity'' is usually measured by Euclidean distance
  \(\|x_i - x_j\|_2^2\), or by angle.

  \begin{itemize}
  \tightlist
  \item
    If \(\|x_i\| = 1\) (normalized), then similarity is measured by
    angle: \[
    \|x_i - x_j\|_2^2 = \|x_i\|^2 + \|x_j\|^2 - 2\langle x_i, x_j\rangle = 2(1 - \cos\theta)
    \;\Rightarrow\;
    \begin{cases} \cos\theta = 1, & \text{similar} \\ \cos\theta = -1, & \text{not similar} \end{cases}
    \]
  \item
    If not normalized, just use \(\|x_i - x_j\|_2^2\).
  \end{itemize}
\item
  In practice, the number of partitions \(k\) is given.
\end{itemize}

\end{remark}

\section{The Easy Case}\label{the-easy-case}

In an ``easy'' clustering case we have the following observation.

\textbf{Observation.} There exists a partition \(X = \bigcup_i X_i\)
such that \[
\max_{\substack{x_i, x_j \in \text{ same } X_k}} \|x_i - x_j\|_2
\;\le\;
\min_{\substack{x_i, x_j \in \text{ different } X_k}} \|x_i - x_j\|_2 .
\]

In this case, define the distance matrix \(D_{ij} = \|x_i - x_j\|_2\),
and let \(\text{LHS} < R < \text{RHS}\) of the inequality above. Then
there exists a permutation
\(\Pi : \{1, 2, \dots, n\} \to \{1, 2, \dots, n\}\) and a corresponding
permutation matrix \[
P_\Pi = \begin{pmatrix} e_{\Pi(1)}^{\mathsf T} \\ e_{\Pi(2)}^{\mathsf T} \\ \vdots \\ e_{\Pi(n)}^{\mathsf T} \end{pmatrix},
\qquad
e_i = \begin{pmatrix} 0 \\ \vdots \\ 1 \\ \vdots \\ 0 \end{pmatrix} \leftarrow i,
\] such that \[
D_\Pi = P_\Pi\, D\, P_\Pi^{\mathsf T}
\] is ``close'' to a block-diagonal matrix: within a diagonal block
\(D_{ij} < R\), and outside the diagonal blocks \(D_{ij} > R\).

\begin{remark}
The above simply means we can rearrange the points so that the points
that are ``close'' are grouped together (clustering).
\end{remark}

\section{\texorpdfstring{\(k\)-means
Clustering}{k-means Clustering}}\label{k-means-clustering}

For \(X = \{x_1, \dots, x_n\} \subseteq \mathbb{R}^d\) with \(k\)
partitions \(X = X_1 \cup X_2 \cup \dots \cup X_k\), and centers
\(\mu_1, \dots, \mu_k \in \mathbb{R}^d\), find the optimal \(\{X_i\}\)
and \(\{\mu_i\}\) such that \[
\min_{\{X_j\}} \left( \min_{\{\mu_j\}} \sum_{i=1}^{k} \sum_{x_j \in X_i} \|x_j - \mu_i\|_2^2 \right).
\]

\begin{lemma}[Optimal centers for a fixed
partition]\protect\hypertarget{lem-kmeans-center}{}\label{lem-kmeans-center}

For fixed \(\{X_j\}\), the minimizing centers are the cluster means: \[
\mu_j^* = \frac{1}{|X_j|} \sum_{x_i \in X_j} x_i,
\qquad
\{\mu_j^*\} = \underset{\{\mu_j\}}{\arg\min} \sum_{j} \sum_{x_i \in X_j} \|x_i - \mu_j\|_2^2 .
\]

\end{lemma}

\begin{proof}
Let \[
L = \sum_{j=1}^{k} \sum_{x_i \in X_j} \|x_i - \mu_j\|_2^2 = \sum_{j=1}^{k} \sum_{x_i \in X_j} \sum_{p} (x_{i,p} - \mu_{j,p})^2,
\] where \(x_{i,p}, \mu_{j,p}\) are the \(p\)-th components of
\(x_i, \mu_j \in \mathbb{R}^d\). Then \[
\frac{\partial L}{\partial \mu_{\ell, q}}
= \sum_{j=1}^{k} \sum_{x_i \in X_j} \sum_{p} 2(x_{i,p} - \mu_{j,p})\, \delta_{\ell j}\, \delta_{pq}
= \sum_{j=1}^{k} \sum_{x_i \in X_j} 2(x_{i,q} - \mu_{j,q})\, \delta_{\ell j}
= \sum_{x_i \in X_\ell} 2(x_{i,q} - \mu_{\ell,q}^*) = 0,
\] hence \[
\sum_{x_i \in X_\ell} x_{i,q} - |X_\ell|\, \mu_{\ell,q}^* = 0
\;\Rightarrow\;
\mu_{\ell,q}^* = \frac{1}{|X_\ell|} \sum_{x_i \in X_\ell} x_{i,q},
\qquad \text{i.e.} \quad
\mu_\ell^* = \frac{1}{|X_\ell|} \sum_{x_i \in X_\ell} x_i .
\]
\end{proof}

\begin{lemma}[Optimal partition for fixed
centers]\protect\hypertarget{lem-kmeans-partition}{}\label{lem-kmeans-partition}

For fixed \(\{\mu_j\}\), the minimizing partition assigns each point to
its nearest center: \[
X_i^* = \Big\{ x_j \in X : \|x_j - \mu_i\|_2 \le \|x_j - \mu_{i'}\|_2,\ \forall\, i' \ne i \Big\},
\] with ties broken arbitrarily.

\end{lemma}

\begin{proof}
For fixed \(\{\mu_i\}\), the objective \[
\min_{\{X_j\}} \sum_{i=1}^{k} \sum_{x_j \in X_i} \|x_j - \mu_i\|_2^2
= \min_{\{X_j\}} \Bigg[
\underbrace{\|x_1^1 - \mu_1\|_2^2 + \cdots + \|x_{|X_1|}^1 - \mu_1\|_2^2}_{x_i^1 \in X_1}
+ \underbrace{\|x_1^2 - \mu_2\|_2^2 + \cdots + \|x_{|X_2|}^2 - \mu_2\|_2^2}_{x_i^2 \in X_2}
+ \cdots
+ \underbrace{\|x_1^k - \mu_k\|_2^2 + \cdots + \|x_{|X_k|}^k - \mu_k\|_2^2}_{x_i^k \in X_k}
\Bigg].
\] Consider a point \(x_j \in X_i\) for which there exists \(\mu_k\)
with \(\|x_j - \mu_k\|_2^2 \le \|x_j - \mu_i\|_2^2\). Then a new
partition with \(X_i \leftarrow X_i - \{x_j\}\) and
\(X_k \leftarrow X_k + \{x_j\}\) drops the objective value. Repeating
this until no such case can be found, the objective is minimized when \[
X_i^* = \Big\{ x_j \in X : \|x_j - \mu_i\|_2 \le \|x_j - \mu_{i'}\|_2,\ \forall\, i' \ne i \Big\}.
\]
\end{proof}

\section{Lloyd's Algorithm}\label{lloyds-algorithm}

These two lemmas suggest an algorithm using \textbf{alternating
optimization}.

\begin{tcolorbox}[enhanced jigsaw, leftrule=.75mm, bottomrule=.15mm, left=2mm, colback=white, colframe=quarto-callout-note-color-frame, toprule=.15mm, breakable, arc=.35mm, rightrule=.15mm, opacityback=0]
\begin{minipage}[t]{5.5mm}
\textcolor{quarto-callout-note-color}{\faInfo}
\end{minipage}%
\begin{minipage}[t]{\textwidth - 5.5mm}

\textbf{Lloyd's algorithm}

\emph{Input:} \(X = \{x_1, \dots, x_n\} \subseteq \mathbb{R}^d\), and
\(k\).

\emph{Initialize:} Choose \(k\) random initial guesses for the centers
\(\mu_1^{(0)}, \dots, \mu_k^{(0)}\), and the correspondingly induced
initial cluster assignment \(\hat{C}_1(0), \dots, \hat{C}_k(0)\) by
associating each \(x_j\) to the closest center.

\emph{For \(t = 1, 2, \dots\) until convergence:}

\begin{enumerate}
\def\labelenumi{\arabic{enumi}.}
\tightlist
\item
  Use the current clusters \(\hat{C}_1(t-1), \dots, \hat{C}_k(t-1)\) to
  update the means: \[
  \mu_j(t) = \frac{1}{|\hat{C}_j(t-1)|} \sum_{x_i \in \hat{C}_j(t-1)} x_i .
  \]
\item
  Use the updated means to update the cluster assignments (assign
  \(x_i\) to its nearest cluster mean).
\end{enumerate}

\end{minipage}%
\end{tcolorbox}

\subsection{Convergence}\label{convergence}

Let \(X = X_1 \cup \dots \cup X_k\) and define the cost for a partition
with optimal centers \[
L(X_1, \dots, X_k) := \sum_{i=1}^{k} \sum_{x_j \in X_i} \Big\| x_j - \tfrac{1}{|X_i|} \sum_{x_k \in X_i} x_k \Big\|^2 .
\]

\begin{proposition}[Monotone
decrease]\protect\hypertarget{prp-lloyd-monotone}{}\label{prp-lloyd-monotone}

In Lloyd's algorithm, \[
L\big(X_1(t), \dots, X_k(t)\big) \ge L\big(X_1(t+1), \dots, X_k(t+1)\big).
\]

\end{proposition}

\begin{proof}
The updated (optimal) center is
\(\hat{\mu}_j(t+1) = \frac{1}{|X_j(t)|} \sum_{x_i \in X_j(t)} x_i\), and
the updated partition assigns each \(x_i\) to the nearest updated
\(\mu_j\), so \(\forall x_i \in X_j\),
\(\|x_i - \hat{\mu}_j\| \le \|x_i - \hat{\mu}_{j'}\|\) for \(j' \ne j\).
By Lemma~\ref{lem-kmeans-center} and Lemma~\ref{lem-kmeans-partition},
\(\{X_j\}\) in Lloyd's algorithm is the minimizer of \(L\), hence \[
L\big(X_1(t), \dots, X_k(t)\big) \ge L\big(X_1(t+1), \dots, X_k(t+1)\big).
\]
\end{proof}

\begin{corollary}[Finite
termination]\protect\hypertarget{cor-lloyd-terminate}{}\label{cor-lloyd-terminate}

Let \(L(X_1(t), \dots, X_k(t))\) be the cost of Lloyd's algorithm. Then
\(\lim_{t \to \infty} L\) exists and is attained in finite time.
Moreover, assuming all partitions have distinct cost (or the
tie-breaking rule is chosen properly), the partition itself converges in
finite time, i.e.~\(\exists\, T > 0\) such that \[
\big(X_1(T), \dots, X_k(T)\big) = \big(X_1(T'), \dots, X_k(T')\big), \qquad T' \ge T.
\]

\end{corollary}

\begin{proof}
The existence of the limit follows because the loss drops monotonically
(Proposition~\ref{prp-lloyd-monotone}) and is bounded below by \(0\).
Finite termination follows because the search space of partitions is
always finite.
\end{proof}

\subsection{Complexity}\label{complexity}

At each iteration, Lloyd's algorithm requires:

\begin{itemize}
\tightlist
\item
  \(O(knd)\) time to assign \(n\) data points in \(\mathbb{R}^d\) to
  \(k\) centers (via \(\|\mu_i - x_j\|_2^2\));
\item
  \(O(nd)\) time to update the mean centers.
\end{itemize}

\begin{remark}
\leavevmode

\begin{itemize}
\tightlist
\item
  Lloyd's algorithm is guaranteed to converge; however, it may only
  reach a \textbf{local} minimizer.
\item
  Random projection can be used to reduce time.
\item
  Lloyd's is sensitive to initialization and may converge to various bad
  local minima.
\end{itemize}

\end{remark}

\section{\texorpdfstring{Initialization:
\(k\)-means++}{Initialization: k-means++}}\label{initialization-k-means}

Since Lloyd's algorithm is initialization-sensitive, one improvement is
to initialize the mean centers as far apart as possible.

\begin{tcolorbox}[enhanced jigsaw, leftrule=.75mm, bottomrule=.15mm, left=2mm, colback=white, colframe=quarto-callout-note-color-frame, toprule=.15mm, breakable, arc=.35mm, rightrule=.15mm, opacityback=0]
\begin{minipage}[t]{5.5mm}
\textcolor{quarto-callout-note-color}{\faInfo}
\end{minipage}%
\begin{minipage}[t]{\textwidth - 5.5mm}

\textbf{Farthest-point initialization}

Pick \(\mu_1 \in \{x_1, \dots, x_N\}\) uniformly at random; set
\(T = \{\mu_1\}\).

While \(|T| \le k\):

\begin{itemize}
\tightlist
\item
  choose \(\mu_i \in X \setminus T\) that maximizes
  \(\min_{t \in T} \|x - t\|_2\);
\item
  \(T.\text{append}(\mu_i)\).
\end{itemize}

\end{minipage}%
\end{tcolorbox}

However, this initialization can be sensitive to outliers (an outlier
very far from the rest would always be picked). The improvement is to
add randomness, choosing farther points with higher probability.

\begin{tcolorbox}[enhanced jigsaw, leftrule=.75mm, bottomrule=.15mm, left=2mm, colback=white, colframe=quarto-callout-note-color-frame, toprule=.15mm, breakable, arc=.35mm, rightrule=.15mm, opacityback=0]
\begin{minipage}[t]{5.5mm}
\textcolor{quarto-callout-note-color}{\faInfo}
\end{minipage}%
\begin{minipage}[t]{\textwidth - 5.5mm}

\textbf{\(k\)-means++}

Pick \(\mu_1 \in X\) uniformly at random; set \(T = \{\mu_1\}\).

While \(|T| < k\):

\begin{itemize}
\tightlist
\item
  pick \(\mu_{|T|+1} \in X \setminus T\) with probability
  \(\propto \min_{t \in T} \|x - t\|_2\);
\item
  \(T \leftarrow T \cup \{\mu_{|T|+1}\}\).
\end{itemize}

\end{minipage}%
\end{tcolorbox}

\begin{remark}
\leavevmode

\begin{itemize}
\tightlist
\item
  The initialization steps take \(O(knd)\), since we need
  \(\|x_j - t\|\) for \(x_j \in X\) (\(n\) of them) and \(t \in T\)
  (\(k\) of them) --- same as one iteration of Lloyd's.
\item
  This initialization does not guarantee anything by itself; however, it
  can be shown that the initial cost is ``close'' to the optimal cost by
  a factor of \(\log(k)\) on average.
\end{itemize}

\end{remark}

\begin{proposition}[\(k\)-means++ approximation
guarantee]\protect\hypertarget{prp-kmeanspp-guarantee}{}\label{prp-kmeanspp-guarantee}

Let \(T \subseteq X\) be the initial means from \(k\)-means++, and
\(T^* \subseteq X\) the optimal \(k\) centers. With the induced
partition \(X_1 \cup X_2 \cup \dots \cup X_k = X\) chosen accordingly,
define \[
L(T) := \min_{X_1 \cup \dots \cup X_k} \sum_{\mu_i \in T} \sum_{x_j \in X_i} \|x_j - \mu_i\|_2^2 .
\] Then \[
\mathbb{E}\big(L(T)\big) \le c \cdot O(\log k)\, L(T^*), \qquad c \text{ some constant.}
\]

\end{proposition}

\section{Geometric Interpretation}\label{geometric-interpretation}

In Lloyd's algorithm the partition for fixed means is updated via \[
X_i = \big\{ x \in X : \|x - \mu_i\| \le \|x - \mu_j\|,\ j \ne i \big\}.
\] Noting that \(\|x - \mu_i\|_2^2 \le \|x - \mu_j\|_2^2\) for
\(i \ne j\), \[
\|x\|^2 + \|\mu_i\|^2 - 2\langle x, \mu_i\rangle \le \|x\|^2 + \|\mu_j\|^2 - 2\langle x, \mu_j\rangle
\;\Rightarrow\;
2\langle x, \mu_j - \mu_i\rangle \le \|\mu_j\|^2 + \|\mu_i\|^2 ,
\] which defines a half-space separated by a hyperplane. Lloyd's
algorithm thus clusters into the intersection of \(k-1\) half-planes.

\begin{remark}
\leavevmode

\begin{itemize}
\tightlist
\item
  It can be shown that the partitions in Lloyd's algorithm form a
  Voronoi diagram, and each cell is convex.
\item
  This means Lloyd's algorithm essentially works with linearly
  separable, convex clusters.
\end{itemize}

\end{remark}

\section{\texorpdfstring{Kernel Trick and Kernel
\(k\)-means}{Kernel Trick and Kernel k-means}}\label{kernel-trick-and-kernel-k-means}

One way to deal with non-convex, non-linear clusters is to use
\textbf{dimension increasing}: map the original data so that it becomes
more linearly separable and convex in a higher dimension.

\begin{example}[Lifting concentric
data]\protect\hypertarget{exm-circle-lift}{}\label{exm-circle-lift}

Consider data in \(\mathbb{R}^2\) (e.g.~two concentric rings). Create
the map \(\phi : \mathbb{R}^2 \to \mathbb{R}^3\), \[
\phi(x_1, x_2) = \begin{pmatrix} x_1^2 \\ \sqrt{2}\, x_1 x_2 \\ x_2^2 \end{pmatrix},
\] after which the two rings become linearly separable in
\(\mathbb{R}^3\).

\end{example}

However, boosting the dimension increases computational complexity. The
\textbf{kernel trick} bypasses this by allowing implicit computation of
the inner product of high-dimensional vectors via a low-dimensional
function.

Example: let \(x, y \in \mathbb{R}^2\) and
\(\phi(x) = (x_1^2, \sqrt{2}\, x_1 x_2, x_2^2)\),
\(\phi(y) = (y_1^2, \sqrt{2}\, y_1 y_2, y_2^2)\). Then \[
\langle \phi(x), \phi(y)\rangle = x_1^2 y_1^2 + 2 x_1 x_2 y_1 y_2 + x_2^2 y_2^2 = (x_1 y_1 + x_2 y_2)^2 = \langle x, y\rangle^2 =: K(x, y).
\]

\begin{definition}[Kernel
function]\protect\hypertarget{def-kernel}{}\label{def-kernel}

A symmetric map \(K : \mathbb{R}^d \times \mathbb{R}^d \to \mathbb{R}\)
is a \textbf{kernel induced by \(\phi\)} if there exists
\(\phi : \mathbb{R}^d \to \mathbb{R}^n\) such that \[
K(x, y) = \langle \phi(x), \phi(y)\rangle .
\]

\end{definition}

\begin{remark}
Examples:

\begin{itemize}
\tightlist
\item
  \(K(x, y) = \langle x, y\rangle^r\);
\item
  \(K(x, y) = (\langle x, y\rangle + c)^r\);
\item
  \(K(x, y) = \exp\!\left(-\dfrac{\|x - y\|^2}{2\sigma^2}\right)\) ---
  the \textbf{Gaussian kernel}.
\end{itemize}

For the Gaussian kernel, the inducing map is
\(\phi : \mathbb{R}^d \to \mathcal{H}\), where \(\mathcal{H}\) is a
Hilbert space.
\end{remark}

\begin{theorem}[Mercer's
theorem]\protect\hypertarget{thm-mercer}{}\label{thm-mercer}

Let \(X \subseteq \mathbb{R}^d\) be a compact set and
\(K : X \times X \to \mathbb{R}\) a continuous kernel. Define the linear
integral operator associated with \(K\) on \(L^2\),
\(T_K : L^2(X) \to L^2(X)\), as \[
T_K(f) := \int_X K(x, x')\, f(x')\, dx',
\] and suppose \(T_K\) is positive semidefinite, i.e. \[
\langle f, T_K f\rangle = \int_{X \times X} f(x)\, K(x, x')\, f(x')\, dx\, dx' \ge 0, \qquad \forall\, f \in L^2(X).
\] Then \(T_K\) has eigenfunctions \(\{\psi_i\}_{i=1,2,\dots} \in L^2\)
with corresponding eigenvalues \(\lambda_i \in \mathbb{R}^+\) such that
\[
K(x, x') = \sum_{i=1}^{\infty} \lambda_i\, \psi_i(x)\, \psi_i(x').
\]

\end{theorem}

\begin{remark}
This means we can view \(K(x, x')\) as induced by the feature map
\(\varphi(x) := \big(\sqrt{\lambda_i}\, \psi_i(x)\big)_{i=1}^{\infty}\).
\end{remark}

\subsection{\texorpdfstring{Kernel
\(k\)-means}{Kernel k-means}}\label{kernel-k-means}

In Lloyd's algorithm, the partition is updated based on the distance to
the \(i\)-th center,
\(\big\| x_j - \tfrac{1}{|X_i|} \sum_{x_i \in X_i} x_i \big\|_2^2\). We
first show this distance can be expressed via inner products: \[
\Big\| x - \tfrac{1}{|X_i|} \sum_{x_j \in X_i} x_j \Big\|_2^2
= \langle x, x\rangle + \frac{1}{|X_i|^2} \sum_{x_j, x_k \in X_i} \langle x_j, x_k\rangle - \frac{2}{|X_i|} \sum_{x_j \in X_i} \langle x_j, x\rangle .
\] If we use a feature map \(\phi(x)\) to boost the dimension, then the
distance in the transformed space can be expressed through the kernel:
\[
\Big\| \phi(x) - \tfrac{1}{|X_i|} \sum_{x_j \in X_i} \phi(x_j) \Big\|_2^2
= K(x, x) + \frac{1}{|X_i|^2} \sum_{x_j, x_k \in X_i} K(x_j, x_k) - \frac{2}{|X_i|} \sum_{x_j \in X_i} K(x_j, x) .
\]

We modify Lloyd's algorithm as follows to include the kernel.

\begin{tcolorbox}[enhanced jigsaw, leftrule=.75mm, bottomrule=.15mm, left=2mm, colback=white, colframe=quarto-callout-note-color-frame, toprule=.15mm, breakable, arc=.35mm, rightrule=.15mm, opacityback=0]
\begin{minipage}[t]{5.5mm}
\textcolor{quarto-callout-note-color}{\faInfo}
\end{minipage}%
\begin{minipage}[t]{\textwidth - 5.5mm}

\textbf{Kernel \(k\)-means}

\emph{Initialize:}
\(X_1^{(0)} \cup X_2^{(0)} \cup \dots \cup X_k^{(0)} = X\).

\emph{For \(t = 1, 2, \dots\) until convergence:}

\[
X_i^{(t+1)} = \Big\{ x_j \in X : \big\| \phi(x_j) - \tfrac{1}{|X_i|} \textstyle\sum \phi(x_k) \big\|_2 \le \big\| \phi(x_j) - \tfrac{1}{|X_{i'}|} \textstyle\sum \phi(x_m) \big\|_2,\ \forall\, i' \ne i \Big\},
\]

where each distance is computed via the kernel expression above.

\end{minipage}%
\end{tcolorbox}

\begin{remark}
\leavevmode

\begin{itemize}
\tightlist
\item
  The distance is calculated using the expression with the kernel.
\item
  The centers are now not explicitly computed in feature space, but
  implicitly via kernels.
\item
  In practice, the kernels are precomputed each iteration.
\item
  Each iteration of kernel Lloyd's takes \(O(N^2)\).
\end{itemize}

\end{remark}

\chapter{Classification}\label{classification}

Given labelled training data, the goal of a \textbf{classification}
problem is to predict a discrete label for each new (test) data point.
This chapter develops the \textbf{support vector machine (SVM)} for
binary classification, the Lagrangian \textbf{duality} theory behind its
kernelized (dual) form, the \textbf{soft-margin} extension for
non-separable data, and finally the reduction of \textbf{multi-class}
classification to a collection of binary problems.

\section{Support Vector Machine for Binary
Classification}\label{support-vector-machine-for-binary-classification}

\subsection{The binary classification
problem}\label{the-binary-classification-problem}

\begin{itemize}
\tightlist
\item
  \textbf{Train data:}
  \((\vec{x}_1, y_1), (\vec{x}_2, y_2), \dots, (\vec{x}_n, y_n)\) where
  \(\vec{x}_i \in \mathbb{R}^d\) and \(y_i \in \{1, -1\}\).
\item
  \textbf{Test data:} \(\vec{x}_1, \vec{x}_2, \dots \in \mathbb{R}^d\).
\item
  \textbf{Goal:} predict a label \(1\) or \(-1\) for each \(\vec{x}_i\).
\end{itemize}

\subsection{Linearly separable data}\label{linearly-separable-data}

A binary problem is \textbf{linearly separable} if there exists an
affine hyperplane \(\langle w, x\rangle - b\) that separates the data,
i.e. \[
\exists\, w \in \mathbb{R}^d,\ b \in \mathbb{R}, \quad\text{s.t.}\quad
\operatorname{sign}\!\big(\langle w, x_i\rangle - b\big) = y_i,
\qquad \forall\, i = 1, 2, \dots, n .
\]

For a linearly separable problem there could be \textbf{multiple}
choices of \(w\) and \(b\). A good choice for better generalization is
to choose a plane that \textbf{maximizes the margin}.

\subsection{Maximizing the margin}\label{maximizing-the-margin}

For linearly separable data, a separating plane
\(\langle w, x\rangle - b = 0\) satisfies \[
\begin{cases}
\langle w, x_i\rangle - b \ge 0, & y_i = 1, \\
\langle w, x_i\rangle - b \le 0, & y_i = -1 .
\end{cases}
\] We may push the plane out on both sides until it touches the boundary
of each class, giving the two \textbf{supporting hyperplanes} \[
\langle w, x\rangle - b = 1
\qquad\text{and}\qquad
\langle w, x\rangle - b = -1,
\qquad\text{s.t.}\qquad
\begin{cases}
\langle w, x_i\rangle - b \ge 1, & y_i = 1, \\
\langle w, x_i\rangle - b \le -1, & y_i = -1 .
\end{cases}
\] The perpendicular distance between these two hyperplanes --- the
\textbf{margin} --- is \[
\frac{2}{\|w\|_2} .
\]

\begin{figure}

\centering{

\pandocbounded{\includegraphics[keepaspectratio]{ml-ds/classification_files/figure-pdf/fig-svm-margin-output-1.pdf}}

}

\caption{\label{fig-svm-margin}Maximum-margin separating hyperplane. The
solid line is the decision boundary \(\langle w,x\rangle - b = 0\); the
two dashed lines are the supporting hyperplanes
\(\langle w,x\rangle - b = \pm 1\); the circled points are the support
vectors that lie on them. The margin (green) has width \(2/\|w\|_2\).}

\end{figure}%

We wish to \textbf{maximize} this distance, so the classification
objective is \[
\max_{w \in \mathbb{R}^d,\; b \in \mathbb{R}} \ \frac{2}{\|w\|_2}
\qquad\text{s.t.}\qquad
\begin{cases}
w^{\mathsf T} x_i - b \ge 1, & \text{if } y_i = 1, \\
w^{\mathsf T} x_i - b \le -1, & \text{if } y_i = -1 .
\end{cases}
\]

Maximizing \(2/\|w\|_2\) is the same as minimizing
\(\tfrac12\|w\|_2^2\), and the two constraints combine into the single
condition \(y_i(\langle w, x_i\rangle - b)\ge 1\).

\begin{definition}[Primal form of
SVM]\protect\hypertarget{def-svm-primal}{}\label{def-svm-primal}

\[
\min_{w,\,b} \ \tfrac12 \|w\|_2^2
\qquad\text{s.t.}\qquad
y_i\big(w^{\mathsf T} x_i - b\big) \ge 1,
\qquad \forall\, i = 1, \dots, n .
\]

\end{definition}

\begin{remark}
\leavevmode

\begin{itemize}
\tightlist
\item
  The constraint condition is equivalent to the two-sided condition
  above.
\item
  This problem can be solved by constrained optimization methods.
\end{itemize}

\end{remark}

\section{Duality of Constrained
Optimization}\label{duality-of-constrained-optimization}

To introduce non-linearity into the SVM one uses the \textbf{kernel
trick}. This requires the \textbf{dual} form of the SVM, so we first
review the duality of a general constrained optimization problem.

\subsection{The Lagrangian and the dual
function}\label{the-lagrangian-and-the-dual-function}

Consider the constrained optimization problem \[
\min_x \ f(x)
\qquad\text{s.t.}\qquad
\begin{cases}
f_i(x) \le 0, & i = 1, \dots, m, \\
h_i(x) = 0, & i = 1, \dots, p .
\end{cases}
\] As in Lagrangian penalty methods, construct the \textbf{Lagrangian}
\(L : D \times \mathbb{R}^m \times \mathbb{R}^p \to \mathbb{R}\), where
\(D\) is the joint domain of the constraints: \[
L(x, \lambda, v) = f(x) + \sum_{i} \lambda_i f_i(x) + \sum_i v_i h_i(x) .
\]

\begin{definition}[Lagrange dual
function]\protect\hypertarget{def-dual-function}{}\label{def-dual-function}

The function \(g : \mathbb{R}^m \times \mathbb{R}^p \to \mathbb{R}\)
defined by \[
g(\lambda, v) := \min_{x \in D} L(x, \lambda, v)
\] is called the \textbf{Lagrange dual function}. When \(L\) is
unbounded below, \(g \to -\infty\).

\end{definition}

\subsection{The primal value as a
min--max}\label{the-primal-value-as-a-minmax}

\begin{remark}
\leavevmode

\begin{itemize}
\tightlist
\item
  The optimal value \(p^\ast\) of the original problem is an
  \textbf{upper bound} of the dual function:
  \(g(\lambda, v) \le p^\ast\) for all \(\lambda \ge 0, v\).
\item
  The optimal value can be expressed as \[
  p^\ast = \min_x \ \max_{\lambda \ge 0,\, v} \ L(x, \lambda, v).
  \]
\end{itemize}

\end{remark}

The second point is valid through the following argument. Observe that
\[
\sup_{\lambda \ge 0,\, v} L(x, \lambda, v) =
\begin{cases}
f_0(x), & x \in D, \\
\infty, & \text{otherwise}.
\end{cases}
\] Indeed, if \(x \notin D\) then there exists \(i\) with \(f_i(x) > 0\)
or \(i\) with \(h_i(x) \ne 0\); either case lets us choose a proper
\(\lambda_i\) or \(v_i\) so that \(L \to \infty\). Thus \[
\min_x \ \max_{\lambda,\, v} L = \min_{x \in D} f_0(x) = p^\ast .
\]

\subsection{The min--max inequality}\label{the-minmax-inequality}

\begin{theorem}[Min--max
inequality]\protect\hypertarget{thm-minmax}{}\label{thm-minmax}

Let \(f : W \times Z \to \mathbb{R}\) with \(W \subseteq \mathbb{R}^n\)
and \(Z \subseteq \mathbb{R}^m\). Then \[
\min_{w \in W} \ \max_{z \in Z} f(w, z) \ \ge\ \max_{z \in Z} \ \min_{w \in W} f(w, z) .
\]

\end{theorem}

\begin{proof}
Let \(w^0 \in W\). Then \[
\max_{z \in Z} f(w^0, z) \ge f(w^0, z), \qquad \forall\, z \in Z .
\] Next let \(z^0 \in Z\). Then \[
\min_{w \in W} f(w, z^0) \le f(w', z^0), \qquad \forall\, w' \in W .
\] Combining, \[
\min_{w \in W} f(w, z^0) \le f(w^0, z^0) \le \max_{z \in Z} f(w^0, z) .
\] Let \(h(z) = \min_{w \in W} f(w, z)\) and
\(g(w) = \max_{z \in Z} f(w, z)\). Then \[
g(w) \ge h(z), \qquad \forall\, w \in W,\ z \in Z ,
\] so \(g(w) \ge \max_{z \in Z} h(z)\) for all \(w \in W\), and
therefore \[
\min_{w \in W} g(w) \ge \max_{z \in Z} h(z)
\quad\Longrightarrow\quad
\min_{w \in W} \max_{z \in Z} f(w, z) \ge \max_{z \in Z} \min_{w \in W} f(w, z) . \qquad\blacksquare
\]
\end{proof}

Applying this inequality to the primal problem yields the Lagrange dual
problem.

\subsection{The Lagrange dual problem}\label{the-lagrange-dual-problem}

\begin{definition}[Dual
problem]\protect\hypertarget{def-dual-problem}{}\label{def-dual-problem}

Let \(L\) be the Lagrangian of the primal problem. The corresponding
\textbf{dual problem} is \[
\max_{\lambda \ge 0,\, v} \ g(\lambda, v) .
\]

\end{definition}

\begin{remark}
From Theorem~\ref{thm-minmax}, the primal and dual are related by \[
p^\ast = \min_x \ \max_{\lambda \ge 0,\, v} L(x, \lambda, v)
\ \ge\
\max_{\lambda \ge 0,\, v} \ \min_x L(x, \lambda, v) = d^\ast ,
\] so the dual optimum \(d^\ast\) bounds the primal optimum \(p^\ast\)
from below.
\end{remark}

\textbf{Weak duality and duality gap.} When \(d^\ast < p^\ast\), we have
\textbf{weak duality}, and \(p^\ast - d^\ast\) is called the
\textbf{duality gap}.

\textbf{Strong duality and Slater's condition.} When \(d^\ast = p^\ast\)
holds, we have \textbf{strong duality}. A specific condition that must
hold for strong duality is called a \textbf{constraint qualification};
one simple example is Slater's condition.

\begin{proposition}[Slater's
condition]\protect\hypertarget{prp-slater}{}\label{prp-slater}

For the primal problem \[
\min_x \ f_0(x)
\qquad\text{s.t.}\qquad
\begin{cases}
f_i(x) \le 0, & i = 1, \dots, m, \\
A x = b,
\end{cases}
\] Slater's condition holds if

\begin{itemize}
\tightlist
\item
  \(f_0, f_1, \dots, f_m\) are convex, and
\item
  there exists \(x \in D\) such that \(f_i(x) < 0\) for
  \(i = 1, \dots, m\) and \(Ax = b\).
\end{itemize}

\end{proposition}

\section{Kernel SVM: The Dual Form}\label{kernel-svm-the-dual-form}

Back to the SVM objective: it is not hard to see that the primal SVM
satisfies Slater's condition. Thus we can find the dual form of the SVM
using the Lagrangian \[
L = \tfrac12 \|w\|_2^2 + \sum_i \lambda_i\Big(-\big(\langle w, x_i\rangle - b\big) y_i + 1\Big) .
\]

We first find \(\min_{w \in \mathbb{R}^d,\, b \in \mathbb{R}} L\) for
the Lagrange dual function. Setting the derivatives to zero, \[
\begin{cases}
\dfrac{\partial L}{\partial w_i} = w_i - \displaystyle\sum_j \lambda_j x_j^{\,i} y_j = 0, \\[2mm]
\dfrac{\partial L}{\partial b} = \displaystyle\sum_i y_i \lambda_i = 0,
\end{cases}
\qquad\Longrightarrow\qquad
\begin{cases}
w^\ast = \tfrac12 \displaystyle\sum_i y_i \lambda_i x_i, \\[2mm]
\displaystyle\sum_i y_i \lambda_i = 0 .
\end{cases}
\]

Substituting back, \[
\min_{w,\,b} L =
\begin{cases}
\tfrac12 \|w^\ast\|_2^2 + \displaystyle\sum_i \lambda_i\Big(-\big(\langle w^\ast, x_i\rangle - b\big) y_i + 1\Big), & \sum_i y_i \lambda_i = 0, \\
-\infty, & \sum_i y_i \lambda_i \ne 0 .
\end{cases}
\] (If \(\sum_i y_i \lambda_i \ne 0\), the term
\(-b \sum_i \lambda_i y_i\) can be driven to \(\to -\infty\).) The two
pieces evaluate to \[
\tfrac12 \|w^\ast\|_2^2 = \tfrac12 \langle w^\ast, w^\ast\rangle
= \tfrac18 \sum_{i,j} y_i y_j \lambda_i \lambda_j \langle x_i, x_j\rangle,
\] \[
\sum_i \lambda_i\Big(-\big(\langle w^\ast, x_i\rangle - b\big) y_i + 1\Big)
= -\sum_i \lambda_i \langle w^\ast, x_i\rangle y_i + \sum_i \lambda_i
= -\tfrac12 \sum_{i,j} \lambda_i \lambda_j y_i y_j \langle x_i, x_j\rangle + \sum_i \lambda_i .
\] Therefore \[
\min_{w,\,b} L =
\begin{cases}
-\tfrac14 \displaystyle\sum_{i,j} \lambda_i \lambda_j y_i y_j \langle x_i, x_j\rangle + \sum_i \lambda_i, & \sum_i y_i \lambda_i = 0, \\
-\infty, & \sum_i y_i \lambda_i \ne 0 .
\end{cases}
\]

\begin{definition}[Dual form of
SVM]\protect\hypertarget{def-svm-dual}{}\label{def-svm-dual}

\[
\max_{\lambda \ge 0} \ -\tfrac14 \sum_{i,j} \lambda_i \lambda_j y_i y_j \langle x_i, x_j\rangle + \sum_i \lambda_i
\qquad\text{s.t.}\qquad
\sum_i y_i \lambda_i = 0 .
\]

\end{definition}

\begin{remark}
In the dual form, \(x_i\) appears only as an \textbf{inner product}
\(\langle x_i, x_j\rangle\), and can be replaced with a \textbf{kernel
function} --- this is the kernel trick.
\end{remark}

\section{Soft Margin for Non-separable
Data}\label{soft-margin-for-non-separable-data}

In many cases the data are not (strictly) linearly separable. The
solution is to introduce a \textbf{``soft'' margin} that allows
misclassification. The constraints are modified to \[
y_i\big(\langle w, x_i\rangle - b\big) \ge 1 - \xi_i .
\]

\begin{remark}
\leavevmode

\begin{itemize}
\tightlist
\item
  Instead of the \textbf{hard} constraint that all data of a class lie
  on one side of the hyperplane, soft margins allow a set of hyperplanes
  that permit some violations.
\item
  We would hope the slack variables \(\xi_i\) are not too large.
\end{itemize}

\end{remark}

\subsection{Primal problem with soft
margin}\label{primal-problem-with-soft-margin}

\[
\min_{w,\,b,\,\xi} \ \tfrac12 \|w\|_2^2 + C \sum_i \xi_i
\qquad\text{s.t.}\qquad
y_i\big(\langle w, x_i\rangle - b\big) \ge 1 - \xi_i, \quad \xi_i \ge 0 .
\]

Here \(C\) (written \(\lambda\) in one place in the source) is the
\textbf{regularization parameter}.

\subsection{Dual problem with soft
margin}\label{dual-problem-with-soft-margin}

Following a similar procedure to the hard-constraint problem, the dual
form of the soft margin is \[
\max_{\lambda \ge 0} \ \sum_i \lambda_i - \tfrac14 \sum_{i=1}^{n} \sum_j y_i y_j \lambda_i \lambda_j \langle x_i, x_j\rangle
\qquad\text{s.t.}\qquad
\sum_i \lambda_i y_i = 0 \ \ \&\ \ \lambda_i \le \frac{1}{2n\lambda} .
\]

\section{Multi-class Classification}\label{multi-class-classification}

The multi-class classification problem can be reduced to the binary case
using the \textbf{``one-vs-all''} scheme.

For a multi-class problem the data set is
\((\vec{x}_1, y_1), \dots, (\vec{x}_n, y_n)\) as before, but now
\(y_i \in \{0, 1, \dots, C-1\}\), for a total of \(C\) distinct classes.

\subsection{One-vs-all scheme}\label{one-vs-all-scheme}

The goal of a multi-class classifier is: given a data point, decide that
it is in the \(i\)-th class and not in any other class. Therefore, for
each class we can construct a \textbf{binary} classification problem ---
\emph{is the data point in this class or not?} --- which gives a
decision boundary between points in this class and those not in it.
Doing this for all classes yields a classification region based on these
boundaries.

\textbf{Example (3 classes).} We build three binary classifiers, one for
each of classes \(1, 2, 3\). For each binary classifier we create the
one-vs-all label \[
\hat{y}_i =
\begin{cases}
1, & \text{if } y_i \text{ is in class } c, \\
-1, & \text{if not},
\end{cases}
\qquad c = 1, 2, 3 .
\] This creates \(3\) decision boundaries, one per class.

Each line separates the space into a region ``\(\in C\)'' and the other
side ``\(\notin C\)''. As before, if the weight for the classifier of
class \(c\) is \(\vec{w}_c\), then \[
\begin{cases}
x_i^{\mathsf T} w_c - b > 0 \ \Rightarrow\ x_i \in C, \\
x_i^{\mathsf T} w_c - b < 0 \ \Rightarrow\ x_i \notin C .
\end{cases}
\] This separates the space into subdomains.

\subsection{Decision regions}\label{decision-regions}

Combining the binary boundaries, a test point falls into one of three
cases.

\textbf{(1) Positive side of exactly one boundary.} The straightforward
case: the point is on the positive side of exactly one boundary and on
the negative side of all others. Then the predicted class is the one on
whose positive side it lies: \[
\hat{y}_i = \arg\max_c \big(x_i^{\mathsf T} w_c - b\big) .
\]

\textbf{(2) Positive side of more than one boundary.} The point could
plausibly belong to two classes; we should choose the one that is more
``plausible''. Note that for a binary classifier, a point
\textbf{further} from the boundary is considered \textbf{more
confident}, whereas a point \textbf{closer} to the boundary is more
likely to be mistakenly labelled. Thus we use the \textbf{distance to
the decision boundary} as a decision indicator: the class is the
boundary that is \textbf{furthest}. The distance from \(x_i\) to the
boundary for class \(c\) is \[
\frac{x_i^{\mathsf T} w_c - b_c}{\|w_c\|_2},
\] so \[
\hat{y}_i = \arg\max_{c = 0, \dots, C-1} \ \frac{x_i^{\mathsf T} w_c - b_c}{\|w_c\|_2} .
\]

\textbf{(3) Negative side of all boundaries.} We decide the class as the
\emph{opposite} of case (2): when a point is further from a boundary, it
is more confident of \textbf{not} being in that class, so we assign the
class whose boundary has the \textbf{smallest} distance. Since
\(\dfrac{x_i^{\mathsf T} w_c - b_c}{\|w_c\|_2}\) is the \textbf{signed}
distance, \[
\hat{y}_i
= \arg\min_c \ -\frac{x_i^{\mathsf T} w_c - b_c}{\|w_c\|_2}
= \arg\max_c \ \frac{x_i^{\mathsf T} w_c - b_c}{\|w_c\|_2} .
\]

All three cases therefore reduce to the same rule: predict the class
with the largest signed distance.

\begin{figure}

\centering{

\pandocbounded{\includegraphics[keepaspectratio]{ml-ds/classification_files/figure-pdf/fig-multiclass-output-1.pdf}}

}

\caption{\label{fig-multiclass}Multi-class one-vs-all decision regions
for \(3\) classes. Each pair of regions is separated by a linear
boundary; a test point is assigned to the class with the largest
normalized signed distance
\(\dfrac{x^{\mathsf T} w_c - b_c}{\|w_c\|_2}\).}

\end{figure}%

\subsection{General procedure}\label{general-procedure}

The 3-class example extends to the general multi-class case, summarized
as follows.

\begin{tcolorbox}[enhanced jigsaw, leftrule=.75mm, bottomrule=.15mm, left=2mm, colback=white, colframe=quarto-callout-note-color-frame, toprule=.15mm, breakable, arc=.35mm, rightrule=.15mm, opacityback=0]
\begin{minipage}[t]{5.5mm}
\textcolor{quarto-callout-note-color}{\faInfo}
\end{minipage}%
\begin{minipage}[t]{\textwidth - 5.5mm}

\vspace{-3mm}\textbf{One-vs-all multi-class classification}\vspace{3mm}

\textbf{Input:} \(\{(\vec{x}_i, y_i)\}\) with
\(y_i \in \{0, 1, \dots, C-1\}\).

For \(c = 0, \dots, C-1\):

\begin{enumerate}
\def\labelenumi{\arabic{enumi}.}
\tightlist
\item
  Relabel each sample \(p\) by \[
  \hat{y}_p =
  \begin{cases}
  +1, & \text{if } y_p = c, \\
  -1, & \text{else}.
  \end{cases}
  \]
\item
  Solve the binary subproblem
  \(\{(\vec{x}_p, \hat{y}_p)\} \to (w_c, b_c)\).
\item
  Normalize the weights:
  \(\displaystyle w_c \leftarrow \frac{w_c}{\|w_c\|_2}, \quad b_c \leftarrow \frac{b_c}{\|w_c\|_2}\).
\end{enumerate}

\textbf{Predict} the label by
\(\hat{y} = \arg\max_c \ x^{\mathsf T} w_c\).

\end{minipage}%
\end{tcolorbox}

\begin{remark}
\leavevmode

\begin{itemize}
\tightlist
\item
  The normalization does not change the decision in case (1).
\item
  In practice one can use \textbf{one-hot encoding}, e.g.~a class-\(1\)
  target \[
  \begin{pmatrix} -1 \\ 1 \\ -1 \\ \vdots \end{pmatrix}
  \]
\end{itemize}

\end{remark}

\chapter{Regression}\label{regression}

Given a training set of predictor--response pairs, the goal of a
\textbf{regression} problem is to learn a function \(f\) that predicts a
(continuous) response from the predictors. This chapter develops
\textbf{linear} and \textbf{ridge} regression, their solution via the
normal equations, the extension to \textbf{polynomial}
(linear-in-parameters) regression, and finally \textbf{kernel}
regression together with the Woodbury identity that makes it
computationally tractable.

\section{Linear (Ridge) Regression}\label{linear-ridge-regression}

\subsection{The general regression
task}\label{the-general-regression-task}

\begin{itemize}
\tightlist
\item
  Let \(y_i \in \mathbb{R}\), \(i = 1, 2, \dots, N\) be the
  \textbf{response variables}.
\item
  Let \(x_i \in \mathbb{R}^d\), \(i = 1, 2, \dots, N\) be the
  \textbf{predictors}.
\item
  The pairs \(\{(x_i, y_i)\}\) form the \textbf{training set}.
\end{itemize}

The regression task wishes to find a function
\(f : \mathbb{R}^d \to \mathbb{R}\) such that \(f(x_i) \approx y_i\),
and that generalizes to the test set. This requires two ingredients:

\begin{itemize}
\tightlist
\item
  a \textbf{function form} (parametric methods) --- i.e.~assume a family
  for \(f\);
\item
  a \textbf{measure / metric} for \(f(x_i) \approx y_i\) (the
  \textbf{loss function}).
\end{itemize}

\subsection{Least squares (linear
regression)}\label{least-squares-linear-regression}

Take \(f\) to be \textbf{linear}: \[
y = w^{\mathsf T} x + b, \qquad w \in \mathbb{R}^d,\ b \in \mathbb{R},
\] or simply \(y = w^{\mathsf T} x\) if we absorb the bias by extending
\[
x \ \longrightarrow\ \begin{pmatrix} x_1 \\ \vdots \\ x_d \\ 1 \end{pmatrix} .
\]

For the \textbf{metric}, use the \textbf{least-square error} \[
\min_{w \in \mathbb{R}^d} \ \sum_i \big(y_i - \langle w, x_i\rangle\big)^2
\qquad\longrightarrow\qquad \text{error between predicted and actual,}
\] or, in matrix form, \(\|Y - Xw\|_2^2\) if we denote \[
Y = \begin{pmatrix} y_1 \\ \vdots \\ y_N \end{pmatrix} \in \mathbb{R}^N,
\qquad
X = \begin{pmatrix} x_1^{\mathsf T} \\ x_2^{\mathsf T} \\ \vdots \\ x_N^{\mathsf T} \end{pmatrix} \in \mathbb{R}^{N \times d} .
\]

\subsection{Ridge regression
(regularization)}\label{ridge-regression-regularization}

In some cases \(f\) could \textbf{overfit} (e.g.~matching noise). This
can be reduced by adding a \textbf{regularization term} to the loss
function, as a penalty for too-large coefficients: \[
\min_{w \in \mathbb{R}^d} \ \|Y - Xw\|_2^2 + \lambda \|w\|_2^2, \qquad \lambda \ge 0 .
\]

From optimization theory, the resulting minimization problem is a
\textbf{strictly convex} function (for \(\lambda > 0\)), so every local
minimum is the global minimum.

\begin{proposition}[Ridge regression solution (normal
equations)]\protect\hypertarget{prp-ridge-solution}{}\label{prp-ridge-solution}

For \(\lambda > 0\), the ridge objective has the unique minimizer \[
w^\ast = \arg\min_w \ \|Xw - Y\|_2^2 + \lambda \|w\|_2^2
= \big(X^{\mathsf T} X + \lambda I\big)^{-1} X^{\mathsf T} y .
\]

\end{proposition}

\begin{proof}
(Sketch.) Set the gradient to zero,
\(\dfrac{\partial L}{\partial w_i} = 0\), and check the Hessian
\(\dfrac{\partial^2 L}{\partial w_i \partial w_j}\). This is
straightforward since \(L\) is in \textbf{quadratic form} (and the
\(\lambda I\) term makes \(X^{\mathsf T} X + \lambda I\) positive
definite, hence invertible). \(\blacksquare\)
\end{proof}

\begin{remark}
\textbf{Choice of regularizer.}

\begin{itemize}
\tightlist
\item
  \(\lambda \|w\|_1 = \lambda \sum_i |w_i|\) --- promotes a
  \textbf{sparse} weight vector.
\item
  \(\|w\|_\ast = \sum_i \sigma_i(w)\) (the nuclear norm) --- promotes a
  \textbf{low-rank} matrix.
\end{itemize}

\end{remark}

\begin{remark}
\textbf{Multiple-label predictors.} If \(y_i \in \mathbb{R}^p\)
(multiple labels), we simply extend to a matrix of weights
\(W \in \mathbb{R}^{p \times d}\): \[
y_i = W x_i,
\qquad
\begin{pmatrix} y_1^i \\ \vdots \\ y_p^i \end{pmatrix}
=
\begin{pmatrix}
w_{11} & w_{12} & \cdots & w_{1d} \\
\vdots & & & \vdots \\
w_{p1} & w_{p2} & \cdots & w_{pd}
\end{pmatrix}
\begin{pmatrix} x_1^i \\ \vdots \\ x_d^i \end{pmatrix} ,
\] and with \[
Y = \begin{pmatrix} y_1^{\mathsf T} \\ \vdots \\ y_N^{\mathsf T} \end{pmatrix},
\qquad
X = \begin{pmatrix} x_1^{\mathsf T} \\ \vdots \\ x_N^{\mathsf T} \end{pmatrix},
\qquad
L = \|X W^{\mathsf T} - Y\|_F^2 + \lambda \|W\|_F^2 .
\]
\end{remark}

\section{Polynomial Regression}\label{polynomial-regression}

\(f\) need \textbf{not} be linear, as long as it is \textbf{linear in
its parameters}.

\textbf{Example (3rd-order polynomial fit).} \[
y = w_3^{\mathsf T} x^3 + w_2^{\mathsf T} x^2 + w_1^{\mathsf T} x
= \big(\, w_3^{\mathsf T} \mid w_2^{\mathsf T} \mid w_1^{\mathsf T} \,\big)
\begin{pmatrix} x^3 \\ x^2 \\ x \end{pmatrix} ,
\] where \(x^k\) denotes the monomial (elementwise-power) feature
expansion. This \textbf{shares the same form} as linear regression, so
the same solution applies.

\section{Kernel Regression}\label{kernel-regression}

The kernel trick can be applied to the features \(x\). Suppose
\(\phi : \mathbb{R}^d \to \mathbb{R}^p\) is a \textbf{feature map} that
transforms the original data to \[
\big(\phi(x_1), y_1\big), \big(\phi(x_2), y_2\big), \dots, \big(\phi(x_n), y_n\big)
\in \mathbb{R}^p \times \mathbb{R} .
\] Then perform linear ridge regression in the transformed feature
space: \[
\min_{w \in \mathbb{R}^p} \ \sum_i \big(w^{\mathsf T} \phi(x_i) - y_i\big)^2 + \lambda \|w\|_2^2 .
\] From Proposition~\ref{prp-ridge-solution}, the optimal solution is \[
w^\ast = \big(\Phi^{\mathsf T}\Phi + \lambda I_p\big)^{-1} \Phi^{\mathsf T} y,
\qquad
\Phi = \begin{pmatrix} \phi(x_1)^{\mathsf T} \\ \phi(x_2)^{\mathsf T} \\ \vdots \\ \phi(x_N)^{\mathsf T} \end{pmatrix}
\in \mathbb{R}^{n \times p} .
\] Just as before, directly computing \(\phi\) and \(w^\ast\) is very
costly when \(p\) is large. We want to apply the kernel trick, which
relies on the following identity.

\subsection{The Woodbury identity}\label{the-woodbury-identity}

\begin{lemma}[Woodbury matrix
identity]\protect\hypertarget{lem-woodbury}{}\label{lem-woodbury}

\[
\big(\Phi^{\mathsf T}\Phi + \lambda I_p\big)^{-1} \Phi^{\mathsf T}
= \Phi^{\mathsf T} \big(\Phi\Phi^{\mathsf T} + \lambda I_n\big)^{-1} .
\]

\end{lemma}

\begin{proof}
Take the SVD \(\Phi = U \Sigma V^{\mathsf T}\) (so
\(\Phi^{\mathsf T} = V \Sigma^{\mathsf T} U^{\mathsf T}\),
\(\Phi^{\mathsf T}\Phi = V \Sigma^{\mathsf T}\Sigma V^{\mathsf T}\), and
\(\Phi\Phi^{\mathsf T} = U \Sigma \Sigma^{\mathsf T} U^{\mathsf T}\)).

\textbf{Left-hand side.} Using \(U^{\mathsf T} U = I\) and
\(I_p = V V^{\mathsf T}\), \[
\begin{aligned}
\text{LHS}
&= \big((V\Sigma^{\mathsf T} U^{\mathsf T})(U\Sigma V^{\mathsf T}) + \lambda I_p\big)^{-1} V\Sigma^{\mathsf T} U^{\mathsf T} \\
&= \big(V\Sigma^{\mathsf T}\Sigma V^{\mathsf T} + \lambda V V^{\mathsf T}\big)^{-1} V\Sigma^{\mathsf T} U^{\mathsf T} \\
&= \big(V(\Sigma^{\mathsf T}\Sigma + \lambda I_p) V^{\mathsf T}\big)^{-1} V\Sigma^{\mathsf T} U^{\mathsf T} \\
&= V(\Sigma^{\mathsf T}\Sigma + \lambda I_p)^{-1} V^{\mathsf T} V\Sigma^{\mathsf T} U^{\mathsf T}
= V(\Sigma^{\mathsf T}\Sigma + \lambda I_p)^{-1}\Sigma^{\mathsf T} U^{\mathsf T} .
\end{aligned}
\]

\textbf{Right-hand side.} Using \(I_n = U U^{\mathsf T}\), \[
\begin{aligned}
\text{RHS}
&= V\Sigma^{\mathsf T} U^{\mathsf T}\big((U\Sigma V^{\mathsf T})(V\Sigma^{\mathsf T} U^{\mathsf T}) + \lambda I_n\big)^{-1} \\
&= V\Sigma^{\mathsf T} U^{\mathsf T}\big(U\Sigma\Sigma^{\mathsf T} U^{\mathsf T} + \lambda U U^{\mathsf T}\big)^{-1} \\
&= V\Sigma^{\mathsf T} U^{\mathsf T}\big(U(\Sigma\Sigma^{\mathsf T} + \lambda I_n) U^{\mathsf T}\big)^{-1} \\
&= V\Sigma^{\mathsf T} U^{\mathsf T} U(\Sigma\Sigma^{\mathsf T} + \lambda I_n)^{-1} U^{\mathsf T}
= V\Sigma^{\mathsf T}(\Sigma\Sigma^{\mathsf T} + \lambda I_n)^{-1} U^{\mathsf T} .
\end{aligned}
\]

A detailed calculation on
\(\Sigma^{\mathsf T}(\Sigma\Sigma^{\mathsf T} + \lambda I_n)^{-1}\) and
\((\Sigma^{\mathsf T}\Sigma + \lambda I_p)^{-1}\Sigma^{\mathsf T}\)
shows that both equal the same \(p \times n\) ``pseudo-diagonal'' matrix
with entries \(\dfrac{\sigma_i}{\sigma_i^2 + \lambda}\). Thus
\(\text{LHS} = \text{RHS}\). \(\blacksquare\)
\end{proof}

\subsection{Kernel ridge regression
prediction}\label{kernel-ridge-regression-prediction}

\begin{theorem}[Kernel ridge regression
predictor]\protect\hypertarget{thm-kernel-prediction}{}\label{thm-kernel-prediction}

Using Lemma~\ref{lem-woodbury}, the prediction at a new point \(x\) is
\[
y = w^{\ast\mathsf T}\phi(x)
= \phi(x)^{\mathsf T} \Phi^{\mathsf T}\big(\Phi\Phi^{\mathsf T} + \lambda I_n\big)^{-1} y
= k_x^{\mathsf T}\big(K + \lambda I_n\big)^{-1} y ,
\] where \(K_{ij} = \phi(x_i)^{\mathsf T}\phi(x_j) = K(x_i, x_j)\) is
the \textbf{kernel (Gram) matrix}, and \(k_x = \Phi\phi(x)\) has entries
\((k_x)_i = \phi(x_i)^{\mathsf T}\phi(x) = K(x_i, x)\).

\end{theorem}

\begin{proof}
Expanding the products directly, \[
\Phi\Phi^{\mathsf T}
= \begin{pmatrix} \phi(x_1)^{\mathsf T} \\ \vdots \\ \phi(x_n)^{\mathsf T} \end{pmatrix}
\big(\phi(x_1), \ \dots, \ \phi(x_n)\big)
= \begin{pmatrix}
\phi(x_1)^{\mathsf T}\phi(x_1) & \phi(x_1)^{\mathsf T}\phi(x_2) & \cdots \\
\vdots & \ddots & \\
\end{pmatrix}
= \big(K(x_i, x_j)\big) = K ,
\] and \[
\Phi\phi(x)
= \begin{pmatrix} \phi(x_1)^{\mathsf T} \\ \vdots \\ \phi(x_n)^{\mathsf T} \end{pmatrix} \phi(x)
= \begin{pmatrix} \phi(x_1)^{\mathsf T}\phi(x) \\ \vdots \\ \phi(x_n)^{\mathsf T}\phi(x) \end{pmatrix}
= \begin{pmatrix} K(x_1, x) \\ \vdots \\ K(x_n, x) \end{pmatrix} = k_x . \qquad\blacksquare
\]
\end{proof}

\subsection{Computational cost}\label{computational-cost}

\begin{tcolorbox}[enhanced jigsaw, leftrule=.75mm, bottomrule=.15mm, left=2mm, colback=white, colframe=quarto-callout-note-color-frame, toprule=.15mm, breakable, arc=.35mm, rightrule=.15mm, opacityback=0]
\begin{minipage}[t]{5.5mm}
\textcolor{quarto-callout-note-color}{\faInfo}
\end{minipage}%
\begin{minipage}[t]{\textwidth - 5.5mm}

\vspace{-3mm}\textbf{Cost of kernel ridge regression}\vspace{3mm}

Using the kernel trick, the cost breaks down as:

\begin{itemize}
\tightlist
\item
  \textbf{Precompute} \(K \in \mathbb{R}^{n \times n}\) --- cost
  \(O(n^2)\) kernel evaluations.
\item
  \textbf{Precompute} \((K + \lambda I_n)^{-1} y\) --- cost \(O(n^3)\)
  flops and \(O(n)\) storage.
\item
  \textbf{Compute} \(k_x\) --- cost \(O(n)\) kernel evaluations.
\item
  \textbf{Compute} \(k_x^{\mathsf T}(K + \lambda I_n)^{-1} y\) --- cost
  \(O(n)\) flops.
\end{itemize}

\end{minipage}%
\end{tcolorbox}

\chapter{Markov Chains}\label{markov-chains}

A \textbf{discrete-time Markov chain} is a sequence of random variables
in which each step depends only on the previous one (the \textbf{Markov
property}), which greatly simplifies the calculation of long-time
transition probabilities. This chapter covers the definition and
transition matrix, \(n\)-step transition probabilities, communicating
classes and irreducibility, stationary (invariant) distributions,
periodicity, and the convergence theorems --- the last of which underpin
Markov Chain Monte Carlo (MCMC) methods.

\section{Discrete-Time Markov Chains}\label{discrete-time-markov-chains}

\begin{definition}[Discrete-time Markov
chain]\protect\hypertarget{def-markov-chain}{}\label{def-markov-chain}

A sequence of random variables \((X_n)_{n \ge 0}\), where
\(X_n : \Omega \to S \subseteq \mathbb{Z}\), is called a \textbf{Markov
chain} with \textbf{initial distribution} \(\lambda\) if

\begin{enumerate}
\def\labelenumi{\alph{enumi})}
\tightlist
\item
  \(X_0\) has distribution \(\lambda\);
\item
  it satisfies the \textbf{Markov property} \[
  P\big(X_{n+1} = i_{n+1} \mid X_n = i_n, \dots, X_0 = i_0\big)
  = P\big(X_{n+1} = i_{n+1} \mid X_n = i_n\big) .
  \]
\end{enumerate}

\end{definition}

\begin{remark}
\leavevmode

\begin{itemize}
\tightlist
\item
  \(X_n\) denotes the \textbf{state} of the chain at time \(n\).
\item
  \(S\) is the \textbf{state space} --- all possible outcomes of each
  state.
\item
  Writing \(P(X_{n+1} = i_{n+1} \mid X_n = i_n) = P_{i_n i_{n+1}}\),
  then for finite \(S\) the matrix \(P \in \mathbb{R}^{|S| \times |S|}\)
  is the \textbf{transition matrix}, and \(P_{i_n i_{n+1}}\) is the
  transition probability from state \(i_n\) to \(i_{n+1}\).
\item
  Being probabilities, the entries satisfy \[
  P_{ij} \ge 0, \qquad \sum_{j \in S} P_{ij} = 1 .
  \]
\end{itemize}

\end{remark}

\section{\texorpdfstring{\(n\)-Step Transition
Probabilities}{n-Step Transition Probabilities}}\label{n-step-transition-probabilities}

The Markov property and the total-probability formula let us compute the
\(n\)-step transition probability from the single-step transition
matrix. We wish to calculate the probability that the chain is at state
\(j\) after \(n\) steps starting from \(i\), and, more generally, the
probability of being at \(j\) after \(n\) steps from any initial state.

\begin{theorem}[\(n\)-step transition
probability]\protect\hypertarget{thm-nstep}{}\label{thm-nstep}

Let \((X_n)_{n \ge 0}\) be a Markov chain with initial distribution
\(\lambda\) and transition matrix \(P\). Then for all \(n, m \ge 0\):

\begin{enumerate}
\def\labelenumi{\arabic{enumi}.}
\tightlist
\item
  \(P(X_n = j \mid X_0 = i) = (P^n)_{ij}\);
\item
  \(P(X_n = j) = (\lambda P^n)_j\).
\end{enumerate}

\end{theorem}

\begin{proof}
The calculation is straightforward using total probability and the
Markov property. For (2), sum over all intermediate paths and apply the
Markov property repeatedly: \[
\begin{aligned}
P(X_n = j)
&= \sum_{i_0 \in S} \cdots \sum_{i_{n-1} \in S}
P\big(X_n = j \mid X_{n-1} = i_{n-1}, \dots, X_0 = i_0\big)\, P\big(X_{n-1} = i_{n-1}, \dots, X_0 = i_0\big) \\
&= \sum_{i_0} \cdots \sum_{i_{n-1}}
P\big(X_n = j \mid X_{n-1} = i_{n-1}\big)\, P\big(X_{n-1} = i_{n-1} \mid X_{n-2} = i_{n-2}\big) \cdots P\big(X_1 = i_1 \mid X_0 = i_0\big)\, P(X_0 = i_0) \\
&= \sum_{i_0} \cdots \sum_{i_{n-1}} \lambda_{i_0}\, P_{i_0 i_1} P_{i_1 i_2} \cdots P_{i_{n-1} j}
= (\lambda P^n)_j .
\end{aligned}
\] For (1), condition on \(X_0 = i\) (so \(\lambda\) is replaced by a
point mass at \(i\)): \[
P(X_n = j \mid X_0 = i)
= \sum_{i_1 \in S} \cdots \sum_{i_{n-1} \in S} P_{i\, i_1} P_{i_1 i_2} \cdots P_{i_{n-1} j}
= (P^n)_{ij} . \qquad \blacksquare
\]
\end{proof}

\begin{remark}
\leavevmode

\begin{itemize}
\tightlist
\item
  Using the matrix form assumes a finite state space; the summation
  identity, however, remains true when the state space is infinite.
\item
  Statement (1) extends to the \textbf{Chapman--Kolmogorov} form \[
  P(X_{n+m} = j \mid X_n = i) = (P^m)_{ij} .
  \]
\end{itemize}

\end{remark}

\section{Communicating Classes and
Irreducibility}\label{communicating-classes-and-irreducibility}

We now examine the asymptotic behavior of the state distribution: as
\(n \to \infty\), does the chain tend to an equilibrium? Different
states can behave very differently. The following example illustrates
this.

\begin{figure}

\centering{

\pandocbounded{\includegraphics[keepaspectratio]{ml-ds/markovChain_files/figure-pdf/fig-markov-classes-output-1.pdf}}

}

\caption{\label{fig-markov-classes}A 4-state chain. State \(1\) is
\textbf{transient}: it has a self-loop (\(0.2\)) and leads to state
\(4\) (\(0.5\)) and to state \(2\) (\(0.3\)). State \(4\) is
\textbf{absorbing} (self-loop \(1\)), forming the closed class
\(\{4\}\). States \(\{2,3\}\) form a second \textbf{closed}
communicating class. Once the chain leaves \(1\) it is trapped in
\(\{4\}\) or \(\{2,3\}\) and never returns to \(1\).}

\end{figure}%

For the chain above, note that:

\begin{enumerate}
\def\labelenumi{\arabic{enumi}.}
\tightlist
\item
  State \(1\) can access \(2, 3, 4\), while \(2, 3, 4\) cannot access
  \(1\).
\item
  Once the chain enters state \(2\), it only transitions within
  \(\{2, 3\}\); similarly \(4\) is absorbing.
\item
  When \(n\) is large enough, the state is ``trapped'' in \(\{4\}\) or
  \(\{2, 3\}\).
\end{enumerate}

The following definitions formalize this observation.

\begin{definition}[Leads to /
accessible]\protect\hypertarget{def-accessible}{}\label{def-accessible}

We say state \(i\) \textbf{leads to} \(j\) (equivalently, \(j\) is
\textbf{accessible} from \(i\)) if there exists a finite sequence
\(i_1, i_2, \dots, i_{n-1}\) such that \[
P_{i\, i_1} > 0, \ P_{i_1 i_2} > 0, \ \dots, \ P_{i_{n-1} j} > 0
\quad\Longleftrightarrow\quad
P(X_n = j \mid X_0 = i) > 0 \ \text{ for some } n \ge 0 .
\]

\end{definition}

\begin{definition}[Communicate]\protect\hypertarget{def-communicate}{}\label{def-communicate}

We say \(i\) and \(j\) \textbf{communicate} if \(i\) leads to \(j\) and
\(j\) leads to \(i\).

\end{definition}

\begin{definition}[Communicating class and closed
class]\protect\hypertarget{def-classes}{}\label{def-classes}

~

\begin{itemize}
\tightlist
\item
  A set of states is a \textbf{communicating class} (or
  \textbf{recurrent class}) if every pair \(i, j\) in it communicate.
  (Equivalently, a \textbf{strongly connected component} when the chain
  is viewed as a graph.)
\item
  A communicating class \(C\) is \textbf{closed} if, for every
  \(i \in C\), whenever \(i\) leads to \(j\) then \(j \in C\).
\end{itemize}

\end{definition}

\begin{remark}
\leavevmode

\begin{itemize}
\tightlist
\item
  Once the chain enters a closed class, it never leaves.
\item
  An \textbf{absorbing state} is simply a singleton \(\{i\}\) that is a
  closed class.
\item
  If a state \(i\) is not in any (closed) communicating class, then
  there is a state \(j\) accessible from \(i\) but with \(i\)
  \textbf{not} accessible from \(j\); for large \(n\) the \(i \to j\)
  transition occurs and the chain never returns. Hence we will focus on
  chains with a \textbf{single} communicating class, called
  \textbf{irreducible}.
\end{itemize}

\end{remark}

\section{Stationary (Invariant)
Distributions}\label{stationary-invariant-distributions}

For an irreducible chain, does the probability of being in each state
stabilize in the long run? By ``stable'' we mean that the state
distribution does not change after a transition --- hence the following
definition.

\begin{definition}[Stationary / invariant
distribution]\protect\hypertarget{def-stationary}{}\label{def-stationary}

Let \(\lambda\) denote the probability distribution of state \(X\), i.e.
\(\lambda_i = P(X = i)\). We say \(\lambda\) is a \textbf{stationary}
(or \textbf{invariant}) distribution if \[
\lambda = \lambda P, \qquad \sum_i \lambda_i = 1 .
\]

\end{definition}

\begin{remark}
The definition means that a one-step transition of the current
distribution, \[
\sum_j \lambda_j P_{ji} = \lambda_i ,
\] leaves it unchanged.
\end{remark}

It can be shown that if the transition probability converges to a number
independent of the initial state, then the limiting distribution is
stationary/invariant (for finite state space).

\begin{theorem}[A convergent limit is
invariant]\protect\hypertarget{thm-converge-invariant}{}\label{thm-converge-invariant}

Let the state space \(S\) of a Markov chain \((X_n)_{n \ge 0}\) be
finite, and suppose that for some \(i \in S\), \[
P(X_n = j \mid X_0 = i) = (P^n)_{ij} \to \pi_j \quad\text{as } n \to \infty, \ \ \forall j \in S .
\] Then \(\pi\) is an invariant distribution.

\end{theorem}

\begin{proof}
Using \((P^n)_{ij} = \sum_{k \in S} (P^{n-1})_{ik} P_{kj}\) and
interchanging limit and (finite) sum, \[
\pi_j = \lim_{n \to \infty} (P^n)_{ij}
= \lim_{n \to \infty} \sum_{k \in S} (P^{n-1})_{ik} P_{kj}
= \sum_{k \in S} \Big(\lim_{n \to \infty} (P^{n-1})_{ik}\Big) P_{kj}
= \sum_{k \in S} \pi_k P_{kj} ,
\] so \(\pi = \pi P\). Also \[
\sum_{j \in S} \pi_j
= \sum_{j \in S} \lim_{n \to \infty} (P^n)_{ij}
= \lim_{n \to \infty} \sum_{j \in S} (P^n)_{ij} = 1 .
\] Thus \(\pi\) is a stationary distribution. \(\blacksquare\)
\end{proof}

\begin{remark}
\leavevmode

\begin{itemize}
\tightlist
\item
  The interchange of limit and summation is permitted for a finite state
  space.
\item
  The theorem guarantees that a \textbf{convergent} distribution is
  invariant; the converse need not hold (invariance does not imply
  convergence --- see below).
\item
  For the probability to converge we must restrict to a \textbf{single}
  communicating class; otherwise \(\lim_{n \to \infty}(P^n)_{ij}\)
  depends on the initial state, as different initial states may lead to
  different recurrent classes.
\end{itemize}

\end{remark}

\subsection{Periodicity}\label{periodicity}

Even for a single communicating class, existence of an invariant
distribution does \textbf{not} guarantee convergence.

\begin{figure}

\centering{

\pandocbounded{\includegraphics[keepaspectratio]{ml-ds/markovChain_files/figure-pdf/fig-markov-periodic-output-1.pdf}}

}

\caption{\label{fig-markov-periodic}A 2-state chain with
\(P_{12} = P_{21} = 1\)
(i.e.~\(P = \left(\begin{smallmatrix} 0 & 1 \\ 1 & 0 \end{smallmatrix}\right)\)).
The distribution \(\pi = (\tfrac12, \tfrac12)\) is invariant, yet the
probabilities \textbf{oscillate} and never converge, since
\(P^{2k} = I\) and \(P^{2k+1} = P\).}

\end{figure}%

To exclude such cases we define the aperiodic property.

\begin{definition}[Aperiodic]\protect\hypertarget{def-aperiodic}{}\label{def-aperiodic}

A state \(i\) is \textbf{aperiodic} if \((P^n)_{ii} > 0\) for all
sufficiently large \(n\). A Markov chain is \textbf{aperiodic} if every
state \(i \in S\) is aperiodic.

\end{definition}

\section{Convergence to Equilibrium}\label{convergence-to-equilibrium}

\begin{theorem}[Irreducible + one aperiodic state \(\Rightarrow\) all
states
aperiodic]\protect\hypertarget{thm-irreducible-aperiodic}{}\label{thm-irreducible-aperiodic}

Suppose a Markov chain is irreducible and has an aperiodic state \(i\).
Then for all states \(j, k \in S\), \((P^n)_{jk} > 0\) for all
sufficiently large \(n\). In particular, all states are aperiodic.

\end{theorem}

\begin{proof}
Since the chain is irreducible, there exist \(r, s\) with
\((P^r)_{ji} > 0\) and \((P^s)_{ik} > 0\). Then for sufficiently large
\(n\), \[
\begin{aligned}
(P^{\,n + r + s})_{jk}
&= \sum_{i_1, \dots, i_n \in S} (P^r)_{j\, i_1}\, P_{i_1 i_2} P_{i_2 i_3} \cdots P_{i_{n-1} i_n}\, (P^s)_{i_n k} \\
&= \sum_{i_1, i_n \in S} (P^r)_{j\, i_1}\, (P^n)_{i_1 i_n}\, (P^s)_{i_n k} \\
&\ge (P^r)_{ji}\, (P^n)_{ii}\, (P^s)_{ik} > 0 ,
\end{aligned}
\] choosing \(i_1 = i_n = i\), since \((P^n)_{ii} > 0\) (\(i\)
aperiodic) and the end factors are positive (irreducibility).
\(\blacksquare\)
\end{proof}

For an aperiodic chain, the existence of an invariant distribution
guarantees convergence to it.

\begin{theorem}[Convergence to the invariant
distribution]\protect\hypertarget{thm-convergence}{}\label{thm-convergence}

Suppose \((X_n)_{n \ge 0}\) is an irreducible and aperiodic Markov chain
with transition matrix \(P\) and initial distribution \(\lambda\). If
\(P\) has an invariant distribution \(\pi\), then \[
P(X_n = j) \to \pi_j \quad\text{as } n \to \infty, \ \ \forall j \in S .
\] In particular, \((P^n)_{ij} \to \pi_j\) as \(n \to \infty\) for all
\(i, j \in S\).

\end{theorem}

\begin{theorem}[Finite irreducible aperiodic
chains]\protect\hypertarget{thm-unique-stationary}{}\label{thm-unique-stationary}

Suppose \((X_n)_{n \ge 0}\) is an irreducible and aperiodic Markov chain
with a \textbf{finite} state space. Then there exists a \textbf{unique}
invariant distribution \(\pi\), and \(P(X_n = i) \to \pi_i\) for
\textbf{any} initial distribution \(\lambda\).

\end{theorem}

\part{Machine Learning in Finance}

\part{Practice Machine Learning with Kaggle}

\part{Statistical Inference}

\part{Regression Analysis}

\chapter{Statistical Inference of Simple Linear
Regression}\label{statistical-inference-of-simple-linear-regression}

\part{Bayesian Modeling}

\part{Stochastic Process and Stochastic Integral}

\chapter{Probability in Measure Theory
Language}\label{probability-in-measure-theory-language}

\textbf{Sample space \(\Omega\):} the set of all possible outcome of an
experiment.

\textbf{Event:} a subsets of the sample space \(\Omega\).

\begin{itemize}
\tightlist
\item
  \(\varnothing\) is impossible event
\item
  \(\Omega\) is certain event.
\end{itemize}

\begin{tcolorbox}[enhanced jigsaw, leftrule=.75mm, bottomrule=.15mm, left=2mm, colback=white, colframe=quarto-callout-note-color-frame, toprule=.15mm, breakable, arc=.35mm, rightrule=.15mm, opacityback=0]
\begin{minipage}[t]{5.5mm}
\textcolor{quarto-callout-note-color}{\faInfo}
\end{minipage}%
\begin{minipage}[t]{\textwidth - 5.5mm}

\textbf{Bertrand's Paradox:} probability of chord intersect with inner
circle.

An example of different

\end{minipage}%
\end{tcolorbox}

\textbf{Remark:} Events are subsets of \(\Omega\), but subsets of
\(\Omega\) may not be an Event

we can therefor select subsets of \(\Omega\), which are events and form
a collection \(\mathcal{F}\) (i.e.,
\(\mathcal{F} \subseteq 2^\Omega \leftarrow\) set of all subset of
\(\Omega\)) and should have following property

\begin{itemize}
\tightlist
\item
  \begin{enumerate}
  \def\labelenumi{(\alph{enumi})}
  \tightlist
  \item
    if \(A, B \in \mathcal{F}\) then \(A \cup B \in \mathcal{F}\) and
    \(A \cap B \in \mathcal{F}\)
  \end{enumerate}
\item
  \begin{enumerate}
  \def\labelenumi{(\alph{enumi})}
  \setcounter{enumi}{1}
  \tightlist
  \item
    if \(A \in \mathcal{F}\) then \(A^c \in \mathcal{F}\).
  \end{enumerate}
\item
  \begin{enumerate}
  \def\labelenumi{(\alph{enumi})}
  \setcounter{enumi}{2}
  \tightlist
  \item
    \(\varnothing \in \mathcal{F}\)
  \end{enumerate}
\end{itemize}

\(\mathcal{F} \subseteq 2^\Omega\) satisfied above is called a field.

\textbf{Remark:}

\begin{itemize}
\tightlist
\item
  \begin{enumerate}
  \def\labelenumi{(\alph{enumi})}
  \tightlist
  \item
    \(\to \bigcup_i A_i \in \mathcal{F}\),
    \(\bigcap_i A_i \in \mathcal{F} \to\) close under finite union or
    intersect.
  \end{enumerate}
\item
  \begin{enumerate}
  \def\labelenumi{(\alph{enumi})}
  \setcounter{enumi}{1}
  \tightlist
  \item
    This properties can be shown via easy coin toss or dice experiment
  \end{enumerate}
\end{itemize}

\textbf{Infinite union and \(\cap\)'s:} Consider an experiment. A coin
is repeatedly tossed until the first head show up, try to find number of
toss before this happen

\[
\Omega = \{\omega_1, \omega_2, \dots\}, \qquad \omega_i := \text{first } i{-}1 \text{ are tail}, \ i \text{ are head}
\]

Event \(A :=\) the first head arrive after even number of tosses.

\[
A := \{\omega_2, \omega_4, \omega_6 \dots\} = \bigcup_{i=2} \omega_i
\]

\(A\) should also in the event space \(\mathcal{F}\) \(\to\)
\(\mathcal{F}\) should closed under infinite \(\bigcup_i \ \bigcap_i\)

\section{\texorpdfstring{\(\sigma\)-algebra / \(\sigma\)-field \(\to\)
set of
Events.}{\textbackslash sigma-algebra / \textbackslash sigma-field \textbackslash to set of Events.}}\label{sigma-algebra-sigma-field-to-set-of-events.}

collection of subset of sample space \(\mathcal{F} \subseteq 2^\Omega\),
is called a \(\sigma\)-algebra / \(\sigma\)-field if it satisfied
following

\begin{enumerate}
\def\labelenumi{\arabic{enumi}.}
\tightlist
\item
  \(\varnothing \in \mathcal{F}\)
\item
  if \(A \in \mathcal{F}\), \(A^c \in \mathcal{F}\)
  \(\ (A^c := \Omega - A) \to \Omega \in \mathcal{F}\)
\item
  if \(A_1, A_2 \cdots \in \mathcal{F}\), then
  \(\bigcup_i^\infty A_i \in \mathcal{F}\),
  \(\bigcap_i^\infty A_i \in \mathcal{F}\)
\end{enumerate}

\section{Probability measure}\label{probability-measure}

suppose we have sample space \(\Omega\), and \(\sigma\)-algebra
\(\mathcal{F}\). A probability measure is a function
\(P : \mathcal{F} \to [0, 1]\) satisfying.

\begin{enumerate}
\def\labelenumi{\arabic{enumi}.}
\tightlist
\item
  \(P(\varnothing) = 0\), \(P(\Omega) = 1\)
\item
  if \(A_1, A_2, \cdots \in \mathcal{F}\) and are disjoint
  \((A_i \cap A_j = \varnothing, \text{ for } \forall\, i \neq j)\) then
\end{enumerate}

\[
P\!\left(\bigcup_i^\infty A_i\right) = \sum_{i=1}^{k} P(A_i)
\]

The \((\Omega, \mathcal{F}, P)\) is called a probability space.

\textbf{lemma:}

\begin{itemize}
\tightlist
\item
  \begin{enumerate}
  \def\labelenumi{(\alph{enumi})}
  \tightlist
  \item
    \(P(A^c) = 1 - P(A)\)
  \end{enumerate}
\item
  \begin{enumerate}
  \def\labelenumi{(\alph{enumi})}
  \setcounter{enumi}{1}
  \tightlist
  \item
    if \(A \subseteq B\), \(P(A) \leq P(B)\)
  \end{enumerate}
\item
  \begin{enumerate}
  \def\labelenumi{(\alph{enumi})}
  \setcounter{enumi}{2}
  \tightlist
  \item
    \(P(A \cup B) = P(A) + P(B) - P(A \cap B)\)
  \end{enumerate}
\item
  \begin{enumerate}
  \def\labelenumi{(\alph{enumi})}
  \setcounter{enumi}{3}
  \tightlist
  \item
    \(P\!\left(\bigcup_{i=1}^n A_i\right) = \sum P(A_i) - \sum_{i<j} P(A_i \cap A_j) + \cdots (-1)^{n+1} P(A_1 \cap A_2 \cdots A_n)\)
  \end{enumerate}
\end{itemize}

\section{\texorpdfstring{Generated \(\sigma\)-field \(\to\) for
uncountable
\(\Omega\).}{Generated \textbackslash sigma-field \textbackslash to for uncountable \textbackslash Omega.}}\label{generated-sigma-field-to-for-uncountable-omega.}

suppose a collection of subset \(A_\alpha \subseteq \Omega\), where
\(\alpha \in T\) not necessarily countable. The smallest
\(\sigma\)-field \(\mathcal{F}\) s.t. \(A_\alpha \in \mathcal{F}\), for
\(\forall\, \alpha \in T\) is denoted as \(\sigma(\{A_\alpha\})\).

Then \(A_\alpha\) is called generator, and \(\sigma(\{A_\alpha\})\) is
called \(\sigma\)-field generated by \(\{A_\alpha\}\)

\[
\sigma\{A_\alpha\} = \bigcap \left\{\, \mathcal{G} : \mathcal{G} \subseteq 2^\Omega \text{ is a } \sigma\text{-field}, \ A_\alpha \in \mathcal{G}, \ \forall\, \alpha \in T \,\right\}
\]

\textbf{Remark:} if \(A_\alpha\) is \(\sigma\)-field for each
\(\alpha \in T\). \(\bigcap_{\alpha \in T} A_\alpha\) is a \(\sigma\)

\section{\texorpdfstring{Borel \(\sigma\)-field on
\(\mathbb{R}\).}{Borel \textbackslash sigma-field on \textbackslash mathbb\{R\}.}}\label{borel-sigma-field-on-mathbbr.}

a Borel \(\sigma\)-field on \(\mathbb{R}\) is a \(\sigma\)-field
generated by \(\{(a,b), \ a<b \in \mathbb{R}\}\)

\[
\mathcal{B} = \sigma\big(\{(a,b), \ a<b \in \mathbb{R}\}\big) \iff \sigma\big(\{O \subseteq \mathbb{R} \text{ open}\}\big) \iff \sigma\big(\{(-\infty, b], \ b \in \mathbb{R}\}\big)
\]

\section{Random Variable (R.V.)}\label{random-variable-r.v.}

A R.V. is a function \(X : \Omega \to \mathbb{R}\) s.t.
\(\forall\, \alpha \in \mathbb{R}\) the set
\(\{\omega : X(\omega) \leq \alpha\}\) is in \(\mathcal{F}\) ( such
function is also called \(\mathcal{F}\)-measurable ) (or just
measurable)

equivently, the measurable means \(X\) is measurable. in \(\mathcal{F}\)
if for \(\forall\, B \in \mathcal{B}(\mathbb{R})\),
\(X^{-1}(B) \in \mathcal{F}\)

\textbf{Remark:} The R.V. create a mapping between outcomes and real
numbers, in a way s.t. out comes that captures by any borel set (could
be record as an event of random real space) is forming an events in
event space of \(\Omega\).

\section{\texorpdfstring{\(\sigma\)-field generated by random
variable.}{\textbackslash sigma-field generated by random variable.}}\label{sigma-field-generated-by-random-variable.}

suppose \(X\) is a R.V. on \((\Omega, \mathcal{F})\) the
\(\sigma\)-field generated by \(X\), \(\sigma(X)\) si the smallest
\(\sigma\)-field \(\mathcal{G} \subseteq \mathcal{F}\), s.t. \(X\) is
still measurable on \((\Omega, \mathcal{G})\).

\textbf{Remark:} \(\{X^{-1}(B) \mid B \in \mathcal{B}\}\) is such
smallest \(\sigma\)-algebra. (by definition, \(X\) is such function
\(\Omega \to \mathbb{R}\), s.t. every \(B \in \mathcal{B}(\mathbb{R})\)
is inside \(\mathcal{F}\), the smallest set from \(X\) is
\(\{X^{-1}(B) \mid B \in \mathcal{B}\}\), only need to show
\(\{X^{-1}(B) \mid B \in \mathcal{B}\}\) is \(\sigma\)-field)

\textbf{Remark:} The mapping of R.V. is chosen s.t.
\(\omega \in \Omega\) enclosed by any \(B \in \mathcal{B}\) is belong to
\(\sigma\)-field \(\mathcal{F}\). Thus \(\mathcal{F}\) contains all the
information relavent to \(X\). the \(\sigma(X)\) is thus the least amout
of information relevant to \(X\).

\section{Simple function :}\label{simple-function}

\[
X := \sum_{n=n}^{N} C_n \mathbf{1}_{A_n} \quad \text{is a R.V. for finite } N \text{ and } C_n \in \mathbb{R}
\]

here \(A_n \in \mathcal{F}\) for any \(n\) and \[
\mathbf{1}_{A_n} := \begin{cases} 1, & \omega \in A_n \\ 0, & \omega \notin A_n. \end{cases}
\]

\textbf{prop:} for every R.V. \(X\), there exist a sequence of simple
functions \(X_n\) s.t. \(X_n(\omega) \to X(\omega)\) as \(n \to \infty\)
for \(\forall\, \omega \in \Omega\).

\textbf{Remark:} The R.V. \(X(\omega)\) is usually denote simple as
\(X\) and sometime \(P(X^{-1}(B))\) as \(P(X \in B)\)

\section{Two R.V. same (almost
surely)}\label{two-r.v.-same-almost-surely}

suppose \(X\) and \(Y\) are R.V. defined on
\((\Omega, \mathcal{F}, P)\), we says \(X\) and \(Y\) are almost surely
(a.s.) the same if \(P(\omega : X(\omega) \neq Y(\omega)) = 0\).

\textbf{Remark:} in many case, \(X \overset{a.s.}{=} Y\) is simply
denote as \(X = Y\).

\begin{itemize}
\tightlist
\item
  The a.s. notation applies to many other cases such as \(X \geq 0\)
  a.s. means \(P(\{\omega : X(\omega) < 0\}) = 0\).
\end{itemize}

\section{Borel (measurable) function}\label{borel-measurable-function}

A R.V. \(g\) defined on \((\mathbb{R}, \mathcal{B})\) is called a Borel
function.

\section{\texorpdfstring{Extend Borel-\(\sigma\)-algebra to
\(\mathbb{R}^n\).}{Extend Borel-\textbackslash sigma-algebra to \textbackslash mathbb\{R\}\^{}n.}}\label{extend-borel-sigma-algebra-to-mathbbrn.}

A Borel-\(\sigma\)-field in \(\mathbb{R}^n\) is simply \[
\mathcal{B}_n = \sigma\big( [a_1, b_1] \times [a_2, b_2] \times \cdots \times [a_n, b_n]; \ a_i, b_i \in \mathbb{R}, \ i = 1, \dots, n \big)
\]

\section{Expected Value.}\label{expected-value.}

Integration w.r.t. measure , lebesgue integral

\begin{enumerate}
\def\labelenumi{\arabic{enumi}.}
\item
  Simple R.V. Def

  suppose \(X\) on \((\Omega, \mathcal{F}, P)\) is Simple value R.V.
  \(X = \sum_{i=1}^{k} C_i \mathbf{1}_{A_i}\), then define \[
  \int_\Omega X \, dP := \sum_{i=1}^{k} C_i P(A_i)
  \]
\item
  if \(X\) is non negative R.V. \[
  \int_\Omega X \, dP := \sup_{Y \leq X, \, Y \text{ simple}} \int_\Omega Y \, dP.
  \]
\item
  if \(X : \Omega \to \mathbb{R}\), then \[
  \int_\Omega X \, dP = \int_\Omega X^+ \, dP - \int_\Omega X^- \, dP. \qquad X^+ := \max(X, 0), \ X^- := \max(-X, 0)
  \]
\end{enumerate}

\textbf{Remark:}

\begin{itemize}
\tightlist
\item
  The integration is based on spliting the range of \(X(\omega)\) to
  small equal interval and approximating \(X(\omega)\) by a constant on
  the correspondy set of \(\omega\). values s.t. \(X(\omega)\) fall into
  the interval.
\item
  Contrast to Rieman where the domain of integration is split into
  intervals, limited to \(\mathbb{R}^n\).
\item
  lebesgue integral allows to deal with general domains.
\item
  If the correspondy Riemann integral exist, then it coincide with
  lebesgue integral.
\end{itemize}

\section{Conditional probability}\label{conditional-probability}

let \((\Omega, \mathcal{F}, P)\) be a prob space and.
\(A, B \in \mathcal{F}\) and \(P(B) > 0\).

The conditional probability of \(A\), given \(B\) is \[
P(A \mid B) := \frac{P(A \cap B)}{P(B)}
\]

\textbf{motivation}

want to find probability of \(A\), given \(B\), so suppose
\(\omega \in \Omega\) and we know \(\omega \in B\), sure that, \(B\) can
be considered as a new probability space,
\((\Omega', \mathcal{F}', \hat{P})\), where \(\Omega' = B\),
\(\mathcal{F}' := \{C \cap B \mid C \in \mathcal{F}\}\) and
\(\hat{P} := \frac{P}{P(B)}\) so that \(\hat{P}(B) = 1\) then the
probability \(\omega\) lies in \(A\) also is
\(\hat{P}(A \cap B) = \frac{P(A \cap B)}{P(B)}\)

\section{Independence}\label{independence}

This means \(P(A \mid B) = P(A)\) (know \(B\) don't give \(A\)). \[
\Rightarrow P(A \cap B) = P(A) P(B)
\]

\section{independent Random Variable}\label{independent-random-variable}

let \(X_i : \Omega \to \mathbb{R}^n\) be random variables
\((i = 1, \dots)\). They are independent if for all \(k \geq 2\),
\(k \in \mathbb{Z}\) and all choice of Borel set
\(B_1, \cdots B_k \subseteq \mathbb{R}^n\). \[
P(X_1 \in B_1, X_2 \in B_2, \cdots X_k \in B_k) = P(X_1 \in B_1) P(X_2 \in B_2) \cdots P(X_k \in B_k).
\] this essentrally means \(\sigma\)-field generated by each \(X_i\) \[
\sigma(X_i)
\]

\section{Conditional Expectation}\label{conditional-expectation}

\(E(X \mid B)\) the expected value of R.V. \(X\), given the event \(B\)
happens.

From previous, we can think \(B\) form a new probability space, with
\(\hat{P} = \frac{P}{P(B)}\)
\(\Rightarrow E(X \mid B) = \frac{1}{P(B)} \int_B X \, dP\).

\(E(X \mid Y)\), the expected value of R.V., given another R.V. \(Y\)
known.

consider the case with in \((\Omega, \mathcal{F}, P)\) where \(Y, X\) is
defined \(Y = \sum_{i=1}^{m} c_i \mathbf{1}_{A_i}\), suppose \(c_i\) are
distinct and \(A_1, A_2 \cdots A_m\) are disjoint.

\[
Y(\omega) = \begin{cases} c_1, & \omega \in A_1 \\ c_2, & \omega \in A_2 \\ \ \vdots \\ c_n, & \omega \in A_n \end{cases}
\qquad \text{now if we know } Y(\omega) \to \text{ which event } A_i
\]

\[
\text{thus } E(X \mid Y) = \begin{cases} \dfrac{1}{P(A_1)} \displaystyle\int_{A_1} X \, dP, & A_1 \\ \ \vdots \\ \dfrac{1}{P(A_n)} \displaystyle\int_{A_n} X \, dP, & A_n \end{cases}
\]

\begin{itemize}
\tightlist
\item
  \(E(X \mid Y)\) is R.V.
\item
  \(E(X \mid Y)\) is \(\sigma(Y)\) measurable
\item
  \(\int_A E(X \mid Y) \, dP = \int_A X \, dP \qquad \forall\, A \in \sigma(Y)\)
\end{itemize}

\textbf{Def:} let \(X, Y\) be R.V. defined on \(\Omega\), the
conditional expectation of \(X\) given \(Y\) (\(E(X \mid Y)\)) is a
\(\sigma(Y)\) measurable random variable \(Z\) s.t. \[
\int_A X \, dP = \int_A Z \, dP \qquad \text{for } \forall\, A \in \sigma(Y)
\]

\textbf{Remark:} Such \(Z\) is unique up to set of probability zero.
thus \(Z\) can just be denote as \(E(X \mid Y)\)

\textbf{Def:} \((\Omega, \mathcal{F}, P)\) be a probability space and
\(\mathcal{V}\) is a \(\sigma\)-field,
\(\mathcal{V} \subseteq \mathcal{F}\), if
\(X : \Omega \to \mathbb{R}^n\) is an random variable, then
\(E(X \mid \mathcal{V}) :=\) any R.V. on \(\Omega\) s.t.

\begin{enumerate}
\def\labelenumi{(\roman{enumi})}
\item
  \(E(X \mid \mathcal{V})\) is \(\mathcal{V}\)-measurable.
\item
  \(\int_A X \, dP = \int_A E(X \mid \mathcal{V}) \, dP\), for
  \(\forall\, A \in \mathcal{V}\)
\end{enumerate}

\textbf{Remark:} \(E(X \mid \mathcal{V})\) can be understand as given
the (information) available in a \(\sigma\)-field \(\mathcal{V}\), from
which we intend to build an estimate of R.V. \(X\)

\begin{enumerate}
\def\labelenumi{(\roman{enumi})}
\item
  \(\to E(X \mid \mathcal{V})\) has to be constructed from information
  in \(\mathcal{V}\).
\item
  \(\to\) The estimation should be consist with \(X\), at least as
  regards in integration.
\end{enumerate}

\textbf{Th:} let \(X\) be an integrable R.V. Then for each
\(\sigma\)-field \(\mathcal{V} \subseteq \mathcal{F}\), the
\(E(X \mid \mathcal{V})\) exists and is unique up to
\(\mathcal{V}\)-measurable set of probability zero.

\chapter{Martingales}\label{martingales}

\textbf{fair game:} Consider \(Y_1, Y_2 \cdots\) are independent
``fair'' gambling game, \(E(Y_i) = 0\).

\(S_n := Y_1 + Y_2 \cdots \cdots + Y_n\) as accumulated. total wining

\(Q\): what is the best guess of \(S_{n+k}\) given
\(S_1, S_2 \cdots S_n\)?

\[
\begin{aligned}
E(S_{n+k} \mid S_1, S_2, \cdots S_n)
&= E(Y_1 + Y_2 \cdots + Y_{n+k} \mid S_1, S_2, \cdots S_n) \\
&= \underline{E(Y_1 + Y_2 + \cdots + Y_n \mid S_1, \cdots S_n)} + E(Y_{n+1} + \cdots Y_{n+k} \mid S_1, \cdots S_n) \\
&= Y_1 + Y_2 + \cdots + Y_n + E(Y_{n+1} + \cdots + Y_{n+k} \mid S_1, \cdots S_n) \\
&= Y_1 + Y_2 + \cdots + Y_n + E(Y_{n+1} + \cdots + Y_{n+k}) \\
&= S_n
\end{aligned}
\]

Thus, for a fair game, the best estimate of future value \(S_{n+k}\)
given the history up to \(n\), is just \(S_n\).

\section{Discrete time martingle}\label{discrete-time-martingle}

\textbf{Def.} consider probability space \((\Omega, \mathcal{F}, P)\).
\((X_m)_{m=0,1,\cdots}\) be a stochastic process, i.e.~\(X_m\) is R.V.
for each \(m \geq 0\), and integrable \((E(|X|) < \infty)\)

Then \((X_m)_{m=0,1,\cdots}\) is called a Martingle if \[
X_m = E(X_{m+1} \mid X_0, X_1, \cdots X_m), \quad \text{for } \forall m \geq 0
\]

a submartingle if \[
X_m \leq E(X_{m+1} \mid X_0, \cdots X_m), \quad \text{for } \forall m \geq 0
\]

a supmartingle if \[
X_m \geq E(X_{m+1} \mid X_0, \cdots X_m) \quad \text{for } \forall m \geq 0
\]

\textbf{Remark:} A martingle means that if we have the observation of
the history of the process \(X_0, \cdots X_m\), then the observed. value
of present is the expected value of future.

\section{Filtration}\label{filtration}

In many practical example, our observation may not be just \(X\), for
example if we have a sequence \((X_m, Y_m)_{m=0,1,\cdots}\) where
\(X_m, Y_m\) might be correlated, in this case, it is reasonable to
define martingle as \[
X_m = E(X_{m+1} \mid (X_0, Y_0), \cdots (X_m, Y_m)) \quad \text{or} \quad = E(X_{m+1} \mid \mathcal{F}_m)
\] where each \((Y_m, X_m)\) is measurable w.r.t. to \(\mathcal{F}_m\),
and of course \(X_m\)

we thus can redefine martingle under a more generalized language, we we.
introduce the concept of filtration

\textbf{Def:} a set of sub-\(\sigma\)-field \(\{\mathcal{F}_n\}\) on
\((\Omega, \mathcal{F})\) is called a filtration s.t.
\(\mathcal{F}_0 \subseteq \mathcal{F}_1 \subseteq \mathcal{F}_2 \subseteq \cdots \mathcal{F}_n \subseteq \mathcal{F}\),
\(\forall \mathcal{F}_n\) is \(\sigma\)-field

\textbf{Remarks:} A filtration represents any procedure of collecting
more and more information as time goes on.

\textbf{Def:} let \((X_m)_m\) on \((\Omega, \mathcal{F}, P)\) and
\((\mathcal{F}_m)_m\) be a filtration. we called. \((X_m)_m\) is adapted
to \((\mathcal{F}_m)\) if \(\sigma(X_n) \subseteq \mathcal{F}_n\)

\section{Martingle w.r.t.
filtration.}\label{martingle-w.r.t.-filtration.}

consider \((\Omega, \mathcal{F}, P)\), \((X_m)_{m=0,1,\cdots}\) and
filtration \((\mathcal{F}_m)_{m=0,1,\cdots}\) then \((X_m)_m\) is called
a martingle w.r.t. to filtration \((\mathcal{F}_m)_m\) if

\begin{itemize}
\tightlist
\item
  \(\forall m \geq 0\), \(X_m\) is \(\mathcal{F}_m\) measurable.
  (adapted)
\item
  \(X_m = E(X_{m+1} \mid \mathcal{F}_m)\)
\end{itemize}

\textbf{Remarks:} This definition keep the essense of martingle that
when the past information is known (conditional), the present value is
the expected future value.

\textbf{Lemma:} let \((X_m)_m\) a martingle w.r.t. \((\mathcal{F}_m)_m\)
on \((\Omega, \mathcal{F}, P)\) then
\(X_m = E(X_n \mid \mathcal{F}_m)\), \(\forall n > m\).

\textbf{proof:} \[
X_m = E[X_{m+1} \mid \mathcal{F}_m] = E[E(X_{m+2} \mid \mathcal{F}_{m+1}) \mid \mathcal{F}_m]. \quad \text{Since } \mathcal{F}_m \subseteq \mathcal{F}_{m+1}
\] \[
= E[X_{m+2} \mid \mathcal{F}_m]
\] continous this process can show the lemma.

\section{Natural filtration}\label{natural-filtration}

\textbf{prop:} suppose we have discrete time process
\(X_0, X_1, \cdots X_m \cdots\), we can define
\(\mathcal{F}_m^X := \sigma(X_0, X_1, \cdots X_m)\), \(m = 0, 1 \cdots\)
natural filtration of \((X_m)_m\), \((\mathcal{F}_m^X)\) is the smallest
filtration s.t. \((X_m)_m\) is adaptive.

\textbf{Remark:} \(\mathcal{F}_m^X\) is the \(\sigma\)-algebra generated
by previous history \((X_0, X_1, \cdots X_m)\), i.e.~the smallest
`information' of the entire history.

\textbf{proof:} need to show that if \(X_m\) is \(\mathcal{F}_m\)
measurable, then \(\mathcal{F}_m^X \subseteq \mathcal{F}_m\),
\(\forall m\).

suppose we have \(\sigma(X_0) \subseteq \mathcal{F}_0\),
\(\sigma(X_1) \subseteq \mathcal{F}_1\),
\(\cdots \sigma(X_m) \subseteq \mathcal{F}_m\). and
\(\mathcal{F}_0 \subseteq \mathcal{F}_1 \subseteq \mathcal{F}_2 \cdots \subseteq \mathcal{F}_m\).

then we know
\(\sigma(X_0) \subseteq \mathcal{F}_0 \subseteq \mathcal{F}_m\),
\(\sigma(X_1) \subseteq \mathcal{F}_1 \subseteq \mathcal{F}_m\),
\(\cdots \sigma(X_m) \subseteq \mathcal{F}_m\).

i.e.~\(X_0, \cdots X_m\) are all \(\mathcal{F}_m\) measurable.

thus
\(\sigma(X_0 \cdots X_m) = \mathcal{F}_m^X \subseteq \mathcal{F}_m\).

\section{Martingle Transformation.}\label{martingle-transformation.}

The martingle transformation of a discrete martingle can be viewed as a
``discrete'' form of later defined stochastic Integral

Consider a discrete model of stock price \((X_m)_{m=0,1,\cdots 2}\)
adaptive to \((\mathcal{F}_m)_{m=0,1,\cdots}\). let \(U_m := \#\) of
share from \(m-1\) to \(m\). This should be \((\mathcal{F}_{m-1})\)
measurable, as we only know the information till \(m-1\) to make
decision how much to hold till next day. Then my profit would be \[
U_1(X_1 - X_0) + U_2(X_2 - X_0) + \cdots + U_n(X_n - X_{n-1})
\]

The formal definition of martingale transformation is as follows

\textbf{Def} let \((X_m)_{m=0,\cdots}\) on \((\Omega, \mathcal{F}, P)\),
adapted to \((\mathcal{F}_m)_m\), let \((U_m)_{m=1,2,\cdots}\)
measurable respet to \(\mathcal{F}_{m-1}\)

let \[
Z_m = U_1(X_1 - X_0) + U_2(X_2 - X_1) + \cdots + U_m(X_m - X_{m-1})
\] \[
Z_0 = 0
\] \((Z_m)_{m=0,1,2,\cdots}\) is called the martingle transformation of
\(X_m\).

As the name suggest, the martingle transformation transform a martingle
to another martingle

\textbf{Th:} if \((X_m)_m\) is a martingle and \((Z_m)_m\) is
\(\mathcal{F}_{m-1}\) measurable \(\forall m = \{0, 1, 2, \cdots\}\),
then \((Z_m)_m\) is a martingle.

\textbf{Proof:} for a martingle
\(M_m = E(M_{m+1} \mid \mathcal{F}_m) \iff E(M_{m+1} - M_m \mid \mathcal{F}_m) = 0\).
\[
\begin{aligned}
E(Z_{m+1} - Z_m \mid \mathcal{F}_m)
&= E(U_{m+1}(X_{m+1} - X_m) \mid \mathcal{F}_m) \\
&= U_{m+1} E(X_{m+1} - X_m \mid \mathcal{F}_m) = 0. \qquad \blacksquare
\end{aligned}
\]

\chapter{Stopping Time}\label{stopping-time}

One advantage of a martingle process is that it provide. properties in
the stopping time of the sequence. First we need to know what is a
stopping time and its definition

\section{Stopping time}\label{stopping-time-1}

Consider the stock model example again. Suppose we buy one share, at
\(t=0\) and. hold tile random time \(T\). Then \[
U_m \equiv \mathbf{1}_{\{m \leq T\}} = 1 - \mathbf{1}_{\{T \leq m-1\}} \qquad \text{and} \qquad Z_m = X_{\min(m,T)} - X_0
\]

\textbf{obs:} \[
U_m \text{ is } \mathcal{F}_{m-1} \text{ measurable} \iff \{T \leq m-1\} \text{ is } \mathcal{F}_{m-1} \text{ measurable}
\] \[
\iff \{T \leq m\} \text{ is } \mathcal{F}_m \text{ measurable}
\]

\textbf{Def} A R.V. \(T : \Omega \to \{0, 1, \cdots 2, \cdots \infty\}\)
is called a stopping time. for filtration \(\{\mathcal{F}_m\}\) if
\(\{\omega \mid T(\omega) \leq n\}\) is \(\mathcal{F}_n\) measurable
for. each \(n \geq 0\).

\textbf{Remark:} Essentialy, a stopping time is to make a decision that
a process stop at time \(n\) or not, based on the information only
avaible so far (not future). The Def above is exactly this in the formal
probability language (\(\{T \leq n\}\) is the to decide stop or not;
\(\mathcal{F}_n\) measuribility make sure the information).

\section{Stopped process}\label{stopped-process}

\textbf{Def} let \((X_m)_m\) be a stochastic process and \(T\) be a
stopping time, Then a stopped process at \(T\) \(X_{\min(n,T)}\) or
\(X_{n \wedge T}\) is \[
X_{\min(n,T)(\omega)}(\omega) = \begin{cases} X_n(\omega), & n \leq T(\omega) \\ X_{T(\omega)}(\omega), & n > T(\omega) \end{cases} \qquad \omega \in \Omega.
\]

\textbf{prob:} let \(T, S\) be stopping time. on
\((\Omega, \mathcal{F}, P)\) with filtration
\((\mathcal{F}_n)_{n=0,\cdots}\), then \(T \wedge S = \min(T, S)\) and
\(T \vee S = \max(T, S)\) are also stopping time

\section{Stopped martingle}\label{stopped-martingle}

For a martingle process, a nice result is that a stopped. martingle (or
sub, sup) is still a martingle (sub, sub)

\textbf{Th} let \((X_n, \mathcal{F}_n)\) be a martingle (sub, sup) on
\((\Omega, \mathcal{F}, P)\), and \(T\) is a stopping time for
\((\mathcal{F}_n)_n\), then \((X_{n \wedge T}, \mathcal{F}_n)\) is also
a martingle (Sub, sup).

\section{Doob's optional stopping}\label{doobs-optional-stopping}

The martingle property allow ones to find the expectation of stoped
value. Conviniently using the following th-

\textbf{Th (Doob's optional stopping)} let \((X_n, \mathcal{F}_n)\) be a
submartingle and \(T < \infty\) a.s. is a stopping time (will stop) s.t.
\((X_{n \wedge T})_n\) is uniformly integrable. then \[
E(X_T) \geq E(X_0).
\] In perticular for a martingle \(E(X_T) = E(X_0)\)

The Doob's optional sampling th is useful for finding the ``probability
to ruin'' in random walk problem.

\textbf{Ex:} Consider a simple symmetric random walk \((\xi_n)_n\) on an
interval \((-a, b)\), starting at \(0\). \(T_{ab}\) be the stopping time
hit \(-a\) or \(b\). \(S_n = \sum_k \xi_k\) is a martingle, thus \[
E(S_{T_{ab}}) = E(S_0) = 0 = -a \, P(T_{ab} = -a) + b \, P(T_{ab} = b) = 0
\] \[
P(T_{ab} = -a) + P(T_{ab} = b) = 1 \implies P(S_{T_{ab}} = -a) = \frac{b}{b+a}
\]

the probability of hits \(-a\) before hit \(b\).

\textbf{Remark:} should make sure the Th condition is statisfied. If in
above example the walk is bound \((-\infty, b)\), then
\((S_{n \wedge T_{ab}})\) is not integrable and.
\(E(S_{T_b}) \neq E(S_0)\)

\chapter{Brownian Motion}\label{brownian-motion}

The stock market can be model as simple as predicting price go up or
down every day, and extend to Continous diffusion path in BS model.
Brownian mothion is fundeful of stochasti Calculons (Ito Integral), and
can be view

\section{Simple Random Walk}\label{simple-random-walk}

Idea of Simple random walk is `Simple'. You go left with prob \(p\),
other wise you go right. with prob \(1-P\). In the following we will be
interested in the symmetric case where. \(p = \tfrac12\).

If each step is denote as \(\xi_i\), then
\(P(\xi_i = 1) = \tfrac12 = P(\xi_i = -1) = 1\).

and \(X_m = \sum_{k=1}^{m} \xi_k\), with \(X_0 = 0\)

\textbf{property of symmetric Random Walk}

\begin{enumerate}
\def\labelenumi{\arabic{enumi}.}
\item
  Independent Increment.
\item
  \(E(M_{k_{i+1}} - M_{k_i}) = 0\),
  \(\operatorname{Var}(M_{k_m} - M_{k_i}) = k_m - k_i\) \[
  0 = k_0 \leq k_1 \leq \cdots \leq k_m.
  \]
\item
  Sym Random Walk is a martingle.
\item
  \([M, M]_k = \sum_{j=1}^{k} (M_j - M_{j-1})^2 = k\) (Quadratic
  Variation)
\end{enumerate}

\textbf{Remark:} The QV is irelavent to prob of step, thus is general to
more symmetric path.

In practice, QV of a path is easy to evaluate is tick-by-tick price
data.

\section{Scaled Random Walk}\label{scaled-random-walk}

The random walk defined above has time step of \(1\), that means the.
time can only be \(1, 2, 3, \cdots\). we can make the time to be
\(\prec \tfrac12\) by rescale it as \[
W^n_t = \frac{1}{\sqrt{n}} M_{nt}, \qquad \text{where } n \in \mathbb{N}^+ \text{ and } nt \text{ is integer}.
\]

essentially we devide the time \([0, 1]\) to \(n\) interval of size
\(\tfrac1n\).

when \(nt\) is not integer, the \$W\^{}n\_t = \$ interpolated value,
between closest integer value.

the \(W^n_t\) share the. similar properties as random walk.

\begin{itemize}
\tightlist
\item
  independed increment
\item
  \(E(W^n(t) - W^n(t_0)) = 0\)
  \(\qquad \operatorname{Var}(W^n_t - W^n_s) = t - s\)
\item
  martingle \(E(W^n_t \mid \mathcal{F}_s) = W^n_s\), \(\quad s \leq t\)
  \& \(ns, nt \in \mathbb{N}^+\)
\item
  Quadratic Variation \([W^n, W^n](t) = t\)
\end{itemize}

\textbf{limiting dist of scaled random walk.}

if we chose large \(n\) and draw histogram of \(W^n_t\), we notice that
the histogram will approximate normal dist. In fact, the \(W^n_t\) will
converge to normal. as \(n \to \infty\)

\textbf{Th.} fix \(t \geq 0\), As \(n \to \infty\), \(W^n(t)\) evaluated
at time \(t\) Converge to \(N(0, t)\).

\textbf{proof:} the proof can be show the characteristic / moment
generating function Converge.

we thus has the first definition of BM.

\section{Brownian Motion (definition)}\label{brownian-motion-definition}

\textbf{Def} \((\Omega, \mathcal{F}, P)\), \((B_t)_{t \geq 0}\) is a
stochastic process, i.e.~for \(\forall t \geq 0\), \(B_t\) is a. R.V.

\((B_t)_{t \geq 0}\) is called BM if \(B_0 = 0\) and
\(\forall\, 0 = t_0 < t_1, \cdots < t_i\) \[
B_{t_1} - B_{t_0}, \ B_{t_2} - B_{t_1}, \ \cdots B_{t_i} - B_{t_{i-1}} \quad \text{all independent}, \quad \forall s < t, \ B_t - B_s \overset{(d)}{=} N(0, t-s)
\]

\textbf{Def2} (using filtration), \((B_t)_{t \geq 0}\) is called BM
w.r.t. \((\mathcal{F}_t)\) if

\begin{itemize}
\tightlist
\item
  \(B_t\) is continous path, \(B_0 = 0\)
\item
  \(\forall t\), \(B_t\) is measurable w.r.t. \(\mathcal{F}_t\)
  \(\rightsquigarrow \mathcal{F}_t^B \subseteq \mathcal{F}_t\)
\item
  \(B_t - B_s \perp \mathcal{F}_s\) and
  \(B_t - B_s \overset{(d)}{=} N(0, t-s)\)
\end{itemize}

from the definition we can show that the
\(\operatorname{Cov}(B_t, B_s) = E(B_t B_s) = s \wedge t\).

\[
\begin{aligned}
\operatorname{Cov}(B_t, B_s) &= E(B_t B_s) - E(B_t)E(B_s) = E(B_t B_s) \\
&= E(B_s(B_t - B_s) + B_s^2) \qquad \text{if } t \geq s \geq 0 \\
&= E(B_s(B_t - B_s)) + E(B_s^2) = s = s \wedge t.
\end{aligned}
\]

Thus the BM can also be defined from Gaussian process.

\section{Gaussian process}\label{gaussian-process}

\textbf{Def Gaussian process}

\((X_t)_{t \geq 0}\) SP on \((\Omega, \mathcal{F}, P)\) is called a
Gaussian process if \((X_{t_1}, X_{t_2}, \cdots X_{t_d})\) form a multi
dimensional Gaussian, \(\forall t_1 \leq t_2 \leq \cdots \leq t_d\).

Then the BM motion is a Gaussian process with \(B_0 = 0\),
\(E[B_t] = 0\) and \(\operatorname{Cov}(B_s, B_t) = s \wedge t\)
\(\ \forall s, t\).

\textbf{Def} \((X_1, \cdots X_n) \in \mathbb{R}^d\) is a
multi-dimensional Gaussian with mean \(\mu \in \mathbb{R}^d\), and cov
matrix \(\Sigma \in \mathbb{R}^{d \times d}\), symmetric, semi def if \[
P(x_1, \cdots x_n) = P(\vec{x}) = \frac{1}{\sqrt{(2\pi)^d \det(\Sigma)}} \, e^{-\frac12 \langle \Sigma^{-1}(\vec{x} - m), (\vec{x} - m)\rangle}
\] or can write further \[
E[e^{i \langle \vec{a}, x\rangle}] = e^{\, i \langle \vec{a}, \mu\rangle - \frac12 \langle \Sigma \vec{a}, \vec{a}\rangle}.
\]

\section{Properties}\label{properties}

\textbf{(1)} \((B_t)_{t \geq 0}\) is a martingle w.r.t. to
\(\mathcal{F}_t\), \[
E(B_t \mid \mathcal{F}_s) = E(B_s + B_t - B_s \mid \mathcal{F}_s) = B_s + E(B_t - B_s \mid \mathcal{F}_s) = B_s + E(B_t - B_s) = B_s
\]

\textbf{(2)} \(B_t^2 - t\) is a martingle w.r.t.
\((\mathcal{F}_t)_{t \geq 0}\) \[
\begin{aligned}
E(B_t^2 - t \mid \mathcal{F}_s) &= E\big((B_s + (B_t - B_s))^2 - t \mid \mathcal{F}_s\big) \\
&= E\big(B_s^2 + 2B_s(B_t - B_s) + (B_t - B_s)^2 \mid \mathcal{F}_s\big) - t \\
&= B_s^2 + 2B_s E(B_t - B_s \mid \mathcal{F}_s) + E\big((B_t - B_s)^2 \mid \mathcal{F}_s\big) - t \\
&= B_s^2 - s
\end{aligned}
\]

\section{Quadratic variation of BM}\label{quadratic-variation-of-bm}

we denote the QV of BM \((B_t)_t\) as on a partition
\(0 = t_0 < t_1 < t_2 \cdots < t_n = T\) as \[
[B, B]_{t_n} := (B_{t_1} - B_{t_0})^2 + (B_{t_2} - B_{t_1})^2 + \cdots + (B_{t_n} - B_{t_{n-1}})^2
\]

\textbf{obs:} denote \[
\langle B, B\rangle_{t_n} := E[(B_{t_1} - B_{t_0})^2 \mid \mathcal{F}_{t_0}] + E[(B_{t_2} - B_{t_1})^2 \mid \mathcal{F}_{t_1}] + \cdots + E[(B_{t_n} - B_{t_{n-1}})^2 \mid \mathcal{F}_{t_n}]
\]

Since \(B_{t_n} - B_{t_{n-1}} \sim N(0, t_n - t_{n-1})\) and are
independent from previous \[
\langle B, B\rangle_{t_n} = (t_1 - t_0) + (t_2 - t_1) + \cdots + (t_n - t_{n-1}) = t_n - t_0 = t_n.
\]

One important property for the QV of BM is that \([B, B]_{t_n} \to t_n\)
when the partition is tend to infinitesmal small.

\textbf{Th:} let \(\|\sigma\| := \max_{i=1 \cdots n} |t_i - t_{i-1}|\),
then we have. \[
\lim_{\|\sigma\| \to 0} [B, B]_{t_n} = \langle B, B\rangle_{t_n} = t_n. \qquad \text{or more precisely.}
\] \[
P\big(\{\omega \mid [B, B]_{t_n}(\omega) = t_n\}\big) = 1.
\]

\textbf{proof:} we first consider \[
E\big([B,B]_{t_n} - t_n\big) = E\Big(\textstyle\sum_i (B_i - B_{i-1})^2 - \sum_i (t_i - t_{i-1})\Big) = \sum_i E\big((B_i - B_{i-1})^2 - (t_i - t_{i-1})\big) = 0.
\]

next consider the variation \[
\begin{aligned}
E\big(([B,B]_{t_n} - t_n)^2\big) &= E\Big(\big[\textstyle\sum_i ((B_i - B_{i-1})^2 - (t_i - t_{i-1}))\big]^2\Big) \\
&= E\Big[\textstyle\sum_i ((B_i - B_{i-1})^2 - (t_i - t_{i-1}))^2 + 2\sum_{i \neq j} ((B_i - B_{i-1})^2 - (t_i - t_{i-1}))((B_j - B_{j-1})^2 - (t_j - t_{j-1}))\Big] \\
&= \textstyle\sum_i E\big[((B_i - B_{i-1})^2 - (t_i - t_{i-1}))^2\big] + 2\sum_{i \neq j} E\big[((B_i - B_{i-1})^2 - (t_i - t_{i-1}))((B_j - B_{j-1})^2 - (t_j - t_{j-1}))\big]
\end{aligned}
\]

from independent intervals, we have.
\(B_i - B_{i-1} \perp B_j - B_{j-1}\), \(i \neq j\) \[
\Rightarrow E\big[((B_i - B_{i-1})^2 - (t_i - t_{i-1}))((B_j - B_{j-1})^2 - (t_j - t_{j-1}))\big] = E[\cdots]\,E[\cdots] = 0.
\]

denote \(B_i - B_{i-1} = N_i\),
\(\operatorname{Var}(N_i) = t_i - t_{i-1}\) \[
\begin{aligned}
\Rightarrow E[\ ] &= \textstyle\sum_i E\big((N_i^2 - \operatorname{Var}(N_i))^2\big) = \sum_i E\big((N_i^2 - E(N_i^2))^2\big) \\
&= \textstyle\sum_i E\big(N_i^4 - 2N_i^2 E(N_i^2) + E^2(N_i^2)\big) \\
&= \textstyle\sum_i 3(t_i - t_{i-1})^2 - 2(t_i - t_{i-1})^2 + (t_i - t_{i-1})^2 \\
&= 2\textstyle\sum_i (t_i - t_{i-1})^2 \leq 2\|\sigma\| \sum_i |t_i - t_{i-1}| \leq 2|\sigma| \, t_n.
\end{aligned}
\]

thus
\(\lim_{\sigma \to 0} E\big[([B, B] - \langle B, B\rangle)^2\big] = 0\)

\textbf{Remark:} It is Convinient to memorize the \([B, B]_t = t\) in
``differential'' form as \[
(dB)^2 = dt
\] The simply means BM cumulate QV at rate one per unit time.

\chapter{Stochastic Integral}\label{stochastic-integral}

suppose we have a stock price \((X_t)_t\) and a trading strategy
\((a_t)_t\), which essentially a Random variable that tells me how much
share to hold. The How do I Calculate my total profit from \(0, to\ T\)

one simple model turn descret time into interval
\(t_0 < t_1 < t_2 \cdots < t_n\), and we hold \(\xi_i\) \# of share
between \([t_{i-1}, t_i]\), decide at \(t_{i-1}\) (no future
information) \(\Rightarrow \xi_i\) is \(\mathcal{F}_{t_{i-1}}\)
measurable. and BM \((B_t)_t\) as price movement

Then the profit is
\(I_t = \xi_1(B_{t_1} - B_{t_0}) + \cdots + \xi_n(B_t - B_{t_{n-1}})\)

when \(\sigma t = t_i - t_{i-1}\) is very small, this looks like an
integral in regular calculous sense. However due to the QV of BM, the
exact stochastic Calculous is more Complicated the intuitive argument
above. but the general idea is similar.

we start with define the stochastic integral (Ito integral specificaly)
for Simple step function \[
0 = t_0 < t_1 < t_2 \cdots < t_n = T. \quad \text{a partition}
\]

\section{Step process}\label{step-process}

\textbf{Def step process} \(G_t\): \[
(G_t)_{t \in [0,T]} := \sum_{k=1}^{n} \xi_{k-1} \mathbf{1}_{\{t_{k-1} < t \leq t_k\}}, \qquad \xi_{k-1} \text{ is } \mathcal{F}(t_{k-1}) \text{ measurable.}
\] \[
\mathbf{1}_{\{t_{k-1} < t \leq t_k\}} = \begin{cases} 1 & t_{k-1} < t < t_k \\ 0 & \text{other} \end{cases}
\]

\section{Itô's stochastic integral (on step
processes)}\label{ituxf4s-stochastic-integral-on-step-processes}

\textbf{Def Itô's stochastic integral} (on step processes) \[
I_t^n = \int_0^t G_u^n \, dB_u = \sum_{i=0}^{k-1} \xi_i(B_{t_{i+1}} - B_{t_i}) + \xi_k(B_t - B_{t_k}), \quad t \in [t_k, t_{k+1}], \quad (B_t)_{t \geq 0} \text{ is BM} \iff dI_t^n = G_t^n \, dB_t
\]

\textbf{Remark:} An intuitive under standy of this definition is to
think \(t_0, t_1 \cdots t_n\) as the trading days, \(B(t_i)\) be the
price at these days and \(\xi_i\) as number of shares acquired at date
\(t_i\), hold tile time \(t_{i+1}\). Then the total gain of asset (which
would stans for `integral') would be
\(\sum(\text{share} \times \text{price})\) all trades.

\section{Properties of the Itô
integral}\label{properties-of-the-ituxf4-integral}

\textbf{Th} \(I_t^n\) is a continous martingle with \(I_0^n = 0\)

\textbf{proof:} (1) let \(0 \leq s < t\), to show
\(I_s^n = E(I_t^n \mid \mathcal{F}_s)\), we need to show
\(E(I_t^n - I_s^n \mid \mathcal{F}_s) = 0\). suppose
\(s \in [t_l, t_{l+1}]\), \(t \in [t_k, t_{k+1}]\), \(l \leq k\)

then \[
\begin{aligned}
I_t^n - I_s^n &= \xi_k(B_t - B_{t_k}) + \xi_{k-1}(B_{t_k} - B_{t_{k-1}}) \cdots + \xi_l(B_{t_{l+1}} - B_{t_l}) - \xi_l(B_s - B_{t_l}) \\
&= \xi_k(B_t - B_{t_k}) + \cdots + \xi_l(B_{t_{l+1}} - B_s)
\end{aligned}
\] \[
E\big(\xi_k(B_t - B_{t_k}) \mid \mathcal{F}_s\big) = E\big(E(\xi_k(B_t - B_{t_k}) \mid \mathcal{F}_{t_k}) \mid \mathcal{F}_s\big) = E\big(\xi_k E(B_t - B_{t_k} \mid \mathcal{F}_{t_k}) \mid \mathcal{F}_s\big) = 0
\] \[
E\big(\xi_l(B_{t_{l+1}} - B_s) \mid \mathcal{F}_s\big) = \xi_l E(B_{t_{l+1}} - B_s \mid \mathcal{F}_s) = 0.
\] \(\to I_t^n\) is continous martingle.

\textbf{Th Ito isometry} \[
E(I_t^{n2}) = E\Big[\int_0^t G_u^{n2} \, du\Big]
\]

\textbf{proof:} Ito isomety
\(E(I^2) = E\big((\sum_i \xi_i \Delta B_i)^2\big)\) \[
= E\Big(\sum_i \xi_i^2 \Delta B_i^2 + 2\sum_{i \neq j} \xi_i \Delta B_i \, \xi_j \Delta B_j\Big)
\] when \(i \neq j\), assume \(i < j\), \(\Delta B_j\) is independent of
\(\Delta B_i, \xi_i, \xi_j\) thus
\(E(\xi_i \xi_j \Delta B_i \Delta B_j) = 0\) \[
\begin{aligned}
E(I^2) &= E\Big(\sum_i \xi_i^2 \Delta B_i^2\Big) = \sum_i E(\xi_i^2 \Delta B_i^2) \\
&= \sum_i E(\xi_i^2) E(\Delta B_i^2) \\
&= E\big(\xi_i^2 (t_{i+1} - t_i)\big) = E\Big(\int_0^t G_s^{n2} \, ds\Big)
\end{aligned}
\]

\textbf{Th The QV of \(I_t^n\)} \[
[I_t^n, I_t^n] = \int_0^t G_u^{n2} \, du, \quad \text{and.} \qquad \langle I_t^n, I_t^n\rangle = \int_0^t G_u^{n2} \, du
\]

\textbf{proof:} Calculate the QV on each piece wise interval. \[
\begin{aligned}
[I_{t_i}^n, I_{t_i}^n] &= \lim_{|\sigma| \to 0} \sum_{j=1}^{m} (I_{h_j} - I_{h_{j-1}})^2, \quad \text{where } t_{i-1} = h_0 < h_1 \cdots < h_m = t_i. \\
&= \lim_{|\sigma| \to 0} \sum_{j=1}^{m} \xi_{t_{i-1}}(B_{h_j} - B_{h_{j-1}})^2 = \xi_{t_{i-1}} \lim_{|\sigma| \to 0} \sum_{j=1}^{m} (B_{h_j} - B_{h_{j-1}})^2 \\
&= \xi_{t_{i-1}}(t_i - t_{i-1})
\end{aligned}
\] notice that the QV on \(0, t\) is then sum of all this \[
\text{thus} \quad [I_t^n, I_t^n] = \sum_j \xi_{t_{j-1}}(t_j - t_{j-1}) = \int_0^t G_u^{n} \, du.
\] \[
\begin{aligned}
\langle I_{t_i}^n, I_{t_i}^n\rangle &= \sum_{c=1}^{k} E\big[\xi_{i-1}^2(B_{t_i} - B_{t_{i-1}})^2 \mid \mathcal{F}_{t_{i-1}}\big] = \sum_{i=1}^{k} \xi_{i-1}^2 E\big[(B_{t_i} - B_{t_{i-1}})^2 \mid \mathcal{F}_{t_{i-1}}\big] \\
&= \sum_{i=1}^{k} \xi_{i-1}^2 (t_i - t_{i-1})^2 = \int_0^t G_t^2 \, dt.
\end{aligned}
\]

\textbf{Remark:} We can denote the integral form
\(I_t^n = \int_0^t G_u^n \, dB_u\) in ``differential form''
\(dI_t^n = G_t^n \, dt\) and the QV \[
(dI_t^n)^2 = (G_t^n \, dB_u)^2 = G_t^{n2}(dB_u)^2 = G_t^{n2} \, dt.
\] or. \(I_t^n\) accumulate \(G_t^{n2}\) QV per unit time.

\textbf{Th Martingle Property} \[
I_t^{n} - \int_0^t G_u^{n2} \, du \quad \text{is a martingle.}
\]

\textbf{Remark:} the QV is computed path-by-path. chosing different
\(a_t\) would leads to different integral. QV can be regarded. as a
measure of risk, and it depends on the size of position \((a_t)\) we
take.

\section{From simple step process to General
process}\label{from-simple-step-process-to-general-process}

use \((G_t^n)_{t \in [0,T]}\), a step process to approximate any \(G\)
to defines \(\int_0^T G_t \, dB_t\)

let \(T\) be fixed, \[
L^2(0,T) = \Big\{(G_t)_{0 \leq t \leq T}, \ G_t \text{ measurable w.r.t. } (\mathcal{F}_t), 0 \leq t \leq T, \text{ and } E\Big[\int_0^T |G_t|^2 dt\Big] < \infty\Big\}
\] \[
\exists (G^n)_t \in L^2(0,T), \text{ bounded step process (sequence of) s.t. } E\Big[\int_0^T |G_u - G_u^n|^2 du\Big] \to 0.
\]

\textbf{Def} let \(G_t \in L^2(0,T)\) and \((G^n)_t\) as bound sequence
of step process above then
\(L_t = \lim_{n \to \infty} \int_0^t G_u^n \, dB_u\) exist and defined
as \(\int_0^t G_u \, du\).

\textbf{Properties:}

\begin{itemize}
\tightlist
\item
  \(L_t\) is continous martingle, \(L_0 = 0\)
\item
  \(E(L_t^2) = E\big(\int_0^t G_u^2 \, du\big)\) --- Ito isomety
\item
  \(\langle L_t, L_t\rangle = \int_0^t G_u^2 \, du = [L_t, L_t] := \lim_{\sigma \to 0} \sum_{k=0} (L_{t_k} - L_{t_{k-1}})^2\)
  --- QV
\item
  \((L_t)^2 - \int_0^t G_u^2 \, du\) is a martingle
\item
  \(E(L_t) = 0\).
\end{itemize}

\textbf{Justification on the QV.} \[
E\Big[\big(\textstyle\sum_k (L_{t_k} - L_{t_{k-1}})^2 - \langle L, L\rangle_t\big)^2\Big] = E\Big[\big(\textstyle\sum_k \underbrace{((L_{t_k} - L_{t_{k-1}})^2 - (\langle L,L\rangle_{t_k} - \langle L,L\rangle_{t_{k-1}}))}_{A_k}\big)^2\Big]
\] observation
\(E[(L_t - L_s)^2 \mid \mathcal{F}_s] = E[L_t^2 - L_s^2 \mid \mathcal{F}_s]\),
let \(s = t_{k-1}\), \(t = t_k\) \[
E\big[(L_{t_k} - L_{t_{k-1}})^2 \mid \mathcal{F}_{t_{k-1}}\big] = E\big[L_{t_k}^2 - L_{t_{k-1}}^2 \mid \mathcal{F}_{t_{k-1}}\big] = E\big[\langle L,L\rangle_{t_k} - \langle L,L\rangle_{t_{k-1}} \mid \mathcal{F}_{t_{k-1}}\big]
\] as
\(E[(L_{t_k} - L_{t_{k-1}})^2 \mid \mathcal{F}_{t_{k-1}}] = \langle L,L\rangle_{t_k} - \langle L,L\rangle_{t_{k-1}} = E[\langle L,L\rangle_{t_k} - \langle L,L\rangle_{t_{k-1}} \mid \mathcal{F}_{t_{k-1}}]\)
\[
\Rightarrow E(A_k \mid \mathcal{F}_{t_{k-1}}) = 0
\] \[
E\big((\textstyle\sum_k A_k)^2\big) = E\big(\textstyle\sum_k A_k^2\big) + 2E\big(\textstyle\sum_{i \leq k} A_i A_k\big) \quad \text{since } E(A_i A_k \mid \mathcal{F}_{k-1}) = A_i E(A_k \mid \mathcal{F}_{k-1}) = 0 \text{ and } E(A_k^2) \text{ close to } 0.
\]

\textbf{Remark:} again we denote \(dL_t = G_t \, dB_t\) and
\((dL_t)^2 = G_t^2 \, dt\).

\section{Itô processes and Itô's chain
rule}\label{ituxf4-processes-and-ituxf4s-chain-rule}

we then extend the integral process defined above to more generalise Ito
processes.

\textbf{Def - Itô processess} let \((B_t)_{t \geq 0}\) a BM on
\((\Omega, \mathcal{F}, P)\) with filtration \(\mathcal{F}_t\) and some
\((\mathcal{F}_t)_{t \geq 0}\), \((a_t)_{t \geq 0}\) sp adapted to
\(\mathcal{F}_t\) and \(F_t \in L^1(0,T)\), \(G_t \in L^2(0,T)\) and \[
X_t = X_0 + \int_0^t F_u \, du + \int_0^t G_u \, dB_u
\] is called the Ito process, which often denote in differential form \[
dX_t = F_t \, dt + G_t \, dB_t.
\]

\textbf{Remark:} 1. the differential form is a notation for the integral
expression, \(dX_t, dt, dB_t\) has no stand alone meaning

\begin{enumerate}
\def\labelenumi{\arabic{enumi}.}
\setcounter{enumi}{1}
\item
  Almost all sp, exlept the one with jumps (Levy-processes) are Ito
  process.
\item
  Tthe \(F_t \, dt\) are refered as drift term (over all motion. due to
  external field) and Brownian term is called diffusive term.
\end{enumerate}

\textbf{Lemma:} \([X, X]_t = \int_0^t G_u^2 \, dt\). \(\to\) QV come
from BM, drift has no QV.

The natural Question next is to ask if there is a ``chain'' rule. in
stochastic Integral that tell us \(d(U(X_t))\) is. The such rule in
stochastic Calculous is called Ito's chain rule. Similar to regular
chain rule with extra term as result of none zero QV of BM.

\textbf{Th - Itô's Chain Rule.} Suppose
\(dX_t = F_t \, dt + G_t \, dB_t\) is an Ito processes defined above,
let \(U: \mathbb{R} \times [0,T] \to \mathbb{R}\) and \(U(x,t)\) is
continous and
\(\frac{\partial U}{\partial t}, \frac{\partial U}{\partial x}, \frac{\partial^2 U}{\partial x^2}\)
exists and continous, then \(Y_t = U(X_t, t)\) has diff eventtutw \[
\begin{aligned}
dY_t &= U_t \, dt + U_x \, dX_t + \tfrac12 U_{xx}(dX_t)^2 \\
&= U_t \, dt + U_x(F_t \, dt + G_t \, dB_t) + \tfrac12 U_{xx} G_t^2 \, dt.
\end{aligned}
\]

Similar question can be asked, if there is a product rules in stochastic
sense, and the answer is Yes, with also additional term due to QV of BM

\textbf{Example:} \(d(B_u^2) = 2 B_u \, dB_u + dt.\) \[
\begin{aligned}
\int_0^t d(B_u^2) &= \lim_{|\sigma| \to 0} \sum_i (B_{t_i}^2 - B_{t_{i-1}}^2) = \lim_{|\sigma| \to 0}\Big(\textstyle\sum_i (B_{t_i} - B_{t_{i-1}})^2 + 2B_{t_{i-1}}\sum_i (B_{t_i} - B_{t_{i-1}})\Big) \\
&= \lim_{|\sigma| \to 0} \sum_i (B_{t_i} - B_{t_{i-1}})^2 + 2\lim_{|\sigma| \to 0} \sum_i (B_{t_i} - B_{t_{i-1}}) B_{t_{i-1}} \\
&= [B, B]_t + 2\int_0^t B_u \, dB_u = \int_0^t dt + \int_0^t 2 B_u \, dB_u
\end{aligned}
\] \[
\iff dB_u^2 = 2 B_u \, dB_u + dt
\]

\textbf{Example:} \(d(t B_t)\) \[
\int_0^t d(u B_u) = \lim \sum_i (t_i B_{t_i} - t_{i-1} B_{t_{i-1}}) = \lim \sum_i t_i(B_{t_i} - B_{t_{i-1}}) + \lim \sum_i (t_i - t_{i-1}) B_{t_{i-1}} = \int_0^t u \, dB_u + \int_0^t B_u \, du
\] \[
\iff d(t B_t) = t \, dB_t + B_t \, dt
\]

\textbf{Th - Itô's product rule.} let
\(dX_t^1 = F_t^1 \, dt + G_t^1 \, dB_t\) be two Ito process on
\((\Omega, \mathcal{F}, P)\), \(dX_t^2 = F_t^2 \, dt + G_t^2 \, dB_t\)
\[
\begin{aligned}
\text{then} \quad d(X_1 X_2) &= X_2 \, dX_1 + X_1 \, dX_2 + dX_1 \, dX_2 \\
&= X_2 \, dX_1 + X_1 \, dX_2 + G_t^1 G_t^2 \, dt.
\end{aligned}
\]

\textbf{lemma:} \(dX_1 \, dX_2 = G_t^1 G_t^2 \, dt\)

\textbf{proof} let \(G_t^1, G_t^2, F_t^1, F_t^2\) has its Converging
step sequence \(G^{1n}, G^{2n}, F^{1n}, F^{2n}\) and let
\(X_n^1, X_n^2\) defined as. \[
X_n^1 = \sum_i F_{t_{i-1}}^1(t_i - t_{i-1}) + \sum_i G_{t_{i-1}}^1(B_i - B_{i-1})
\] \[
X_n^2 = \sum_i F_{t_{i-1}}^2(t_i - t_{i-1}) + \sum_i G_{t_{i-1}}^2(B_i - B_{i-1})
\] \[
\text{then} \quad \sum_i (X_{t_i}^1 - X_{t_{i-1}}^1)(X_{t_i}^2 - X_{t_{i-1}}^2) = \sum_i \big(F_{t_{i-1}}^1(t_i - t_{i-1}) + G_{t_{i-1}}^1(B_i - B_{i-1})\big)\big(F_{t_{i-1}}^2(t_i - t_{i-1}) + G_{t_{i-1}}^2(B_i - B_{i-1})\big)
\]

taking \(|\sigma| \to 0\), the two drift-drift and drift-diffusion sums
vanish (in \(L^2\) sense): \[
\lim_{|\sigma| \to 0} \sum_i F^1 F^2 (\Delta_i t)^2 = 0, \qquad \lim_{|\sigma| \to 0} \sum_i F^1 G^2 \, \Delta_i B \, \Delta_i t = 0
\] \[
\lim_{|\sigma| \to 0} \sum_i (X_i^1 - X_{i-1}^1)(X_i^2 - X_{i-1}^2) = \lim_{|\sigma| \to 0} \sum_i G^1 G^2 \, \Delta_i B^2 = \int_0^t G_u^1 G_u^2 \, du
\]

\textbf{proof:} \[
\begin{aligned}
\int_0^t d(X^1 X^2) &= \lim_{|\sigma| \to 0} \sum_i (X_{t_i}^1 X_{t_i}^2 - X_{t_{i-1}}^1 X_{t_{i-1}}^2) \\
&= \lim_{|\sigma| \to 0} \sum_i \big[X_{t_{i-1}}^1(X_{t_i}^2 - X_{t_{i-1}}^2) + X_{t_{i-1}}^2(X_{t_i}^1 - X_{t_{i-1}}^1) + (X_{t_i}^1 - X_{t_{i-1}}^1)(X_{t_i}^2 - X_{t_{i-1}}^2)\big] \\
&= \int_0^t X_u^1 \, dX_u^2 + \int_0^t X_u^2 \, dX_u^1 + \int_0^t G_u^1 G_u^2 \, du
\end{aligned}
\]

\textbf{Def} \(dX_t = F_t \, dt + G_t \, dB_t\), and \((a_t)_t\) be a
adapted process \[
\int_0^t a_u \, dX_u := \int_0^t a_u F_u \, du + \int_0^t a_u G_u \, dB_u.
\] or. \(a_t \, dX_t = a_t F_t \, dt + a_t G_t \, dB_t.\)

\textbf{proof - Itô's Chain Rule.} using Itô's product, an proof
chainrule

① using product rule can show for \(U(x) = x^m\), we ①,\(\cdots\) \[
dX_t^m = m X^{m-1} \, dX_t + \tfrac12 m(m-1) X^{m-2} G_t^2 \, dt.
\] (using induction), which satisfy the chain rule. The linearity of
integration give that all polynomial satisfied.

② then we show for \(f(x), g(t)\) that are both polynomial, \[
\begin{aligned}
d(U(X_t)) = d[f(x) g] &= f(x) \, dg + g(t) \, df(x) \\
&= f(x) g' \, dt + g(t)\big(f'(x) \, dx + \tfrac12 f''(x) G_t^2 \, dt\big) \\
&= \frac{\partial U}{\partial t} \, dt + \frac{\partial U}{\partial x}(x) \, dx + \tfrac12 \frac{\partial^2 U}{\partial x^2} G_t^2 \, dt.
\end{aligned}
\] due to linearity, \(U(x,t) = \sum_i f_i(x) g_i(t)\) also Satisfies
Ito formula. That is, all polynomial function \(U\) of \(x, t\) satis
fy.

③ For then \(U\) in Ito's chain rule statement,
\(\exists (U^n)_{n=0,1,\cdots}\), 'a sequence of polynomial function
s.t. \(U^n \to U, U_t^n \to U_t, U_x^n \to U_x, U_{xx}^n \to U_{xx}\)
uniformly on compact subset of \(\mathbb{R} \times [0,T]\). \[
\text{sun} \quad \int_0^t dU^n = \int_0^t U_t^n \, du + \int_0^t U_x^n \, dX + \tfrac12 \int_0^t U_{xx}^n G_u^2 \, du \quad \text{a.s.}
\] passing the limit yet chain rule in general.

\textbf{Summary:} The way remember this is to reconize that
\(dB_t \, dB_t = dt.\) instead of \(0\), as in regular calculus
\(\ dB_t \approx \sqrt{dt}\)

\begin{itemize}
\item
  then for product \[
  \begin{aligned}
  d X_1 X_2 &= X_1 \, dX_2 + X_2 \, dX_1 + dX_1 \, dX_2 \\
  &= X_1 \, dX_2 + X_2 \, dX_1 + (F^1 dt + G^1 dB)(F^2 dt + G^2 dB) \\
  &= \cdots + G^1 G^2 \, dt
  \end{aligned}
  \] the result we higher order term, ignored as usual.
\item
  Itô's chain rule then can be view as lowest order Tyler expansion. \[
  f(X_t) = f(0,0) + \frac{\partial f}{\partial x} \, dX_t + \frac{\partial f}{\partial t} \, dt + \tfrac12 \frac{\partial^2 f}{\partial x^2}(dX_t)^2.
  \] as \((dX_t)^2 = G_t^2 \, dt\), and should be keep.
\end{itemize}

\section{General Itô's chain rule}\label{general-ituxf4s-chain-rule}

Itô's chain rule can be similarly extend to multip variable.

let \(dX_t^i = F_t^i \, dt + G_t^i \, dB_t\) on
\((\Omega, \mathcal{F}, P)\) with \(F^i \in L^1(0,T)\),
\(G^i \in L^2(0,T)\) adapted, \(i = 1, \cdots n\), and
\(u: \mathbb{R}^n \times [0,T] \to \mathbb{R}\) is continous, with
continous \(u_t, u_{x_i}, u_{x_i x_j}\) \(\forall t_i\), then \[
d u(x^1, x^2, \cdots, t) = u_t \, dt + \sum_i \frac{\partial u}{\partial x_i} \, dX_t^i + \tfrac12 \sum_{ij} \frac{\partial^2 u}{\partial x_i \partial x_j} G_t^i G_t^j \, dt
\]

\section{Examples}\label{examples}

\textbf{1. Geometric Brownian Motion.} \[
dX_t = G_t \, dB_t + (\alpha_t - \tfrac12 \sigma^2) \, dt, \qquad B_t \text{ is BM w.r.t. } \mathcal{F}_t \text{ filtration and } G_t, \alpha_t \text{ be adapted process.}
\]

and consider an asset price given by \[
S_t = S_0 \exp(X_t)
\] \[
\begin{aligned}
\text{then} \quad dS_t &= d(S_0 \exp(X_t)) = S_0\big(\exp(X_t) \, dX_t + \tfrac12 \exp(X_t) \, dX_t \, dX_t\big) \\
\text{and} \quad dX_t \, dX_t &= G_t^2 (dB_t)^2 = G_t^2 \, dt. \\
\Rightarrow \quad dS_t &= S_0 \exp(X_t)\big(G_t \, dW_t + (\alpha_t - \tfrac12 \sigma^2) \, dt\big) + \tfrac12 S_0 \exp(X_t) G_t^2 \, dt \\
&= \alpha_t S_t \, dt + G_t S_t \, dB_t.
\end{aligned}
\]

the asset price would have instantaneous mean rate of return.
\(\alpha_t\) and volatility \(G(t)\).

This example includes all possible asset price model that is always
positive, has no jumps and driven by a single BM.

\chapter{Stochastic Analysis and Lévy
Process}\label{stochastic-analysis-and-luxe9vy-process}

\section{Characteristic function}\label{characteristic-function}

\textbf{Def - characteristic function} \(\hat{\mu}(z)\) of probability
measure \(\mu\) on \(\mathbb{R}^d\) is \[
\hat{\mu}(z) = \int_{\mathbb{R}^d} e^{i\langle z, x\rangle} \, \mu(dx)
\]

\textbf{Prop:} the characteristic function of the distribution \(P_X\)
of a R.V. \(X\) on \(\mathbb{R}^d\) is \[
\hat{P}_X(z) = \int_{\mathbb{R}^d} e^{i\langle z, x\rangle} \, P_X(dx) = E[e^{i\langle z, X\rangle}]
\]

\textbf{Properties of \(\hat{\mu}(z)\)}

\begin{itemize}
\tightlist
\item
  If \(\mu = \mu_1 \otimes \mu_2\), then
  \(\hat{\mu}(z) = \hat{\mu}_1(z)\hat{\mu}_2(z)\). If \(X_1\) and
  \(X_2\) are independent R.V. on \(\mathbb{R}^d\), then
  \(\hat{P}_{X_1 + X_2}(z) = \hat{P}_{X_1}(z)\hat{P}_{X_2}(z)\)
\item
  If \(X = (X_j)_{j=1,\cdots n}\) be an \(\mathbb{R}^{n \times d}\)
  valued R.V. where \(X_1, \cdots X_n\) are \(\mathbb{R}^d\)-valued R.V.
  Then \(X_1, \cdots X_n\) are independent iff \[
  \hat{P}_X(z) = \hat{P}_{X_1}(z_1)\hat{P}_{X_2}(z_2) \cdots \hat{P}_{X_n}(z_n), \quad z = (z_j)_{j=1,\cdots n}, \ z_j \in \mathbb{R}^d.
  \]
\item
  If \(\mu_n \to \mu\), then \(\hat{\mu}_n(z) \to \hat{\mu}(z)\)
  uniformly on compact.
\end{itemize}

\section{Poisson process}\label{poisson-process}

\textbf{Poisson dist:} \(\lambda > 0\), \[
1 = \sum_{n=0}^{\infty} e^{-\lambda} \frac{\lambda^n}{n!} \implies P(N = n) = e^{-\lambda} \frac{\lambda^n}{n!}, \quad n = 0, \cdots, \quad \lambda \text{ is called intensity}
\]

\textbf{Poisson process} \(\lambda > 0\), Consider
\(\xi_1, \xi_2, \cdots\) iid \(\sim \operatorname{Exp}(\lambda)\),
i.e.~\(P(\xi \leq t) = e^{-\lambda t}\), \(t \geq 0\),
\(E(\xi) = \tfrac1\lambda\).

consider R.V. \(N_t := \#\) of ``arrival'' at/before \(t\). \[
N_t := \sum_{n=1}^{\infty} n \, \mathbf{1}_{\{\xi_1 + \cdots \xi_{n-1} \leq t, \; \xi_1 + \cdots \xi_n > t\}}
\]

\textbf{prob:} \[
\begin{cases}
N_t - N_s \overset{(d)}{=} \operatorname{Poisson}(\lambda(t-s)), & 0 \leq s \leq t. \\
\text{independent increment.}
\end{cases}
\] \((N_t)_{t \geq 0}\) is then called poisson process with intensity
\(\lambda\)

\section{Compound Poisson process}\label{compound-poisson-process}

\textbf{Compound Poisson R.V.} \(\to\) each arrival associated with R.V.

let \(N \overset{(d)}{=} \operatorname{Poisson}(\lambda)\), and
\(Y_1, \cdots Y_n\) be iid. with dist \(f(dx)\) \[
X := \sum_{n=0}^{N} Y_n \quad \text{is called compound poisson R.V.}
\]

\textbf{chara fun of \(X\)} \[
\begin{aligned}
E[e^{i\langle z, X\rangle}] &= E\big[e^{i\langle z, \sum_n Y_n\rangle}\big] = \sum_{k=0}^{\infty} e^{-\lambda} \frac{\lambda^k}{k!} E\big[e^{i\langle z, \sum_k Y_k\rangle}\big] \\
&= \sum_{k=0}^{\infty} e^{-\lambda} \frac{\lambda^k}{k!} \big(\underbrace{E[e^{i\langle z, Y\rangle}]}_{\hat{G}(z)}\big)^k = \sum_{k=0}^{\infty} e^{-\lambda} \frac{\lambda^k}{k!} \hat{G}(z)^k = \exp\big(-\lambda + \lambda \hat{G}(z)\big) \\
&= \exp\Big(\lambda\big(\int_{\mathbb{R}^d} e^{i\langle z, x\rangle} f(dx) - 1\big)\Big) = \exp\Big(\int_{\mathbb{R}^d}(e^{i\langle z, x\rangle} - 1) \lambda f(dx)\Big).
\end{aligned}
\]

\part{Financial Basics}

\part{Stochastic Process in Finance}

\chapter{Things need to be done}\label{things-need-to-be-done}

\begin{itemize}
\tightlist
\item
  Black-Scholes formula and risk neutral measure

  \begin{itemize}
  \tightlist
  \item
    Use monte-carlo simulation to price the option
  \item
    Binomial tree model building the model
  \item
    finding estimation of parameters
  \item
    deviation of BS model from real world using testing methods (normal
    distribution assumption, etc.)

    \begin{itemize}
    \tightlist
    \item
      Shapiro-Wilk test; Jarque-Bera test; Kolmogorov-Smirnov test;
      Anderson-Darling test; Q-Q plot; etc.
    \end{itemize}
  \end{itemize}
\item
  Option pricing

  \begin{itemize}
  \tightlist
  \item
    Option pricing formula from BS model
  \item
    numerical methods for option pricing - monte-carlo simulation,
    binomial tree model, finite difference method (numerical
    integration), etc.
  \item
    Introducing other financial derivatives
  \end{itemize}
\item
  Characteristic function and its application

  \begin{itemize}
  \tightlist
  \item
    inversion theorem
  \item
    use characteristic function to price the option
  \item
    numerical methods for characteristic function - fast Fourier
    transform (FFT), etc.
  \end{itemize}
\item
  Other stochastic model in finance

  \begin{itemize}
  \tightlist
  \item
    Heston model
  \item
    Levy processes
  \item
    CIR model
  \item
    for each model, include numerical methods for simulation;option
    pricing; properties and what difference comparing to BS model;
    parameter estimation; etc.
  \end{itemize}
\end{itemize}

\chapter{Risk-Neutral Pricing}\label{risk-neutral-pricing}

\section{Change of measure}\label{change-of-measure}

we wish to change how we measure probability from \(P\) to \(\hat{P}\)
using Random Variable \(Z\).

\textbf{Th} let \((\Omega, \mathcal{F}, P)\) be a prob space and \(Z\)
be a a.s. non-negative. random variable with \(E(Z) = 1\). for
\(A \in \mathcal{F}\), define \[
\hat{P}(A) = \int_A Z(\omega) \, dP(\omega)
\] then \(\hat{P}\) is a probability measure

\textbf{Coll} let \(\hat{P}\) be the prob measure defined, and \(X\) be
a nonnegative random variable, then \[
\hat{E}(X) = \int_\Omega X(\omega) \, d\hat{P}(\omega) = \int_\Omega X(\omega) Z(\omega) \, dP(\omega) = E(ZX).
\] if \(Z\) is a.s. stricly positive, then \[
E(X) = \hat{E}\Big(\frac{X}{Z}\Big) \qquad \text{for all nonnegative R.V. } X.
\]

Next, defines the meaning of equivalent between \(P\) and \(\hat{P}\)

\textbf{Def.} let \(P\) \& \(\hat{P}\) be two prob-meas on
\((\Omega, \mathcal{F})\), \(P\) \& \(\hat{P}\) are said equivalent if
they agree which set in \(\mathcal{F}\) have probability zero.

\textbf{Remark:}

\begin{itemize}
\tightlist
\item
  Since the set of zero prob and set of prob \(1\) are complementary,
  equivalent \(P\) \& \(\hat{P}\) agrees on event with prob \(1\) as
  well
\item
  \(P\) \& \(\hat{P}\) defined in Th are equivalent
\item
  Equivalent means that two probability measure predict the same on the
  event that happens a.s. or a.s. not happens
\end{itemize}

In finance, the model has a sample space, where all possible future
scenario lies. The probability measure for the event can be estimated
through empirical data in real world, or in the price of derivatives
securities. a risk-neutral measure \(\hat{P}\). The risk netural measure
helps to solve the equation. though \(\hat{P}\) assign different prob to
price paths, the equivalent property make sure that a hedging strategy
works a.s. in \(\hat{P}\) will also work in \(P\)

Notice that from the fact \[
\int_\Omega X(\omega) \, d\hat{P}(\omega) = \int_\Omega X(\omega) Z(\omega) \, dP(\omega) \qquad \text{and} \qquad \int_\Omega Y(\omega) \, dP(\omega) = \int_\Omega \frac{Y(\omega)}{Z(\omega)} \, d\hat{P}(\omega)
\] we can ``regard''
\(Z(\omega) = \dfrac{d\hat{P}(\omega)}{dP(\omega)}\) as in regular
calculous, as a somehow ``ratio of probability''. \(Z(\omega)\) is
refered as Radon-Nikodym derivative.

\section{Risk-Neutral Measure}\label{risk-neutral-measure}

Next, we perform similar change of measure. for the process. on
\((\Omega, \mathcal{F}, P)\) and filtration \(\mathcal{F}_t\), defined
on \(0 \leq t \leq T\). let \(Z\) be the a.s. positive R.V. with
\(E(Z) = 1\) and \(\hat{P}\) defined above. Then can define
Radon-Nikodym derivative process as follow \[
Z_t := E(Z \mid \mathcal{F}_t)
\]

\textbf{Coll} \(Z_t\) is a martingle

\textbf{proof:} let \(t > s \in [0,T]\), and \[
E(Z_t \mid \mathcal{F}_s) = E\big((E(Z \mid \mathcal{F}_t)) \mid \mathcal{F}_s\big) = E(Z \mid \mathcal{F}_s) = Z_s.
\]

\textbf{lemma:} let \(t \in [0,T]\) and \(Y\) is \(\mathcal{F}_t\)
measurable, then \(\hat{E} Y = E(Y Z_t)\)

\textbf{proof:}
\(\hat{E} Y = E(YZ) = E[E(YZ \mid \mathcal{F}_t)] = E(Y E(Z \mid \mathcal{F}_t)) = E(Y Z_t)\)

\textbf{lemma:} let \(s, t\) satisfy \(0 \leq s \leq t \leq T\) and
\(Y\) be a \(\mathcal{F}_t\) measurable R.V. then \[
\hat{E}(Y \mid \mathcal{F}_s) = \frac{1}{Z(s)} E(Y Z(t) \mid \mathcal{F}_s)
\]

\textbf{proof:} \(\frac{1}{Z(s)} E(Y Z(t) \mid \mathcal{F}_s)\) is
\(\mathcal{F}_s\) measurable, to show that it equal to conditional
expectation, we need. \[
\int_A \hat{E}(Y \mid \mathcal{F}_s) \, d\hat{P} = \int_A \frac{1}{Z(s)} E(Y Z(t) \mid \mathcal{F}_s) \, d\hat{P} = \int_A Y \, d\hat{P} \qquad \forall A \in \mathcal{F}_s
\] \[
\begin{aligned}
\int_A \frac{1}{Z(s)} E(Y Z(t) \mid \mathcal{F}_s) \, d\hat{P} &= \hat{E}\Big(\mathbf{1}_A \frac{1}{Z(s)} E(Y Z(t) \mid \mathcal{F}_s)\Big) \\
&= E\big(\mathbf{1}_A E(Y Z(t) \mid \mathcal{F}_s)\big) = E\big(E(\mathbf{1}_A Y Z(t) \mid \mathcal{F}_s)\big) \\
&= E(\mathbf{1}_A Y Z(t)) = \hat{E}(\mathbf{1}_A Y) = \int_A Y \, d\hat{P}. \qquad \blacksquare
\end{aligned}
\]

\section{Change of measure for 1D Brownian
motion}\label{change-of-measure-for-1d-brownian-motion}

\textbf{Th (change of measure for 1D BM)} Let
\((W_t)_{0 \leq t \leq T}\) be a BM on a probability space
\((\Omega, \mathcal{F}, P)\) and \((\mathcal{F}_t)_{0 \leq t \leq T}\)
be its filtration, and \((\Theta_t)_{0 \leq t \leq T}\) be another
adapted process to this filtration, then define \[
Z_t = \exp\Big(-\int_0^t \Theta_u \, dW_u - \tfrac12 \int_0^t \Theta_u^2 \, du\Big)
\] \[
\widetilde{W}_t = W_t + \int_0^t \Theta_u \, du
\] assuming \[
E\Big(\int_0^T \Theta_u^2 Z_u^2 \, du\Big) < \infty
\] Then, we have \(E(Z_T) = 1\) and \(\widetilde{W}_t\) is BM under
\(\hat{P}\) defined above.

\section{Stock under risk-neutral
measure}\label{stock-under-risk-neutral-measure}

consider a stock price \(S_t\) described by \[
dS_t = \alpha_t S_t \, dt + G_t S_t \, dW_t, \qquad 0 \leq t \leq T
\] or in the integration form \[
S_t = S_0 \exp\Big(\int_0^t G_u \, dW_u + \int_0^t (\alpha_u - \tfrac12 G_u^2) \, du\Big)
\] now consider an interest rate process \(R(t)\) and defines the
discount process as \(D_t = e^{-\int_0^t R_u \, du}\).
\(\implies dP_t = -P_t R_t \, dt\)

\textbf{Discounted stock price.} \[
D_t S_t = S_0 \exp\Big(\int_0^t G_u \, dW_u + \int_0^t (\alpha_u - \tfrac12 G_u^2 - R_u) \, du\Big)
\] \[
\begin{aligned}
d(P_t S_t) &= S_t \, dP_t + P_t \, dS_t + dS_t \, dP_t = -S_t P_t R_t \, dt + P_t(\alpha_t S_t \, dt + G_t S_t \, dW_t) \\
&= P_t S_t(\alpha_t - R_t) \, dt + G_t S_t D_t \, dW_t = G_t S_t P_t(\Theta_t \, dt + dW_t)
\end{aligned}
\] where \(\Theta_t = \dfrac{\alpha_t - R_t}{G_t}\) is the market price
of risk.

\textbf{Remark:} for discounted stock price. the mean rate of return of
the stock reduce while the volatility remains the same.

now we consider the discounted price in the prob measure \(\hat{P}\)
defined in theorem with \(\Theta_t\) used the market price of risk. then
\[
d(D_t S_t) = G_t D_t S_t \, d\widetilde{W}_t \iff D_t S_t = S_0 + \int_0^t G_u P_u S_u \, d\widetilde{W}_u.
\] and thus the discounted Stock price is martingle under this measure.
The new probability measure is called Risk-Neutral measure.

\textbf{Stock price under risk-neutral measure.} \[
dW_t = -\Theta_t \, dt + d\widetilde{W}_t \implies dS_t = R_t S_t \, dt + G_t S_t \, d\widetilde{W}_t.
\]

\textbf{Remark:}

\bookmarksetup{startatroot}

\chapter*{References}\label{references}
\addcontentsline{toc}{chapter}{References}

\markboth{References}{References}

\phantomsection\label{refs}




\end{document}
